\section{Standard form and index}

\begin{definition}[Reduced full-support source datum]
    \label{def:mpu_reduced_full_support_source_datum}
    \notready
    % **Scope restriction (source data for the same tensor):** this datum is an
    % explicit hypothesis, not a construction of a reduced representative.
    % Documented in docs/paper-gaps/mpu_canonical_form_full_support.tex.
    \uses{def:mpu, def:mpu_normalized_flattening, def:cpsv_cfii, def:mpu_full_support, thm:mpu_reduced_cfii_transfer_stabilization, thm:mpu_normalized_transfer_fin_one}
    Let $\widetilde{\mathcal U}=d^{-1/2}\mathcal U$ be the normalized
    flattening, and let $E$ be its transfer matrix. A reduced full-support source
    datum for $\mathcal U$, with physical dimension $d>0$ and positive bond
    dimension, consists of an integer $J>0$, a diagonal positive definite matrix
    $\rho$ with $\tr(\rho)=1$, and
    $(\Phi\rvert=\operatorname{vec}(\Id_D)^{\mathsf T}$ such that
    \begin{align}
        E^J & =\lvert\rho)(\Phi\rvert.
    \end{align}
    The one-dimensional case is immediate: if $D=1$, take $J=1$
    and $\rho=\Id_1$ using
    Lemma~\ref{thm:mpu_normalized_transfer_fin_one}. If $D>1$, chosen
    canonical-form-II data with full support give a positive trace-one ambient
    fixed point with $J=D^2-1$ by
    Lemma~\ref{thm:mpu_reduced_cfii_transfer_stabilization}. Diagonality of
    this ambient fixed point is an additional hypothesis: the theorem gives a
    diagonal matrix in the unique block coordinates, but conjugation into the
    ambient coordinates need not preserve diagonality.
    For $D>1$, this states the fixed-pair convention in
    \cite[equations labelled~\texttt{Erightleft} and~\texttt{II\_UTransfer}, and
        proof of Proposition~III.2(ii)]{Cirac2017MPU}. The $D=1$ branch is the direct
    one-dimensional specialization. When $\mathcal U$ is an MPU,
    Lemma~\ref{thm:mpu_canonical_form_ii_stabilization_converse} gives a
    canonical-form-II presentation of this same tensor whose ambient fixed
    matrix is the given $\rho$; it makes no assertion about continuous
    choices along a path.
    % Source lines: paper_v2.tex 269--280 and 397--405.
\end{definition}

\begin{lemma}[Staggered Gram expansion of the source tensor $u$]
    \label{thm:mpu_source_u_staggered_gram}
    \lean{MPOTensor.SourceFactors.sourceU_gram_eq_staggered_four_u}
    \leanok
    \uses{def:mpu_source_uv, def:mpu_weighted_source_factors}
    Let $M_i=X_iY_i$ be given source factorizations for $\mathcal U$ and
    the weight $\rho$. For retained physical pairs $p=(p_1,p_2)$ and
    $q=(q_1,q_2)$, one has
    \begin{align}
        \sum_{l,r}u_{(l,r),q}\overline{u_{(l,r),p}}
        ={} & \sum_{\alpha,\alpha',\beta,\beta',\gamma,j_1,j_2}
        \overline{\mathcal U^{j_1}_{p_2,\alpha,\beta}}
        \rho_{\beta,\beta'}
        \mathcal U^{j_1}_{q_2,\alpha',\beta'} \notag            \\
            & \qquad\qquad\cdot
        \overline{\mathcal U^{j_2}_{p_1,\gamma,\alpha}}
        \mathcal U^{j_2}_{q_1,\gamma,\alpha'}.
    \end{align}
    The $p$ entries are starred and the $q$ entries are unstarred. This is
    precisely the staggered reflected--direct four-tensor expression obtained
    in \cite[equations~\texttt{X1X2b} and~\texttt{uUnitary}]{Cirac2017MPU} from $u=Y_2\mathbin{-}Y_1$ and the two source-factor
    identities.
\end{lemma}

\begin{proof}\leanok
    \uses{def:mpu_source_uv, thm:mpu_source_factorizations, thm:mpu_source_factor_normalizations}
    Since the first cut has rows $(i,\beta)$ and columns $(\alpha,j)$,
    its weighted normalization gives
    \begin{align}
        \sum_r\overline{(Y_1)_{r,(\alpha,p_2)}}
        (Y_1)_{r,(\alpha',q_2)}
         & =\sum_{\beta,\beta',j_1}
        \overline{\mathcal U^{j_1}_{p_2,\alpha,\beta}}
        \rho_{\beta,\beta'}
        \mathcal U^{j_1}_{q_2,\alpha',\beta'}.
    \end{align}
    The unweighted normalization of the second cut gives
    \begin{align}
        \sum_l\overline{(Y_2)_{l,(p_1,\alpha)}}
        (Y_2)_{l,(q_1,\alpha')}
         & =\sum_{\gamma,j_2}
        \overline{\mathcal U^{j_2}_{p_1,\gamma,\alpha}}
        \mathcal U^{j_2}_{q_1,\gamma,\alpha'}.
    \end{align}
    Expanding both copies of
    $u_{(l,r),(i_1,i_2)}=\sum_\alpha
        (Y_2)_{l,(i_1,\alpha)}(Y_1)_{r,(\alpha,i_2)}$, substituting these two
    Gram identities, and reordering the finite sums yields the stated formula.
\end{proof}

\begin{lemma}[Transpose fixed-pair trace closure]
    \label{thm:mpu_source_u_transpose_fixed_pair_trace}
    \lean{MPOTensor.SourceFactors.sourceU_gram_eq_transpose_fixed_pair_trace}
    \leanok
    \discussion{7632}
    \uses{def:mpu_source_uv, def:mpu_weighted_source_factors, def:mpu_double_layer}
    Let $\mathcal U$ be an MPO tensor, let $\rho$ be a bond matrix, and let
    $M_i=X_iY_i$ be given source factorizations for $\mathcal U$ and $\rho$.
    Construct the source tensor $u=Y_2\mathbin{-}Y_1$. For retained physical
    pairs $p=(p_1,p_2)$ and $q=(q_1,q_2)$, let $W^{ij}$ denote the $(i,j)$
    double-layer letter, let $\lvert\rho^{\mathsf T})$ be the column-stacking
    vector of $\rho^{\mathsf T}$, and set
    $(\Phi\rvert=\operatorname{vec}(\Id_D)^{\mathsf T}$. Then
    \begin{align}
        \sum_{l,r}u_{(l,r),q}\overline{u_{(l,r),p}}
         & =\operatorname{tr}\!\left(
        W^{p_2q_2}\lvert\rho^{\mathsf T})(\Phi\rvert W^{p_1q_1}
        \right).
    \end{align}
    Thus the fixed-pair weight is genuinely transposed at this stage. In the
    diagonal canonical-form-II coordinates fixed above,
    $\rho^{\mathsf T}=\rho$; this is the identification used in the complete
    contraction below.
\end{lemma}

\begin{proof}\leanok
    \uses{thm:mpu_source_u_staggered_gram, thm:mpu_source_v_four_u_expansions}
    Enumerate each doubled-bond coordinate as
    $\operatorname{flat}(\alpha,\alpha')$. Expanding a diagonal entry of the
    matrix under the trace gives
    \begin{align}
        \sum_{\beta,\beta',\gamma,j_1,j_2}
        \overline{\mathcal U^{j_1}_{p_2,\alpha,\beta}}
        \mathcal U^{j_1}_{q_2,\alpha',\beta'}
        (\rho^{\mathsf T})_{\beta',\beta}
        \overline{\mathcal U^{j_2}_{p_1,\gamma,\alpha}}
        \mathcal U^{j_2}_{q_1,\gamma,\alpha'}.
    \end{align}
    Since $(\rho^{\mathsf T})_{\beta',\beta}=\rho_{\beta,\beta'}$,
    summing over $\alpha,\alpha'$ gives the identity in
    Lemma~\ref{thm:mpu_source_u_staggered_gram}.
\end{proof}

\begin{lemma}[Retained-interior source-$u$ closure]
    \label{thm:mpu_source_u_retained_interior}
    \lean{MPOTensor.SourceFactors.sourceU_gram_eq_closed_doubleLayer_trace_of_isSymm,
        MPOTensor.SourceFactors.sourceU_gram_eq_normalized_input_tail}
    \leanok
    \discussion{7633}
    \uses{thm:mpu_source_u_transpose_fixed_pair_trace,
        thm:mpu_normalized_input_tail_closed_trace}
    Let $\mathcal U$ be an MPO tensor with physical dimension $d>0$, let $W$
    be its double layer, and put $E=d^{-1}\sum_iW^{ii}$. Let $\rho$ be
    symmetric, set $(\Phi\rvert=\operatorname{vec}(\Id_D)^{\mathsf T}$, and
    suppose that for some $K\geq0$,
    \begin{align}
        E^K & =\lvert\rho)(\Phi\rvert.
    \end{align}
    For source factors built from the same weight $\rho$, the paper
    gate $u=Y_2\mathbin{-}Y_1$ satisfies
    \begin{align}
        \sum_{l,r}u_{(l,r),q}\overline{u_{(l,r),p}}
         & =\operatorname{tr}\!\left(W^{p_1q_1}W^{p_2q_2}E^K\right).
        \label{eq:mpu_source_u_trace}
    \end{align}
    If $\rho$ is diagonal, the same quantity is
    \begin{align}
        \sum_{l,r}u_{(l,r),q}\overline{u_{(l,r),p}}
         & = d^{-K}\sum_{\tau,\eta}
        \overline{U^{(K+2)}_{\eta,(p,\tau)}}
        U^{(K+2)}_{\eta,(q,\tau)}.
        \label{eq:mpu_source_u_tail_sum}
    \end{align}
    Thus only the $K$-site interior is blocked; the two endpoint letters remain
    one-site tensors. The $p$ entries are conjugated and the $q$ entries are
    not. In \cite[Lemma~III.7]{Cirac2017MPU}, this identity is the fixed-pair
    contraction.
    % Source locator: paper_v2.tex, lines 269--280 and 545--557.
    % Orientation: docs/paper-gaps/mpu_source_cut_orientation.tex.
\end{lemma}

\begin{proof}\leanok
    \uses{thm:mpu_source_u_transpose_fixed_pair_trace,
        thm:mpu_normalized_input_tail_closed_trace}
    Symmetry changes the inserted vector
    $\lvert\rho^{\mathsf T})$ to $\lvert\rho)$ in
    Lemma~\ref{thm:mpu_source_u_transpose_fixed_pair_trace}; substituting the
    fixed pair and cycling the first endpoint letter yields~(\ref{eq:mpu_source_u_trace}). The normalized input-tail identity Lemma~\ref{thm:mpu_normalized_input_tail_closed_trace} expands that trace into~(\ref{eq:mpu_source_u_tail_sum}).
\end{proof}

\begin{lemma}[Complete source-$u$ network contraction]
    \label{thm:mpu_source_u_complete_network}
    \lean{MPOTensor.SourceFactors.sourceU_gram_eq_normalized_mpo_input_tail_of_isDiag}
    \leanok
    \discussion{5982}
    \uses{def:mpu, def:mpu_normalized_flattening, def:mpu_source_uv, def:mpo_family}
    % Revalidation boundary: docs/paper-gaps/mpu_source_cut_orientation.tex.
    Let $\mathcal U$ be an MPU tensor with physical dimension $d>0$ and bond
    dimension $D>0$. Let $\rho>0$ be diagonal with $\tr(\rho)=1$, and put
    $(\Phi\rvert=\operatorname{vec}(\Id_D)^{\mathsf T}$. Suppose that for some
    $J>0$ the normalized transfer matrix satisfies
    \begin{align}
        E^J & =\lvert\rho)(\Phi\rvert.
    \end{align}
    Construct the source factors and $u$ from $\mathcal U$ using $\rho$, and set
    $K=JD^2$. For retained physical pairs
    $p,q\in\{0,\ldots,d-1\}^2$, the complete source-$u$ contraction satisfies
    \begin{align}
        \sum_{l,r}u_{(l,r),q}\overline{u_{(l,r),p}}
         & =d^{-K}\sum_{\tau,\eta}
        \overline{U^{(K+2)}_{\eta,(p,\tau)}}
        U^{(K+2)}_{\eta,(q,\tau)}.
    \end{align}
    Here $\tau$ ranges over input words of length $K$, and $\eta$ ranges over
    output words of length $K+2$. This is the fixed-pair form of the
    complete-network equality in \cite[Lemma~III.7]{Cirac2017MPU} for $u^\dagger u$.
    % Source locator: paper_v2.tex, lines 545--557.
\end{lemma}

\begin{proof}\leanok
    \uses{thm:mpu_source_u_retained_interior,
        thm:mpu_blocked_transfer_power_rank_one}
    With $K=JD^2$, trace normalization makes the given rank-one transfer
    matrix idempotent, and the normalized blocked-transfer identity therefore
    gives
    \begin{align}
        E^K & =\lvert\rho)(\Phi\rvert.
    \end{align}
    Diagonality gives $\rho^{\mathsf T}=\rho$, so
    Lemma~\ref{thm:mpu_source_u_retained_interior} applies at this $K$ and
    yields the displayed normalized input-tail sum.
\end{proof}

\begin{lemma}[Source-$v$ isometry of a simple tensor]
    \label{thm:mpu_simple_source_v_isometry}
    \lean{MPOTensor.IsMPUCanonicalFormII.sourceV_isIsometry,
        MPOTensor.IsMPUCanonicalFormII.rightRank_mul_leftRank_le}
    \leanok
    \uses{def:mpu_simple, def:mpu_source_uv, def:mpu_double_layer, def:mpu_canonical_form_ii}
    Let $\mathcal U$ be simple and in canonical form II, with ambient fixed
    matrix $\rho$. Let $v$ be the source operator
    of $\mathcal U$ built from that same weight $\rho$. Then
    $v^\dagger v=\Id_{r\ell}$, and consequently $r\ell\le d^2$. This is the
    implication $1\to2$ in \cite[proof of Theorem~III.8]{Cirac2017MPU},
    before the standing isometry of $u$ is used.
    % Source lines: paper_v2.tex 577--588.
\end{lemma}

\begin{proof}\leanok
    \uses{thm:mpug_simple_canonical_boundary_contractions,
        thm:mpu_source_v_complete_contraction, thm:mpu_source_v_dressed_gram_cancel}
    Lemma~\ref{thm:mpug_simple_canonical_boundary_contractions} gives
    the insertion identity for the recorded canonical boundaries:
    $W^{ij}W^{k\ell}=W^{ij}\lvert\rho)(\Phi\rvert W^{k\ell}$. Since $\rho$
    is diagonal, $\rho^{\mathsf T}=\rho$, so
    Lemma~\ref{thm:mpu_source_v_complete_contraction} gives the dressed Gram
    equation for the factors built from the same weight $\rho$. Then
    Lemma~\ref{thm:mpu_source_v_dressed_gram_cancel} cancels the dressing. Since $v^\dagger v=\Id$ on
    $\C^r\otimes\C^\ell$, the map $v$ is injective, so $r\ell\le d^2$.
\end{proof}

\begin{lemma}[Source-$v$ isometry after positive blocking]
    \label{thm:mpu_source_v_isometry_after_blocking}
    \lean{MPOTensor.IsMPUCanonicalFormII.sourceV_blockTensor_isIsometry,
        MPOTensor.IsMPUCanonicalFormII.exists_sourceV_blockTensor_isIsometry}
    \leanok
    \uses{thm:mpu_canonical_form_ii_positive_blocking,
        thm:mpu_simple_source_v_isometry, thm:mpu_bounded_simple_blocking}
    Let $\mathcal U$ be an MPU in canonical form II, with recorded weight
    $\rho$. If $p>0$ and $\mathcal U^{[p]}$ is simple, then the source operator
    $v$ computed from the source cuts of $\mathcal U^{[p]}$, using this same
    $\rho$, satisfies $v^\dagger v=\Id$. In particular, some integer
    $1\leq p\leq D^2$ gives a simple blocked tensor and this source-$v$
    isometry. The statement includes $D=1$.
\end{lemma}

\begin{proof}\leanok
    \uses{thm:mpu_canonical_form_ii_positive_blocking,
        thm:mpu_simple_source_v_isometry, thm:mpu_bounded_simple_blocking,
        thm:mpu_canonical_form_ii_positive_dimensions}
    The blocked tensor retains canonical form II with the same recorded
    matrix $\rho$. The source-$v$ contraction for a simple tensor therefore
    applies to its own source cuts and gives the isometry. The blocking
    length $D^2$ is simple for every MPU, including in bond dimension one.
\end{proof}

\begin{lemma}[Simplicity from the source-$v$ isometry]
    \label{thm:mpu_simple_of_source_v_isometry}
    \lean{MPOTensor.IsMPU.isMPUSimple_of_sourceV_isIsometry}
    \leanok
    \uses{def:mpu, def:mpu_simple, def:mpu_source_uv}
    Let $\mathcal U$ be an MPU tensor, let $\rho>0$, and let $v$ be its
    source operator built from $\rho$. If $v^\dagger v=\Id_{r\ell}$, then
    $\mathcal U$ is simple, with boundary vectors $(\Phi\rvert$ and
    $\lvert\rho^{\mathsf T})$. In \cite[proof of Theorem~III.8]{Cirac2017MPU}, this step proves the implication
    $4\to1$.
    % Source lines: paper_v2.tex 596--600.
\end{lemma}

\begin{proof}\leanok
    \uses{thm:mpu_source_v_dressed_gram_cancel, thm:mpu_source_v_complete_contraction, thm:mpu_all_later_simple_blockings}
    By Lemma~\ref{thm:mpu_source_v_dressed_gram_cancel}, $v^\dagger v=\Id$
    gives the dressed Gram equation, and
    Lemma~\ref{thm:mpu_source_v_complete_contraction} turns it into
    $W^{ij}W^{k\ell}=W^{ij}\lvert\rho^{\mathsf T})(\Phi\rvert W^{k\ell}$.
    The boundary-vector implication in
    Lemma~\ref{thm:mpu_all_later_simple_blockings} gives
    $(\Phi\rvert W^{ij}\lvert\rho^{\mathsf T})=\delta_{ij}$ from unitarity,
    so both simple contractions hold.
\end{proof}

\begin{lemma}[The source operator $u$ is an isometry]
    \label{lemuisometry}
    \lean{MPOTensor.IsMPUCanonicalFormII.sourceU_isIsometry,
        MPOTensor.IsMPUCanonicalFormII.mul_self_le_rightRank_mul_leftRank}
    \leanok
    \discussion{5708}
    \uses{def:mpu_source_uv, def:mpu_canonical_form_ii}
    For any MPU tensor $\mathcal U$ under the source's standing
    canonical-form-II convention, let $r$ and $\ell$ be the two source-cut
    ranks and let $u$ be the source operator. Then
    \begin{align}
        u^\dagger u & =\Id.
    \end{align}
    Consequently $u$ is an isometry and $r\ell\geq d^2$, which is the
    unrestricted statement of \cite[Lemma~III.7]{Cirac2017MPU} on the isometricity of $u$.
    % Source lines: paper_v2.tex 545--557.
\end{lemma}

\begin{proof}\leanok
    \uses{lem:mpu_admissible_source_u_isometry, thm:mpu_canonical_form_ii_transfer_stabilization}
    The convention records a positive diagonal ambient weight $\rho$ and,
    by Lemma~\ref{thm:mpu_canonical_form_ii_transfer_stabilization}, the
    stabilized power $E^J=\lvert\rho)(\Phi\rvert$ at the positive exponent
    $J=\max(D^2-1,1)$. Lemma~\ref{lem:mpu_admissible_source_u_isometry},
    applied to the source factors built from that same weight, gives
    $u^\dagger u=\Id$ and $d^2\le r\ell$.
\end{proof}

\begin{lemma}[Fixed-pair source-$u$ isometry]
    \label{lem:mpu_admissible_source_u_isometry}
    \lean{MPOTensor.SourceFactors.sourceU_isIsometry_of_isMPU,
        MPOTensor.SourceFactors.mul_self_le_rightRank_mul_leftRank_of_isMPU}
    \leanok
    \discussion{5708}
    % This lemma is the reusable generic form: it assumes the exact rank-one
    % identity for an arbitrary positive diagonal weight, and is the step used
    % by the source-labelled Lemma III.7 above. Orientation documented
    % in docs/paper-gaps/mpu_canonical_form_full_support.tex and
    % docs/paper-gaps/mpu_source_cut_orientation.tex.
    \uses{def:mpu, def:mpu_source_uv, def:mpu_weighted_source_factors,
        def:mpu_double_layer}
    Let $\mathcal U$ be an MPU tensor with physical dimension $d>0$. Let $W$
    be its double layer, put $E=d^{-1}\sum_iW^{ii}$, and let $\rho>0$ be
    diagonal. Put $(\Phi\rvert=\operatorname{vec}(\Id_D)^{\mathsf T}$ and
    suppose that for some $K\geq0$,
    \begin{align}
        E^K & =\lvert\rho)(\Phi\rvert.
    \end{align}
    Let $r$ and $\ell$ be the source-cut ranks, and construct the compact-SVD
    source factors and $u$ from $\mathcal U$ using $\rho$. Then
    \begin{align}
        u^\dagger u & =\Id,       \\
        d^2         & \leq r\ell.
    \end{align}
    Thus the compact-SVD paper gate is an isometry, giving the
    fixed-pair version of \cite[Lemma~III.7]{Cirac2017MPU} for $u$.
    % Source lines: paper_v2.tex 269--280 and 545--557.
\end{lemma}

\begin{proof}\leanok
    \uses{thm:mpu_source_u_retained_interior,
        thm:mpu_normalized_input_tail_isometry}
    The retained-interior identity and input-first MPU unitarity give, for all
    physical pairs $p,q$,
    \begin{align}
        (u^\dagger u)_{p,q}
         & =d^{-K}\sum_{\tau,\eta}
        \overline{U^{(K+2)}_{\eta,(p,\tau)}}
        U^{(K+2)}_{\eta,(q,\tau)}
        =\delta_{p,q}.
    \end{align}
    Hence $u^\dagger u=\Id$. Since $u$ maps a space of dimension $d^2$
    isometrically into one of dimension $r\ell$, one also has $d^2\leq r\ell$.
\end{proof}

\begin{theorem}[Simple-tensor equivalence theorem]
    \label{ThmFund1}
    \lean{MPOTensor.IsMPUCanonicalFormII.isMPUSimple_tfae}
    \leanok
    \discussion{5708}
    \uses{lemuisometry, def:mpu_simple, def:mpu_source_uv, def:mpu_canonical_form_ii}
    Let $\mathcal U$ be an MPU tensor under the source's standing
    canonical-form-II convention. Let $r$ and $\ell$ be its source-cut ranks,
    and let $u$ and $v$ be its source operators, both built from the ambient
    weight $\rho$ recorded by the convention. Then the following conditions
    are equivalent:
    \begin{enumerate}
        \item $\mathcal U$ is simple;
        \item $r\ell=d^2$;
        \item $u$ is unitary;
        \item $v$ is unitary.
    \end{enumerate}
    In \cite[Theorem~III.8]{Cirac2017MPU}, this equivalence is the unrestricted characterization of simple
    tensors.
    % Source lines: paper_v2.tex 563--601.
\end{theorem}

\begin{proof}\leanok
    \uses{lemuisometry,
        thm:mpu_simple_source_v_isometry,
        thm:mpu_source_gates_two_site_factorization,
        thm:mpu_simple_of_source_v_isometry}
    The convention gives $u^\dagger u=\Id$ and
    $d^2\leq r\ell$ before any of the four conditions is assumed. Follow the
    cycle $1\to2\to3\to4\to1$.

    If $\mathcal U$ is simple, the source-$v$ contraction for the same weight
    $\rho$ gives $v^\dagger v=\Id$, hence $r\ell\leq d^2$ and therefore
    $r\ell=d^2$. Conversely, this equality makes the standing isometry $u$
    square, hence unitary. The two-site factorization gives
    \begin{align}
        \operatorname{reindex}(U^{(2)})
         & =v\,\operatorname{swap}_{r,\ell}(u).
    \end{align}
    Reindexing preserves unitarity. Since $U^{(2)}$ and $u$ are unitary,
    cancellation makes $v$ unitary. Finally, $v^\dagger v=\Id$ gives the
    second simple contraction, while the MPU identity gives the first, so
    $\mathcal U$ is simple.
\end{proof}

\begin{lemma}[Three-legged decomposition of a simple tensor]
    \label{thm:mpu_three_legged_decomposition}
    \lean{MPOTensor.SourceFactors.IsThreeLegged,
        MPOTensor.SourceFactors.X₁_leftInverse,
        MPOTensor.SourceFactors.X₂_leftInverse,
        MPOTensor.IsMPUCanonicalFormII.sourceFactors_isThreeLegged,
        MPOTensor.IsMPUCanonicalFormII.exists_threeLegged}
    \leanok
    \uses{ThmFund1, def:mpu_weighted_source_factors, def:mpu_source_uv,
        def:mpu_canonical_form_ii, def:mpu_simple}
    Let $\mathcal U$ be a simple MPU tensor of physical dimension $d$ under
    the standing canonical-form-II convention, whose periodic operator is
    unitary on every chain of length at least two, and let $\rho>0$ be its
    recorded diagonal fixed point. Then the weighted source factors
    $X_1,Y_1,X_2,Y_2$ for $\rho$ satisfy:
    \begin{enumerate}
        \item the two three-legged decompositions
              $\mathcal U^{ij}_{\alpha\beta}=(X_1Y_1)_{(i,\beta),(\alpha,j)}
                  =(X_2Y_2)_{(\alpha,i),(j,\beta)}$, joined along the curly leg of
              dimension $d_r=r$ and the thick leg of dimension $d_l=\ell$;
        \item the gates $u=Y_2\mathbin{-}Y_1$ and $v=X_1\mathbin{-}X_2$ are
              unitary;
        \item the normalized contractions $X_2^\dagger X_2=\Id_{d_l}$ and
              $X_1^\dagger(I_d\otimes\rho)X_1=\Id_{d_r}$ hold;
        \item $X_1$ and $X_2$ are left-invertible, with left inverses
              $X_1^\dagger(I_d\otimes\rho)$ and $X_2^\dagger$, and $Y_1,Y_2$ are
              right-invertible;
        \item $d_ld_r=d^2$.
    \end{enumerate}
    This is the structure theorem for simple matrix product unitaries in
    canonical form II of \cite{FrancoRubio2025MPUGauging}, which asks
    unitarity on every nonempty chain; that convention implies the one used
    here.
    % Source lines: 2502.20257/main.tex 657--809.
\end{lemma}

\begin{proof}\leanok
    \uses{ThmFund1, def:mpu_weighted_source_factors}
    Items 1, 3 and the one-sided inverses of item 4 are the defining
    identities of the weighted source factors. Since $\mathcal U$ is simple,
    Theorem~\ref{ThmFund1} gives $r\ell=d^2$ and unitarity of $u$ and $v$,
    which are items 2 and 5.
\end{proof}

\begin{definition}[Standard form]
    \label{SF}
    \leanok
    \discussion{6019}
    \uses{ThmFund1, def:mpu_source_uv,
        lem:mpu_simple_reduced_literal_standard_form}
    Let $\mathcal U$ be a simple tensor with positive physical and bond
    dimensions. Choose its simple canonical representative $\mathcal V$,
    of no larger bond dimension and with the same periodic operator at every
    positive length. Its standard form is the depth-two form of the two-site
    block $\mathcal V_2$ determined by the literal source unitaries
    \begin{align}
        u & \colon\mathbb C^d\otimes\mathbb C^d
        \longrightarrow\mathbb C^\ell\otimes\mathbb C^r, \\
        v & \colon\mathbb C^r\otimes\mathbb C^\ell
        \longrightarrow\mathbb C^d\otimes\mathbb C^d.
    \end{align}
    On four original sites its unitary is
    \begin{align}
        U & =v_{23}v_{41}u_{12}u_{34}.
    \end{align}
    In \cite[Definition~III.9]{Cirac2017MPU}, this decomposition is the unrestricted definition of the
    standard form.
    % Source lines: paper_v2.tex 603--622.
\end{definition}

\begin{definition}[Two-site standard form with given gates]
    \label{def:mpu_supplied_two_site_standard_form}
    \lean{MPOTensor.TwoSiteStandardFormData}
    \leanok
    Let $d,D,\ell,r$ be positive integers and let
    $W$ be a two-site tensor of physical dimension $d^2$ and bond dimension
    $D$. A standard form with given gates consists of unitary maps
    \begin{align}
        u & \colon \mathbb C^d\otimes\mathbb C^d
        \longrightarrow \mathbb C^\ell\otimes\mathbb C^r, \\
        v & \colon \mathbb C^r\otimes\mathbb C^\ell
        \longrightarrow \mathbb C^d\otimes\mathbb C^d,
    \end{align}
    and half factors $X_1,X_2$ such that
    \begin{align}
        v^{(i_1,i_2)}_{s,l}
         & =\sum_{\beta=0}^{D-1}
        (X_1)_{(i_1,\beta),s}(X_2)_{(\beta,i_2),l}, \\
        W^{(i_1,i_2),(j_1,j_2)}_{\alpha,\gamma}
         & =\sum_{l=0}^{\ell-1}\sum_{s=0}^{r-1}
        (X_2)_{(\alpha,i_1),l}\,
        u_{(l,s),(j_1,j_2)}\,
        (X_1)_{(i_2,\gamma),s}.
    \end{align}
    This is the open contraction underlying
    \cite[Definition~III.9]{Cirac2017MPU}. No unitary periodic operator is
    assumed for $W$.
    % Source: paper_v2.tex lines 532--543 and 603--617. The generic open
    % equation is unreflected; the alternative closed reflection is specific
    % to the source tensor and is not a condition on supplied gates.
\end{definition}

\begin{lemma}[Dimensions of supplied standard-form gates]
    \label{lem:mpu_supplied_gate_dimension_identity}
    \lean{MPOTensor.TwoSiteStandardFormData.left_mul_right}
    \leanok
    Every supplied standard form of Definition~\ref{def:mpu_supplied_two_site_standard_form}
    satisfies $\ell r=d^2$. This is the dimension identity of its actual
    unitary gate in \cite[Definition~III.9]{Cirac2017MPU}.
\end{lemma}

\begin{proof}\leanok
    \uses{def:mpu_supplied_two_site_standard_form}
    The isometry and coisometry of $u$ give the two dimension inequalities.
\end{proof}

\begin{center}
    % Formula: W^{(i_1,i_2),(j_1,j_2)}_{\alpha,\gamma}=\sum_{l,s}
    % (X_2)_{(\alpha,i_1),l}\,u_{(l,s),(j_1,j_2)}\,(X_1)_{(i_2,\gamma),s}, the
    % second displayed equation of this definition; compare the source figure
    % StandardForm of \cite{Cirac2017MPU}.
    % Ink-to-index map: W is the two-site tensor with outputs i_1,i_2, inputs
    % j_1,j_2 and bonds \alpha,\gamma.  On the right, X_2 carries \alpha and
    % i_1, X_1 carries i_2 and \gamma, and the gate u takes the inputs j_1,j_2
    % to the source indices l and s.
    % Contraction set: the two vertical source indices l and s.
    % Boundary signature: (virtual-west=1 | virtual-east=1 | physical-up=2 |
    % physical-down=2) on both sides.
    \tenkzkernel
    \begin{tenkz}[cols=2, bonds=none]
        \tn[at=(1,1), skin=box, wide=2,
            ports={90@1:physical:$i_1$, 90@2:physical:$i_2$, 270@1:physical:$j_1$, 270@2:physical:$j_2$,
                    180:virtual:$\alpha$, 0:virtual:$\gamma$}]{W}
    \end{tenkz}
    \quad $=$ \quad
    \begin{tenkz}[rows={wire,wire}, cols=2, bonds=none]
        \tn[at=(1,1), name=x2, skin=box, label pos=n,
            ports={180:virtual:$\alpha$, 90:physical:$i_1$, 270:virtual}]{X_2}
        \tn[at=(1,2), name=x1, skin=box,
            ports={90:physical:$i_2$, 0:virtual:$\gamma$, 270:virtual}]{X_1}
        \tn[at=(2,1), name=u, skin=box, wide=2,
            ports={90@1:virtual, 90@2:virtual, 270@1:physical:$j_1$, 270@2:physical:$j_2$}]{u}
        \tnwire[name=wl]{x2.270}{u.90@1}
        \tnwire[name=ws]{x1.270}{u.90@2}
        \tnmark[form=label, label pos=w]{on wl 0.5}{$l$}
        \tnmark[form=label, label pos=e]{on ws 0.5}{$s$}
    \end{tenkz}
\end{center}

\begin{definition}[The two layers of a given standard form]
    \label{def:mpu_two_site_circuit_layers}
    \lean{MPOTensor.twoSiteUnshiftedLayer,
        MPOTensor.twoSiteShiftedLayer,
        MPOTensor.twoSiteShiftedPairEquiv,
        MPOTensor.twoSitePairConfigEquiv}
    \leanok
    Let $N>0$. Write each physical output and input label as an ordered pair
    $i_x=(i_{x,1},i_{x,2})$ and $j_x=(j_{x,1},j_{x,2})$, and each intermediate
    label as $h_x=(h_{x,1},h_{x,2})\in\{0,\ldots,\ell-1\}\times
        \{0,\ldots,r-1\}$. Indices $x$ are taken modulo $N$. Define the
    unshifted and shifted layer matrices by
    \begin{align}
        (C_{u,N})_{h,j}
         & =\prod_{x\in\mathbb Z/N\mathbb Z}
        u_{(h_{x,1},h_{x,2}),(j_{x,1},j_{x,2})}, \\
        (C_{v,N})_{i,h}
         & =\prod_{x\in\mathbb Z/N\mathbb Z}
        v_{(i_{x,2},i_{x+1,1}),(h_{x,2},h_{x+1,1})}.
    \end{align}
    The cyclic relabeling of pairs is
    $h\mapsto\bigl((h_{x,2},h_{x+1,1})\bigr)_x$, with inverse
    $q\mapsto\bigl((q_{x-1,2},q_{x,1})\bigr)_x$. The alphabet of the
    two-site tensor is identified with ordered pairs $(i_{x,1},i_{x,2})$.
    In \cite[Equation~\texttt{StandardForm}]{Cirac2017MPU}, this convention specifies the two-site
    ordering.
\end{definition}

\begin{lemma}[Unitarity of the two circuit layers]
    \label{thm:mpu_two_site_circuit_layers_unitary}
    \lean{MPOTensor.twoSiteUnshiftedLayer_isUnitaryBetween,
        MPOTensor.twoSiteShiftedLayer_isUnitaryBetween}
    \leanok
    \uses{def:mpu_two_site_circuit_layers}
    If the given gates $u$ and $v$ are unitary between their stated
    coordinate spaces, then for every $N>0$ both $C_{u,N}$ and $C_{v,N}$
    are unitary between their corresponding configuration spaces.
\end{lemma}

\begin{proof}\leanok
    \uses{def:mpu_two_site_circuit_layers}
    The two Gram matrices of the product of a rectangular gate factor
    sitewise: for example,
    \begin{align}
        (C_{u,N}^{\dagger}C_{u,N})_{j,j'}
         & =\prod_{x\in\mathbb Z/N\mathbb Z}
        (u^{\dagger}u)_{j_x,j'_x}=\delta_{j,j'}.
    \end{align}
    The reverse Gram identity is analogous. The shifted layer is the
    product of $v$ on relabeled configurations, and the cyclic relabeling
    preserves both Gram identities.
\end{proof}

\begin{lemma}[Periodic circuit of a given two-site standard form]
    \label{thm:mpu_supplied_two_site_periodic_circuit}
    \lean{MPOTensor.TwoSiteStandardFormData.mpo_apply_eq_shifted_mul_unshifted,
        MPOTensor.TwoSiteStandardFormData.mpo_eq_shifted_mul_unshifted}
    \leanok
    \uses{def:mpu_supplied_two_site_standard_form,
        def:mpu_two_site_circuit_layers}
    Let $W$ have given two-site standard-form data $(u,v,X_1,X_2)$.
    For every $N>0$ and pair configurations $i,j$,
    \begin{align}
        U^{(N)}(W)_{i,j}
         & =\sum_h\prod_{x\in\mathbb Z/N\mathbb Z}
        v_{(i_{x,2},i_{x+1,1}),(h_{x,2},h_{x+1,1})}
        \prod_{x\in\mathbb Z/N\mathbb Z}
        u_{(h_{x,1},h_{x,2}),(j_{x,1},j_{x,2})}
        =(C_{v,N}C_{u,N})_{i,j}.
    \end{align}
    Thus, after the stated identification of pair configurations with the
    two-site alphabet, $U^{(N)}(W)=C_{v,N}C_{u,N}$. This includes $N=1$.
    The formula is the positive-length closed contraction of the given
    open form in \cite[Equations~\texttt{uuvv} and
        \texttt{StandardForm}]{Cirac2017MPU}.
\end{lemma}

\begin{proof}\leanok
    \uses{def:mpu_supplied_two_site_standard_form,
        def:mpu_two_site_circuit_layers}
    Expanding the periodic trace as a sum over closed bond configurations and
    inserting the given $X_2$--$u$--$X_1$ expression at each site shows that, for
    each $x$, the bond label between sites $x$ and $x+1$ contracts
    $(X_1)_{(i_{x,2},\beta),h_{x,2}}$ with
    $(X_2)_{(\beta,i_{x+1,1}),h_{x+1,1}}$; their sum is precisely
    $v_{(i_{x,2},i_{x+1,1}),(h_{x,2},h_{x+1,1})}$. The remaining factors
    form $C_{u,N}$, giving the displayed matrix product in its stated order.
\end{proof}

\begin{lemma}[Positive-length unitarity of a given standard form]
    \label{thm:mpu_supplied_two_site_periodic_unitary}
    \lean{MPOTensor.TwoSiteStandardFormData.mpo_mem_unitaryGroup,
        MPOTensor.TwoSiteStandardFormData.isMPU}
    \leanok
    \uses{thm:mpu_supplied_two_site_periodic_circuit,
        thm:mpu_two_site_circuit_layers_unitary}
    If $W$ has given two-site standard-form data, then
    \begin{align}
        U^{(N)}(W)^{\dagger}U^{(N)}(W)
         & =U^{(N)}(W)U^{(N)}(W)^{\dagger}=\Id
    \end{align}
    for every $N>0$. In particular, $W$ is an MPU under the source's
    $N>1$ length convention.
\end{lemma}

\begin{proof}\leanok
    \uses{thm:mpu_supplied_two_site_periodic_circuit,
        thm:mpu_two_site_circuit_layers_unitary}
    The periodic operator is the product $C_{v,N}C_{u,N}$ of two unitary
    layers, and the physical-pair identification preserves unitarity.
\end{proof}

\begin{lemma}[Open-word factorization with free virtual ends]
    \label{thm:mpu_two_site_open_word_factorization}
    \lean{MPOTensor.twoSiteLeftHalf,
        MPOTensor.twoSiteRightHalf,
        MPOTensor.twoSiteInputGate,
        MPOTensor.twoSiteOutputLink,
        MPOTensor.twoSiteOpenGateTail,
        MPOTensor.twoSite_letter_factorization,
        MPOTensor.twoSite_adjacent_halves,
        MPOTensor.twoSite_evalWord_open_gate_factorization,
        MPOTensor.twoSite_trace_open_gate_with_endpoints}
    \leanok
    \uses{def:mpu_supplied_two_site_standard_form, def:mpo_eval_word}
    For a two-site tensor $W$ with given standard-form data, write
    $L(i)_{\alpha t}=X_2(\alpha,i_1)_t$,
    $R(i)_{s\gamma}=X_1(i_2,\gamma)_s$,
    $u(j)_{ts}=u_{(t,s),j}$, and
    $v(i,i')_{st}=v_{(i_2,i'_1),(s,t)}$.
    For every nonempty word of complete two-site blocks, with the product
    ordered in increasing $x$,
    \begin{align}
        W^{i_0,j_0}\cdots W^{i_{n-1},j_{n-1}}
         & =L(i_0)u(j_0)
        \left(\prod_{x=1}^{n-1}v(i_{x-1},i_x)u(j_x)\right)R(i_{n-1}).
    \end{align}
    The equality is an equality of matrices with both virtual indices free;
    arbitrary matrices may therefore be inserted at the two ends before
    taking the trace.
\end{lemma}

\begin{center}
    % Formula: the case n=3 of W^{i_0,j_0}W^{i_1,j_1}W^{i_2,j_2}
    % =L(i_0)u(j_0)v(i_0,i_1)u(j_1)v(i_1,i_2)u(j_2)R(i_2).
    % Ink-to-index map: the three boxes W on the left are consecutive two-site
    % letters, upper legs the outputs i_{x,1},i_{x,2} and lower legs the inputs
    % j_{x,1},j_{x,2}.  On the right, the lower layer carries the three gates
    % u(j_x); the upper layer carries L(i_0) on the output i_{0,1}, the two
    % gates v(i_{x-1},i_x) on the outputs i_{x-1,2},i_{x,1}, and R(i_2) on the
    % output i_{2,2}.
    % Contraction set: on the left, the two virtual bonds between letters; on
    % the right, the six vertical source indices between the two layers.
    % Boundary signature: (virtual-west=1 | virtual-east=1 | physical-up=6 |
    % physical-down=6) on both sides; the two virtual ends stay free, as stated.
    \tenkzkernel
    \begin{tenkz}[cols=6, bonds=none]
        \tn[at=(1,1), name=w0, skin=box, wide=2,
            ports={90@1:physical, 90@2:physical, 270@1:physical, 270@2:physical, 180:virtual, 0:virtual}]{W}
        \tn[at=(1,3), name=w1, skin=box, wide=2,
            ports={90@1:physical, 90@2:physical, 270@1:physical, 270@2:physical, 180:virtual, 0:virtual}]{W}
        \tn[at=(1,5), name=w2, skin=box, wide=2,
            ports={90@1:physical, 90@2:physical, 270@1:physical, 270@2:physical, 180:virtual, 0:virtual}]{W}
        \tnwire{open w}{w0.180} \tnwire{w0.0}{w1.180} \tnwire{w1.0}{w2.180} \tnwire{w2.0}{open e}
    \end{tenkz}
    \quad $=$ \quad
    \begin{tenkz}[rows={wire,wire}, cols=6, bonds=none]
        \tn[at=(1,1), name=L, skin=box, ports={180:virtual, 90:physical, 270:virtual}]{L}
        \tn[at=(1,2), name=v1, skin=box, wide=2,
            ports={90@1:physical, 90@2:physical, 270@1:virtual, 270@2:virtual}]{v}
        \tn[at=(1,4), name=v2, skin=box, wide=2,
            ports={90@1:physical, 90@2:physical, 270@1:virtual, 270@2:virtual}]{v}
        \tn[at=(1,6), name=R, skin=box, ports={0:virtual, 90:physical, 270:virtual}]{R}
        \tn[at=(2,1), name=u0, skin=box, wide=2,
            ports={90@1:virtual, 90@2:virtual, 270@1:physical, 270@2:physical}]{u}
        \tn[at=(2,3), name=u1, skin=box, wide=2,
            ports={90@1:virtual, 90@2:virtual, 270@1:physical, 270@2:physical}]{u}
        \tn[at=(2,5), name=u2, skin=box, wide=2,
            ports={90@1:virtual, 90@2:virtual, 270@1:physical, 270@2:physical}]{u}
        \tnwire{open w}{L.180} \tnwire{R.0}{open e}
        \tnwire{L.270}{u0.90@1} \tnwire{v1.270@1}{u0.90@2}
        \tnwire{v1.270@2}{u1.90@1} \tnwire{v2.270@1}{u1.90@2}
        \tnwire{v2.270@2}{u2.90@1} \tnwire{R.270}{u2.90@2}
    \end{tenkz}
\end{center}

\begin{proof}\leanok
    \uses{def:mpu_supplied_two_site_standard_form}
    A single letter factors as $W^{i,j}=L(i)u(j)R(i)$. For adjacent
    letters, $R(i)L(i')=v(i,i')$ by the definition of the shifted gate.
    Induction over the remaining letters gives the displayed product.
    Multiplying it by the two end matrices and taking the trace preserves
    the equality.
\end{proof}

\begin{definition}[Open three-block coefficient and neighboring observable]
    \label{def:mpu_two_site_open_three_block}
    \lean{MPOTensor.twoSiteThreeBlockOpen,
        MPOTensor.twoLinkOutputObservable,
        MPOTensor.twoSiteThreeBlockTrace,
        MPOTensor.twoSiteBufferedTrace}
    \leanok
    \uses{def:mpu_supplied_two_site_standard_form, def:mpo_eval_word}
    Let $W$ have given two-site standard-form data. Fix virtual indices
    $\alpha,\gamma$, exterior output half-sites $e_-,e_+$, and exterior
    input pairs $j_0,j_2$. For two neighboring output pairs
    $a=(a_0,a_1)$, $b=(b_0,b_1)$ and a middle input pair $j$, set
    \begin{align}
        F_{\alpha,\gamma}(a,b;j)
         & =\bigl[W^{(e_-,a_0),j_0}
            W^{(a_1,b_0),j}W^{(b_1,e_+),j_2}\bigr]_{\alpha,\gamma}.
    \end{align}
    For an observable $A$ on the middle input pair, let
    $C=uAu^{\dagger}$ and define the matrix on the two neighboring output
    pairs by
    \begin{align}
        B(A) & =(v\otimes v)\,\widetilde C\,(v\otimes v)^{\dagger}, \\
        \widetilde C_{(s_0,t_1,s_1,t_2),(s'_0,t'_1,s'_1,t'_2)}
             & =\delta_{s_0,s'_0}
        C_{(t_1,s_1),(t'_1,s'_1)}\delta_{t_2,t'_2}.
    \end{align}
    The link coordinates are ordered as $(r,\ell,r,\ell)$. For arbitrary
    virtual matrices $P,Q$ and complete-block word pairs
    $(\sigma^{L},\tau^{L})$ and $(\sigma^{R},\tau^{R})$, set
    \begin{align}
        F_{P,Q}(a,b;j)
         & =\operatorname{tr}\bigl(P\,W^{(e_-,a_0),j_0}
        W^{(a_1,b_0),j}W^{(b_1,e_+),j_2}\,Q\bigr),         \\
        F_{\sigma,\tau}(a,b;j)
         & =\operatorname{tr}\bigl(W^{\sigma^{L},\tau^{L}}
        W^{(e_-,a_0),j_0}W^{(a_1,b_0),j}W^{(b_1,e_+),j_2}
        W^{\sigma^{R},\tau^{R}}\bigr),
    \end{align}
    where $W^{\sigma,\tau}$ denotes word evaluation. In \cite[Equations~\texttt{uuvv} and \texttt{StandardForm}]{Cirac2017MPU}, these expressions are the finite
    open
    contractions.
\end{definition}

\begin{lemma}[Open three-block locality]
    \label{lem:mpu_two_site_open_three_block_locality}
    \lean{MPOTensor.twoSite_threeBlock_open_intertwiner}
    \leanok
    \uses{def:mpu_two_site_open_three_block}
    For every given two-site standard form and every observable $A$ on
    the middle input pair,
    \begin{align}
        F_{\alpha,\gamma}A & =B(A)F_{\alpha,\gamma}
    \end{align}
    for every choice of $\alpha,\gamma,e_-,e_+,j_0,j_2$. In particular,
    neither virtual end is closed in this identity, and $B(A)$ acts only on
    the four output half-sites forming the two neighboring pairs.
\end{lemma}

\begin{proof}\leanok
    \uses{def:mpu_supplied_two_site_standard_form}
    The two internal contractions $X_1X_2$ give one $v$ gate on each
    neighboring output pair. Thus, for suitable exterior factors $L,R$,
    \begin{align}
        F_{\alpha,\gamma}(a,b;j)
         & =\sum_{s_0,t_1,s_1,t_2}
        L_{\alpha}(s_0)\,v_{a,(s_0,t_1)}
        u_{(t_1,s_1),j}\,v_{b,(s_1,t_2)}R_{\gamma}(t_2),
    \end{align}
    where $L$ and $R$ retain the exterior half factors and input gates.
    Since $u$ and $v$ are unitary between their stated spaces,
    $uA=(uAu^{\dagger})u$ and
    $(v\otimes v)\widetilde C=B(A)(v\otimes v)$.
    Inserting these two identities into the displayed contraction proves
    the matrix equality without a virtual trace.
\end{proof}

\begin{lemma}[Virtual-end and word-buffer locality]
    \label{lem:mpu_two_site_open_three_block_buffers}
    \lean{MPOTensor.twoSite_threeBlock_trace_intertwiner_with_endpoints,
        MPOTensor.twoSite_threeBlock_buffered_intertwiner}
    \leanok
    \uses{lem:mpu_two_site_open_three_block_locality,
        def:mpu_two_site_open_three_block, def:mpo_eval_word}
    Replace the open word in $F_{\alpha,\gamma}$ by its virtual trace with
    arbitrary matrices $P,Q$ at the left and right ends. The resulting
    matrix $F_{P,Q}$ obeys
    \begin{align}
        F_{P,Q}A & =B(A)F_{P,Q}.
    \end{align}
    The same conclusion holds when $P$ and $Q$ are the matrix products of
    arbitrary complete-block prefix and suffix words, respectively. Their
    physical row and column words remain explicit and may have different
    lengths on the two sides of the three-block window.
\end{lemma}

\begin{proof}\leanok
    \uses{lem:mpu_two_site_open_three_block_locality,
        lem:mpo_eval_word_append}
    If $E$ is the open three-letter word, cyclicity of the virtual trace
    gives
    \begin{align}
        \operatorname{tr}(PEQ)
         & =\sum_{\gamma,\alpha}
        (QP)_{\gamma,\alpha}E_{\alpha,\gamma}.
    \end{align}
    The open identity therefore survives this linear combination. For
    complete-block prefix and suffix words, apply word concatenation twice
    to identify their product with $PEQ$, leaving both factors intact.
\end{proof}

\begin{lemma}[Left-canonical normalization of a canonical MPU]
    \label{thm:mpu_cfii_left_canonical}
    \lean{MPOTensor.IsMPUCanonicalFormII.isLeftCanonical_normalizedFlattening}
    \leanok
    \uses{def:mpu_canonical_form_ii, def:mpu_normalized_flattening}
    Let $\mathcal U$ be a matrix product unitary in full-support canonical
    form II, with bond dimension $D$ and physical dimension $d$. Its normalized
    flattening $A^{ij}=d^{-1/2}U^{ij}$ is left-canonical:
    \begin{align}
        \sum_{i,j=0}^{d-1} A^{ij\dagger}A^{ij} & =\Id_D.
    \end{align}
    In \cite[lines 269--281 and 344--356]{Cirac2017MPU}, this left-canonical identity is the implicit ambient
    normalization.
\end{lemma}

\begin{proof}\leanok
    \uses{thm:mpu_normalized_transfer_trace, thm:mpu_shifted_transfer_trace_single_block, thm:mpu_transfer_eigenvalue_one, thm:mpu_ambient_block_isometry_intertwining, thm:mpu_unique_full_support_geometry}
    The normalized transfer traces satisfy $\tr(E^N)=1$ for $N>1$. Hence the
    canonical form has one retained block $B$, with coefficient $\mu$
    satisfying $|\mu|^2=1$. Let $W$ be its inclusion into the ambient bond
    space. Full support gives $W^\dagger W=\Id$ and $WW^\dagger=\Id$, while
    the reconstruction gives $A^{ij}W=W(\mu B^{ij})$. Therefore
    \begin{align}
        \sum_{i,j}A^{ij\dagger}A^{ij}
         & =W\left(\sum_{i,j}(\mu B^{ij})^\dagger(\mu B^{ij})\right)W^\dagger
        =WW^\dagger=\Id_D.
    \end{align}
\end{proof}

\begin{lemma}[Virtual unitary gauge from equal periodic operators]
    \label{thm:mpu_cfii_equal_periodic_virtual_unitary}
    \lean{MPOTensor.IsMPUCanonicalFormII.exists_unitary_virtual_gauge_of_mpo_eq}
    \leanok
    \uses{def:mpu_canonical_form_ii, def:mpu}
    Let $\mathcal U$ and $\mathcal V$ be matrix product unitaries on the same
    physical space, each in full-support canonical form II, with bond
    dimensions $D_{\mathcal U}$ and $D_{\mathcal V}$. Suppose that their
    periodic operators agree at every length $N>1$:
    \begin{align}
        U^{(N)} & =V^{(N)}.
    \end{align}
    Then $D_{\mathcal U}=D_{\mathcal V}$, and, after identifying their bond
    spaces by this equality, there is a unitary
    $z\in\mathrm U(D_{\mathcal V})$ such that, for every $i,j$,
    \begin{align}
        V^{ij} & =zU^{ij}z^\dagger.
    \end{align}
    This is the virtual-bond implication underlying the local relation in
    \cite[Equations (SFuu) and (SFvv), lines 624--652]{Cirac2017MPU}, with \cite{gap:mpu_canonical_form_full_support} recording the
    full-support
    convention.
\end{lemma}

\begin{proof}\leanok
    \uses{thm:mpu_canonical_form_ii_normal, thm:normal_proportional_mpv_geometric_phase, thm:cpsv_normal_mpv_phase_gauge, thm:mpu_cfii_left_canonical, lem:left_canonical_irreducible_same_phase_unitary_gauge}
    The rescaled flattenings $A_{\mathcal U}^{ij}=d^{-1/2}U^{ij}$ and
    $A_{\mathcal V}^{ij}=d^{-1/2}V^{ij}$ are both normal tensors. Equality
    of periodic operators for $N>1$ makes their matrix product vectors
    equal at these lengths. The normal-tensor phase alternative gives a unit
    scalar $a$ such that
    \begin{align}
        V^{(N)}(A_{\mathcal U})_\sigma
         & =a^N V^{(N)}(A_{\mathcal V})_\sigma
    \end{align}
    at every positive length. At $N=2,3$ the normalized states are nonzero,
    so the given equality yields $a^2=a^3=1$, hence $a=1$. Gauge recovery for
    normal tensors now gives equal bond dimensions and an invertible $X$ and
    nonzero scalar $\zeta$ satisfying
    $A_{\mathcal V}^{ij}=\zeta XA_{\mathcal U}^{ij}X^{-1}$. Applying this
    relation to lengths two and three, where the common states are nonzero,
    gives $\zeta^2=\zeta^3=1$, so $\zeta=1$.
    Lemma~\ref{thm:mpu_cfii_left_canonical} makes both tensors
    left-canonical, so the irreducible unitary-gauge theorem replaces $X$
    by a unitary $z$. Cancelling the common factor $d^{-1/2}$ gives the
    asserted relation for $U^{ij}$ and $V^{ij}$.
\end{proof}

\begin{definition}[Open standard form from source cuts]
    \label{def:mpu_source_open_two_site_standard_form}
    \lean{MPOTensor.IsMPUCanonicalFormII.twoSiteStandardFormData}
    \leanok
    \uses{def:mpu_supplied_two_site_standard_form, ThmFund1,
        def:mpu_source_uv}
    Let $\mathcal U$ be simple and have a reduced full-support
    canonical-form-II presentation with positive diagonal fixed matrix $\rho$.
    The two-site tensor $\mathcal U_2$ has a two-site standard form with
    the source gates $u,v$ and the source half factors $X_1,X_2$. In
    particular,
    \begin{align}
        (\mathcal U_2)^{(i_1,i_2),(j_1,j_2)}_{\alpha,\gamma}
         & =\sum_{l,s}(X_2)_{(\alpha,i_1),l}
        u_{(l,s),(j_1,j_2)}(X_1)_{(i_2,\gamma),s},                  \\
        v^{(i_1,i_2)}_{s,l}
         & =\sum_{\beta}(X_1)_{(i_1,\beta),s}(X_2)_{(\beta,i_2),l}.
    \end{align}
    The note \cite{gap:mpu_canonical_form_full_support} records the full-support
    restriction.
    % Source: paper_v2.tex equations uuvv and StandardForm, lines 532--543
    % and 603--622. This is only the open construction, not the closed
    % four-site equality in Definition SF.
\end{definition}

\begin{proof}\leanok
    \uses{ThmFund1, thm:mpu_source_u_open_two_site,
        thm:mpu_canonical_form_ii_positive_dimensions}
    The source-cut factorization gives the open contraction of
    $\mathcal U_2$ through $X_2,u,X_1$, while the definition of $v$ gives
    its $X_1,X_2$ contraction. The simple-tensor equivalence makes both
    source gates unitary. Its rank identity $r\ell=d^2$ also makes the two
    intermediate dimensions positive.
\end{proof}

\begin{lemma}[Source circuit on the original even chain]
    \label{thm:mpu_source_even_periodic_circuit}
    \lean{MPOTensor.twoSiteChainConfigEquiv,
        MPOTensor.IsMPUCanonicalFormII.mpo_even_eq_source_circuit}
    \leanok
    \uses{def:mpu_source_open_two_site_standard_form,
        def:mpu_two_site_circuit_layers}
    Let $\mathcal U$ be simple and have a reduced full-support
    canonical-form-II presentation with positive diagonal fixed matrix $\rho$.
    Let $u$ and $v$ be its source gates. For $N>0$, flatten each pair
    configuration by the bijection $\Phi_N$ given by
    \begin{align}
        (\Phi_N i)_{2x}   & =i_{x,1}, &
        (\Phi_N i)_{2x+1} & =i_{x,2}.
    \end{align}
    Then the original chain of $2N$ sites satisfies
    \begin{align}
        U^{(2N)}(\mathcal U)_{\Phi_N i,\Phi_N j}
         & =(C_{v,N}C_{u,N})_{i,j}.
    \end{align}
    In particular, numbering four original sites from $1$ to $4$ gives
    \begin{align}
        U^{(4)}(\mathcal U) & =v_{23}v_{41}u_{12}u_{34}.
    \end{align}
    % This is the supplied-CFII corollary of the source equation, not the
    % unrestricted representative theorem in Definition SF. See
    % docs/paper-gaps/mpu_canonical_form_full_support.tex and issue 7127.
\end{lemma}

\begin{proof}\leanok
    \uses{def:mpu_source_open_two_site_standard_form,
        thm:mpu_supplied_two_site_periodic_circuit,
        thm:mpdo_block_two_as_positive_blocking,
        thm:mpdo_closed_chain_physical_blocking}
    Applying the periodic circuit identity to the two-site block of $\mathcal U$,
    the two physical-blocking identities identify its length-$N$ periodic
    operator with $U^{(2N)}(\mathcal U)$ under $\Phi_N$. For $N=2$, the two
    shifted pairs are $(2,3)$ and $(4,1)$.
\end{proof}

\begin{lemma}[Bounded block with an open source standard form]
    \label{thm:mpu_bounded_source_open_standard_form}
    \lean{MPOTensor.twoSiteDirectBlockEquiv,
        MPOTensor.reindexPhysical_blockTwo_blockTensor_two_mul,
        MPOTensor.IsMPUCanonicalFormII.exists_twoSiteStandardFormData_blockTensor}
    \leanok
    \uses{def:mpu_source_open_two_site_standard_form}
    Let $\mathcal U$ have a reduced full-support canonical-form-II
    presentation. There is an integer $k$ with $0<k\leq D^2$ such that
    $\mathcal U^{[k]}$ is simple and its two-site block has the open source
    standard form of Definition~\ref{def:mpu_source_open_two_site_standard_form}.
    If $\eta_{d,k}$ identifies a word of length $2k$ with two consecutive
    words of length $k$, the corresponding physical relabeling satisfies
    \begin{align}
        \operatorname{reindex}_{\eta_{d,k}}
        \bigl((\mathcal U^{[k]})_2\bigr)
         & =\mathcal U^{[2k]}.
    \end{align}
    The source cuts and both gates are formed on the same tensor
    $\mathcal U^{[k]}$, using the fixed matrix of the given presentation.
    % Source: paper_v2.tex lines 378--427 and 603--622.
\end{lemma}

\begin{proof}\leanok
    \uses{def:mpu_source_open_two_site_standard_form,
        thm:mpu_bounded_simple_blocking,
        thm:mpu_canonical_form_ii_positive_blocking,
        thm:mpu_canonical_form_ii_positive_dimensions,
        thm:mpdo_block_two_as_positive_blocking,
        lem:mpu_iterated_physical_blocking}
    At the simple blocking length $k=D^2$,
    canonical form II persists under positive blocking with the same
    fixed matrix, so Definition~\ref{def:mpu_source_open_two_site_standard_form} applies to
    $\mathcal U^{[k]}$. Two-site blocking is direct blocking by $2k$ after
    the displayed relabeling: first identify the ordinary two-site block
    with length-two blocking, then group the direct word into consecutive
    length-$k$ subwords and use $k\cdot2=2k$.
\end{proof}

\begin{lemma}[Original-site residual words around two-site blocks]
    \label{lem:mpu_original_site_residual_two_site_word}
    \lean{MPOTensor.mpo_apply_prefix_twoSiteBlock_suffix}
    \leanok
    \uses{def:mpdo_physical_blocking, def:mpdo_two_site_blocking,
        def:mpo_eval_word, def:mpo_family}
    Let $U$ be any MPO tensor, let $k,p,m,s\in\mathbb N$, and write
    $W=(U^{[k]})_2$. For a word $\sigma$ of $m$ blocks of length $2k$,
    let $\eta\sigma$ relabel each block as two consecutive words of length
    $k$; use the same convention for $\tau$. If $i^-,j^-$ are words of
    length $p$ and $i^+,j^+$ are words of length $s$, then
    \begin{align}
        \rho^{(p+2km+s)}(U)_{i^-\widehat\sigma i^+,\,
            j^-\widehat\tau j^+}
         & =\operatorname{tr}\!\left(
        U^{i^-,j^-}\cdot W^{\eta\sigma,\eta\tau}
        \cdot U^{i^+,j^+}\right).
    \end{align}
    Both residual products remain present, including when $m=0$; no
    unitarity, canonical form, or divisibility of the total length is
    assumed.
\end{lemma}

\begin{proof}\leanok
    \uses{lem:mpdo_physical_blocking_residuals,
        thm:mpu_bounded_source_open_standard_form}
    Residual-site blocking gives the trace with middle factor
    $(U^{[2k]})^{\sigma,\tau}$. The two-site relabeling identifies this
    factor, letter by letter and in the same order, with
    $W^{\eta\sigma,\eta\tau}$; the prefix and suffix products are unchanged.
\end{proof}

\begin{lemma}[Original-chain coefficient at a three-block window]
    \label{thm:mpu_original_site_three_block_window_entry}
    \lean{MPOTensor.originalSite_threeBlock_window_entry}
    \leanok
    \uses{def:mpdo_physical_blocking, def:mpu_two_site_open_three_block,
        def:mpo_family}
    Suppose $W=(U^{[k]})_2$ has given two-site standard-form data.
    Let three consecutive blocks of length $2k$ lie between an arbitrary
    original-site prefix of length $p$ and suffix of length $s$. Under the
    two-site relabeling, suppose their output labels are
    $(e_-,a_0),(a_1,b_0),(b_1,e_+)$ and their input labels are
    $j_0,j,j_2$. Let $P=U^{i^-,j^-}$ and $Q=U^{i^+,j^+}$ be the residual
    products. The corresponding entry of the original length-$p+6k+s$
    MPO is
    \begin{align}
        \rho^{(p+6k+s)}(U)_{i^-\widehat\sigma i^+,\,
            j^-\widehat\tau j^+}
         & = F_{P,Q}(a,b;j).
    \end{align}
    The equality holds for every choice of the residual words and window
    labels; no unitarity or MPU hypothesis on $U$ is required.
\end{lemma}

\begin{proof}\leanok
    \uses{lem:mpu_original_site_residual_two_site_word}
    Applying the residual-word identity to the three middle blocks shows that their
    relabeled words are precisely the three displayed letters, so the
    middle product is the one defining $F_{P,Q}$.
\end{proof}

\begin{lemma}[Buffered original-chain three-block locality]
    \label{thm:mpu_original_site_buffered_three_block_locality}
    \lean{MPOTensor.originalSite_threeBlock_window_with_buffers_entry,
        MPOTensor.originalSiteThreeBlockWindow,
        MPOTensor.originalSite_threeBlock_window_with_buffers_intertwiner}
    \leanok
    \uses{def:mpdo_physical_blocking, def:mpu_two_site_open_three_block,
        def:mpo_family}
    Suppose $W=(U^{[k]})_2$ has given two-site standard-form data.
    Let $i^-,j^-$ be original-site words of length $p$ and $i^+,j^+$
    original-site words of length $s$. Let
    $(\sigma^{L},\tau^{L})$ and $(\sigma^{R},\tau^{R})$ be complete-block
    word pairs, with $|\sigma^{L}|=|\tau^{L}|$. For each choice of the
    window labels $a,b,j$, let $\sigma(a,b)$ and $\tau(j)$ be words of $m$
    blocks of length $2k$ whose relabelings in the sense of the
    residual-word identity are
    \begin{align}
        \eta\sigma(a,b)
         & =\sigma^{L}\,\bigl((e_-,a_0),(a_1,b_0),(b_1,e_+)\bigr)\,
        \sigma^{R},                                                 \\
        \eta\tau(j)
         & =\tau^{L}\,(j_0,j,j_2)\,\tau^{R}.
    \end{align}
    Put
    \begin{align}
        P & =U^{i^-,j^-}\,W^{\sigma^{L},\tau^{L}}, \qquad
        Q=W^{\sigma^{R},\tau^{R}}\,U^{i^+,j^+},                    \\
        F^U(a,b;j)
          & =\rho^{(p+2km+s)}(U)_{i^-\widehat{\sigma(a,b)}\,i^+,\,
                j^-\widehat{\tau(j)}\,j^+}.
    \end{align}
    Then the original-chain coefficient matrix $F^U$ satisfies
    \begin{align}
        F^U(a,b;j) & = F_{P,Q}(a,b;j),
        \label{eq:mpu_orig_coeff_pq}   \\
        F^U A      & = B(A)F^U.
        \label{eq:mpu_orig_coeff_intertwiner}
    \end{align}
    Both equalities are finite-chain identities with every exterior
    physical label fixed; the virtual boundary is not discarded.
\end{lemma}

\begin{proof}\leanok
    \uses{lem:mpu_original_site_residual_two_site_word,
        lem:mpo_eval_word_append,
        lem:mpu_two_site_open_three_block_buffers}
    The residual-word identity first leaves both original-site products
    in the trace. Two word concatenations split the relabeled middle
    product into the left complete-block word, the three window letters,
    and the right complete-block word, yielding~(\ref{eq:mpu_orig_coeff_pq})
    entrywise with $P$ and $Q$ the displayed end products. Applying the
    arbitrary-virtual-end three-block identity to these two matrices yields~(\ref{eq:mpu_orig_coeff_intertwiner}).
\end{proof}

\begin{lemma}[Three-block locality with original-site residuals]
    \label{lem:mpu_two_site_original_residual_locality}
    \lean{MPOTensor.twoSite_threeBlock_originalSite_residual_intertwiner}
    \leanok
    \uses{def:mpu_two_site_open_three_block, def:mpo_eval_word}
    Let $W=(U^{[k]})_2$ have given two-site standard-form data. For
    arbitrary original-site prefix and suffix word pairs, put
    $P=U^{i^-,j^-}$ and $Q=U^{i^+,j^+}$. The three-block coefficient
    contracted against these end matrices satisfies
    \begin{align}
        F_{P,Q}A & =B(A)F_{P,Q}.
    \end{align}
    Thus the local observable on the two neighboring output pairs is
    independent of both residual words.
\end{lemma}

\begin{proof}\leanok
    \uses{lem:mpu_two_site_open_three_block_buffers}
    The virtual-end identity with the stated word products as its
    two end matrices gives the factorization.
\end{proof}

\begin{theorem}[Fundamental Theorem of MPU]
    \label{FundamentalMPU}
    \notready
    \discussion{6020}
    \uses{SF}
    Let $\mathcal U$ and $\mathcal V$ be simple tensors written in standard
    form, with source unitaries $(u_{\mathcal U},v_{\mathcal U})$ and
    $(u_{\mathcal V},v_{\mathcal V})$. They generate the same MPU for every
    $N$,
    \begin{align}
        U^{(N)} & =V^{(N)},
    \end{align}
    if and only if their standard forms are related by the local unitary gauges
    of \cite[Eqns.~(SFuu) and (SFvv)]{Cirac2017MPU}: there are unitaries $x$,
    $y$, and $z$ on the three corresponding source and canonical-form-II legs.
    In particular, the source unitaries obey the local relation
    \begin{align}
        u_{\mathcal V} & =(x\otimes y)u_{\mathcal U},
    \end{align}
    while the two canonical-form-II tensors are related on their virtual leg by
    $z$ and on their source legs by $x$ and $y$. This records the source's
    unblocked all-length wording. Its converse has the odd-length parity
    obstruction of Lemma~\ref{thm:mpu_standard_form_parity_obstruction};
    the blocked-family interpretation is distinguished in
    \cite{gap:mpu_standard_form_parity_gap}. The correct form of this
    theorem is Lemma~\ref{thm:mpu_fundamental_letters} together with
    Lemmas~\ref{thm:mpu_fundamental_gates_necessity}
    and~\ref{thm:mpu_fundamental_gates_sufficiency}: equality of the
    operators at every length is equivalent to a unitary conjugation of the
    letters, the gate relations follow from it, and the gate relations give
    equality of the operators on rings of even length only.
    % The source's converse from the half-factor relations alone has no
    % formal counterpart; this entry therefore keeps \notready.
    % Source label: FundamentalMPU. Source lines: paper_v2.tex 624--648.
\end{theorem}

\begin{lemma}[Fundamental theorem of MPU, letters]
    \label{thm:mpu_fundamental_letters}
    \lean{MPOTensor.IsMPUCanonicalFormII.mpo_eq_iff_exists_unitary_conj,
        MPOTensor.mpo_eq_of_unitary_conj}
    \leanok
    \uses{def:mpu_canonical_form_ii, def:mpo_family}
    Let $\mathcal U$ and $\mathcal V$ be tensors in canonical form II with
    bond dimensions $D_1$ and $D_2$. Then $U^{(N)}=V^{(N)}$ for every
    $N\ge 2$ if and only if $D_1=D_2$ and there is a unitary $z$ on
    $\mathbb C^{D_1}$ with $V^{ij}=zU^{ij}z^\dagger$ for all $i,j$.
    % Source: paper_v2.tex 624--648 with the remark at 647--648 that $z$ is
    % unitary; the letter-level equivalence itself is not printed in the source.
\end{lemma}
\begin{proof}\leanok
    \uses{thm:mpu_cfii_equal_periodic_virtual_unitary}
    The forward direction is the unitary virtual gauge obtained from equal
    operators. For the converse, conjugating every letter by $z$ leaves each
    trace $\operatorname{tr}(V^{i_1j_1}\cdots V^{i_Nj_N})$ unchanged.
\end{proof}

\begin{lemma}[Fundamental theorem of MPU, gates, necessity]
    \label{thm:mpu_fundamental_gates_necessity}
    \lean{MPOTensor.IsMPUCanonicalFormII.exists_source_gate_unitary_gauges_of_mpo_eq,
        MPOTensor.IsMPUCanonicalFormII.eq_virtualSandwich_of_mpo_eq,
        MPOTensor.IsMPUCanonicalFormII.fundamentalTheorem}
    \leanok
    \uses{def:mpu_canonical_form_ii, def:mpu_simple, def:mpu_source_uv,
        def:mpu_virtual_sandwich}
    Let $\mathcal U$ and $\mathcal V$ be simple tensors in canonical form II
    of the same bond dimension with $U^{(N)}=V^{(N)}$ for every $N\ge2$. Then
    $\mathcal V=z\,\mathcal U\,z^\dagger$ letterwise for a unitary $z$, the
    source ranks of the two tensors agree, and after identifying the rank
    spaces there are unitaries $x$ on $\mathbb C^{\ell}$ and $y$ on
    $\mathbb C^{r}$ with $u_{\mathcal V}=(x\otimes y)u_{\mathcal U}$ and
    $v_{\mathcal V}=v_{\mathcal U}(y^\dagger\otimes x^\dagger)$, the gates
    being those of the selected factorizations.
    % Source: paper_v2.tex 624--648, equation SFuu and the contraction of SFvv.
\end{lemma}
\begin{proof}\leanok
    \uses{thm:mpu_fundamental_letters,
        thm:mpu_source_gates_virtual_unitary_gauge}
    Lemma~\ref{thm:mpu_fundamental_letters} gives $z$; the selected source
    gates of a unitary virtual conjugate are related to the reindexed gates
    by unitaries on the two rank spaces.
\end{proof}

\begin{lemma}[Fundamental theorem of MPU, gates, sufficiency]
    \label{thm:mpu_fundamental_gates_sufficiency}
    \lean{MPOTensor.TwoSiteStandardFormData.mpo_eq_of_gate_gauges,
        MPOTensor.IsMPUCanonicalFormII.mpo_two_mul_eq_of_source_gate_gauges,
        MPOTensor.IsMPUCanonicalFormII.mpo_blockTwo_eq_of_source_gate_gauges,
        MPOTensor.TwoSiteStandardFormData.reindexRanks,
        MPOTensor.siteProduct}
    \leanok
    \uses{def:mpu_canonical_form_ii, def:mpu_simple, def:mpu_source_uv}
    Let $\mathcal U$ and $\mathcal V$ be simple tensors in canonical form II
    of the same bond dimension whose selected source gates satisfy, after an
    identification of the rank spaces, $u_{\mathcal V}=(x\otimes y)u_{\mathcal U}$
    and $v_{\mathcal V}=v_{\mathcal U}(y^\dagger\otimes x^\dagger)$ for
    unitaries $x$, $y$. Then $U^{(2N)}=V^{(2N)}$ for every $N\ge1$. The same
    holds for any two tensors carrying two-site standard-form data with gates
    so related on the same internal legs.
    % Source: paper_v2.tex 624--648, the "if" direction read for the standard
    % forms; the parity obstruction limits it to even lengths.
\end{lemma}
\begin{proof}\leanok
    \uses{thm:mpu_standard_form_parity_obstruction}
    Expressing both operators as the two-layer circuit of the standard form shows that in
    the circuit of $\mathcal V$ the unitary placed on the internal leg of
    each site by the layer of $u$ meets its adjoint placed there by the layer
    of $v$, so the circuits coincide. The relabelling of the rank spaces is
    carried through the standard-form data.
\end{proof}

\begin{lemma}[The gate statements hold for every decomposition]
    \label{thm:mpu_gate_statements_any_svd}
    \lean{MPOTensor.TwoSiteStandardFormData.ofGateGauges,
        MPOTensor.IsMPUCanonicalFormII.sourceU_sourceFactorsOf_isIsometry,
        MPOTensor.IsMPUCanonicalFormII.isMPUSimple_iff_sourceU_sourceFactorsOf_isUnitaryBetween,
        MPOTensor.IsMPUCanonicalFormII.isMPUSimple_iff_sourceV_sourceFactorsOf_isUnitaryBetween,
        MPOTensor.IsMPUCanonicalFormII.exists_twoSiteStandardFormData_sourceFactorsOf,
        MPOTensor.IsMPUCanonicalFormII.exists_source_gate_unitary_gauges_of_mpo_eq_sourceFactorsOf,
        MPOTensor.SourceFactors.mpo_two_reindex_eq_sourceV_mul_sourceU_swap,
        MPOTensor.exists_source_factor_unitary_gauges_of_virtualSandwich,
        MPOTensor.IsMPUCanonicalFormII.exists_selected_source_factor_unitary_gauges_sourceFactorsOf,
        MPOTensor.IsMPUCanonicalFormII.mpo_two_mul_eq_of_source_gate_gauges_sourceFactorsOf}
    \leanok
    \uses{thm:mpu_source_factors_change_of_svd, lemuisometry, ThmFund1, SF, FundamentalMPU,
        thm:mpu_fundamental_gates_necessity, thm:mpu_fundamental_gates_sufficiency}
    Let $\mathcal U$ be a matrix product unitary in canonical form II with
    $\rho$ its fixed matrix, and let the factors and gates be built from any
    compact decompositions of its two source cuts. Then $u$ is an isometry;
    $\mathcal U$ is simple if and only if $u$ is unitary, if and only if $v$ is
    unitary; if $\mathcal U$ is simple, its two-site block carries standard-form
    data with these gates; and if $\mathcal U$, $\mathcal V$ are simple, in
    canonical form II, of the same bond dimension, with equal periodic
    operators at every length greater than one, then for any decompositions of the cuts of each there are
    identifications of the rank spaces and unitaries $x$, $y$ with
    $u_{\mathcal V}=(x\otimes y)u_{\mathcal U}$ and
    $v_{\mathcal V}=v_{\mathcal U}(y^\dagger\otimes x^\dagger)$; conversely, gates so
    related give equal operators on every ring of even length. The two-site
    operator is $v$ times the swapped $u$ for the factors of any decomposition,
    and the factor relations under a bond unitary hold for any decompositions
    of the two tensors. Standard-form data transform along gate gauges: from data with gates $u,v$ and unitaries
    $x$, $y$ one obtains data with gates $(x\otimes y)u$, $v(y^\dagger\otimes x^\dagger)$.
    % Source: the source never fixes a decomposition (line 479); these are its
    % gate statements, lines 545--622 and 624--648, for every choice.
\end{lemma}

\begin{proof}
    \leanok
    \uses{thm:mpu_source_factors_change_of_svd}
    The gates of any choice are those of the fixed choice multiplied by the
    unitaries $x_2\otimes x_1$ and $x_1^\dagger\otimes x_2^\dagger$; isometry and
    unitarity are preserved, the standard-form data transform along the gauge,
    and the gate relation of the fundamental theorem composes with the
    change-of-decomposition unitaries of each tensor.
\end{proof}

\begin{lemma}[Scalar invariance of the double layer]
    \label{thm:mpu_double_layer_scalar_invariance}
    \lean{MPOTensor.physicalAdjointTensor_smul,
        MPOTensor.doubleLayerTensor_smul_of_star_mul_self}
    \leanok
    \uses{def:mpu_double_layer}
    For any tensor $\mathcal U$ and scalar $c\in\mathbb C$,
    $(c\mathcal U)^\sharp=\overline c\,\mathcal U^\sharp$.
    If $\overline c c=1$, the double layers of $c\mathcal U$ and
    $\mathcal U$ agree.
\end{lemma}
\begin{proof}\leanok
    Conjugation sends $c$ to $\overline c$, and contracting the two
    layers multiplies their scalars to $\overline c c=1$.
\end{proof}

\begin{lemma}[Simplicity of the identity tensor]
    \label{thm:mpu_identity_tensor_simple}
    \lean{MPOTensor.isMPUSimple_idTensor}
    \leanok
    \uses{def:mpu_simple, def:mpo_stacked_layers}
    For every physical dimension $d$, the bond-one identity tensor
    $\mathcal I^{ij}=\delta_{ij}$ is simple.
\end{lemma}
\begin{proof}\leanok
    Its double-layer letters are $\delta_{ij}$, so both simplicity
    identities hold with boundary vectors $a=b=1$.
\end{proof}

\begin{lemma}[Parity obstruction to an unblocked standard-form converse]
    \label{thm:mpu_standard_form_parity_obstruction}
    \lean{MPOTensor.exists_simple_blockTwo_eq_mpo_three_ne}
    \leanok
    \uses{def:mpu_simple, def:mpdo_physical_blocking, def:mpo_family}
    There are simple matrix product unitaries $\mathcal U,\mathcal V$ of
    physical dimension two and bond dimension one whose two-site blocked
    tensors agree, but whose periodic operators differ at length three:
    \begin{align}
        \mathcal U_2 & =\mathcal V_2, & U^{(3)} & \ne V^{(3)}.
    \end{align}
    This is the parity obstruction to reading the converse of
    \cite[Theorem~FundamentalMPU, lines~624--648]{Cirac2017MPU} as an
    assertion about the original tensors at every length. Definition~SF of
    that source concerns the two-site blocks, and \cite{gap:mpu_standard_form_parity_gap} records the explicit source factors
    and local unitary relations for this
    example.
\end{lemma}

\begin{proof}\leanok
    \uses{def:mpu_simple, def:mpo_family, def:mpo_stacked_layers,
        thm:mpu_double_layer_scalar_invariance, thm:mpu_identity_tensor_simple,
        thm:mpdo_block_tensor_scalar_scaling,
        thm:mpdo_mpo_scalar_scaling}
    Let $\mathcal U^{ij}=\delta_{ij}$ on its one-dimensional bond space and
    set $\mathcal V^{ij}=-\mathcal U^{ij}$. Their double-layer letters are
    both $\delta_{ij}$, so the simplicity contractions hold with unit
    boundary vectors. Two-site blocking removes the pair of minus signs,
    whereas closing a chain of length three leaves one minus sign:
    \begin{align}
        \mathcal V_2 & =\mathcal U_2, &
        U^{(3)}      & =I_8,          & V^{(3)} & =-I_8.
    \end{align}
\end{proof}
\begin{lemma}[The parity gates coincide]
    \label{thm:mpu_parity_gates_coincide}
    \lean{MPOTensor.parityIdentitySourceFactors, MPOTensor.parityNegativeSourceFactors,
        MPOTensor.sourceU_parity_eq, MPOTensor.sourceV_parity_eq,
        MPOTensor.parityIdentityTensor_exists_canonicalFormII,
        MPOTensor.parityNegativeIdentityTensor_exists_canonicalFormII,
        MPOTensor.rightRank_parityIdentityTensor, MPOTensor.leftRank_parityIdentityTensor,
        MPOTensor.rightRank_parityNegativeIdentityTensor,
        MPOTensor.leftRank_parityNegativeIdentityTensor}
    \leanok
    \uses{thm:mpu_standard_form_parity_obstruction, def:mpu_right_rank}
    The two tensors of the parity obstruction are in canonical form II with
    $\rho=1$, and have cut matrices $\pm I_2$ and ranks $r=\ell=2$. With the factors $X_1=X_2=I$, $Y_1=Y_2=Z_1=Z_2=I$ for the identity
    tensor and $X_1=X_2=I$, $Y_1=Y_2=Z_1=Z_2=-I$ for its negative, both satisfying the
    six factor identities with $\rho=1$, the gates $u$ and $v$ computed from the
    factors are the same for both tensors: $u_{(l,r),(i_1i_2)}=\delta_{li_1}\delta_{ri_2}$
    and $v_{(j_1j_2),(r,l)}=\delta_{j_1r}\delta_{j_2l}$.
    % Source: the chapter's parity example; gap note mpu_standard_form_parity_gap.tex
    % (which places the signs on $X$ instead; both choices give identity gates).
\end{lemma}

\begin{proof}\leanok
    With bond dimension one the gate formulas $u=Y_2Y_1$ and $v=X_1X_2$ have one
    term each, and $(-I)(-I)=I$.
\end{proof}

\begin{remark}
    The source-labelled Fundamental Theorem above remains unformalized.
    The source's displayed standard forms concern
    $\mathcal U_2,\mathcal V_2$; their local relations can control the blocked
    families, but do not distinguish the original families at odd lengths.
    The explicit local-gauge calculation and the possible corrected scopes
    are given in \cite{gap:mpu_standard_form_parity_gap}. No odd-length
    premise has been inserted into the source-labelled statement.
\end{remark}

\begin{definition}[Specified-tensor source-index value]
    \label{def:mpu_specified_source_index_value}
    \lean{MPOTensor.sourceIndexValue}
    \leanok
    \uses{def:mpu_right_rank, def:mpu_left_rank}
    For a specified tensor $\mathcal U$ whose right and left source-cut ranks
    $r$ and $\ell$ are positive, define
    \begin{align}
        \operatorname{ind}_{\mathrm{src}}(\mathcal U)
         & =\frac12(\log_2 r-\log_2\ell).
    \end{align}
    This definition makes no choice of a simple blocking.
    % Auxiliary to paper_v2.tex lines 681--688; issue 6255.
\end{definition}

\begin{lemma}[Auxiliary source-index identities]
    \label{thm:mpu_specified_source_index_identities}
    \lean{MPOTensor.sourceIndexValue_eq_zero_of_rightRank_eq_leftRank,
        MPOTensor.sourceIndexValue_eq_of_common_rank_scale,
        MPOTensor.sourceIndexValue_eq_add_of_common_rank_product,
        MPOTensor.sourceIndexValue_eq_logb_rightRank_div,
        MPOTensor.sourceIndexValue_eq_neg_logb_leftRank_div}
    \leanok
    \uses{def:mpu_specified_source_index_value, def:mpu_right_rank, def:mpu_left_rank}
    A specified tensor with equal positive source ranks
    $r_{\mathcal U}=\ell_{\mathcal U}$ has
    $\operatorname{ind}_{\mathrm{src}}(\mathcal U)=0$.
    Let $\mathcal U$ and $\mathcal V$ be specified tensors. Suppose that their
    source ranks obey
    \begin{align}
        r_{\mathcal V}    & =c r_{\mathcal U},   &
        \ell_{\mathcal V} & =c\ell_{\mathcal U},
    \end{align}
    where $c$, $r_{\mathcal U}$, and $\ell_{\mathcal U}$ are positive. Then
    \begin{align}
        \operatorname{ind}_{\mathrm{src}}(\mathcal V)
         & =\operatorname{ind}_{\mathrm{src}}(\mathcal U).
    \end{align}
    Let $\mathcal U$, $\mathcal V$, and $\mathcal W$ instead have positive
    source ranks and suppose that
    \begin{align}
        r_{\mathcal W}    & =c r_{\mathcal U}r_{\mathcal V},      \\
        \ell_{\mathcal W} & =c\ell_{\mathcal U}\ell_{\mathcal V},
    \end{align}
    where $c$ is a positive integer. Then
    \begin{align}
        \operatorname{ind}_{\mathrm{src}}(\mathcal W)
         & =\operatorname{ind}_{\mathrm{src}}(\mathcal U)
        +\operatorname{ind}_{\mathrm{src}}(\mathcal V).
    \end{align}
    If $d$, $r_{\mathcal U}$, and $\ell_{\mathcal U}$ are positive and a
    specified tensor has $r_{\mathcal U}\ell_{\mathcal U}=d^2$, then
    \begin{align}
        \operatorname{ind}_{\mathrm{src}}(\mathcal U)
         & =\log_2\frac{r_{\mathcal U}}d
        =-\log_2\frac{\ell_{\mathcal U}}d.
    \end{align}
    % Auxiliary to paper_v2.tex lines 681--704 and IndexTh lines 837--845;
    % issues 6255 and 6923.
\end{lemma}

\begin{proof}\leanok
    \uses{def:mpu_specified_source_index_value}
    By the definition of the index value, equal ranks give a vanishing difference of
    logarithms. For common scaling, apply
    $\log_2(cx)=\log_2c+\log_2x$ to both ranks and cancel the two copies of
    $\log_2c$. For the three-tensor identity, apply the product formula twice
    to each rank and cancel the common term. For the equivalent formulas,
    apply the same product identity to $r_{\mathcal U}\ell_{\mathcal U}=d^2$
    and use $\log_2(x/y)=\log_2x-\log_2y$.
\end{proof}

\begin{lemma}[Source-index value of an independent tensor product]
    \label{thm:mpu_source_index_value_tensor_product}
    \lean{MPOTensor.sourceIndexValue_tensorProduct}
    \leanok
    \uses{def:mpu_specified_source_index_value, thm:mpu_tensor_product_source_ranks}
    If $\mathcal U$ and $\mathcal V$ have positive right and left source
    ranks, then
    \begin{align}
        \operatorname{ind}_{\mathrm{src}}(\mathcal U\boxtimes\mathcal V)
         & =\operatorname{ind}_{\mathrm{src}}(\mathcal U)
        +\operatorname{ind}_{\mathrm{src}}(\mathcal V).
    \end{align}
    This is the specified-tensor form of the tensoring case of
    \cite[Theorem~IV.6(ii)]{Cirac2017MPU}; it makes no choice of a simple
    blocking.
    % Scope restriction (specified tensors): documented in
    % docs/paper-gaps/mpu_shift_specified_tensor_index_scope.tex.
    % Source lines: paper_v2.tex 837.
\end{lemma}

\begin{proof}\leanok
    \uses{thm:mpu_tensor_product_source_ranks, thm:mpu_specified_source_index_identities}
    Both source ranks are multiplicative, so the logarithms add.
\end{proof}

\begin{lemma}[Conditional blocking invariance of the specified source-index value]
    \label{thm:mpu_specified_source_index_blocking_invariance}
    \lean{MPOTensor.sourceIndexValue_blockTensor_eq_of_products}
    \leanok
    \uses{def:mpu_specified_source_index_value}
    Let $d>0$ and $k_0\leq k$. Suppose that the endpoint source ranks satisfy
    \begin{align}
        r_{k_0}\ell_{k_0} & =d^{2k_0}, \\
        r_k\ell_k         & =d^{2k}.
    \end{align}
    Then the specified source-index values agree:
    \begin{align}
        \operatorname{ind}_{\mathrm{src}}(\mathcal U_k)
         & =\operatorname{ind}_{\mathrm{src}}(\mathcal U_{k_0}).
    \end{align}
    This statement assumes the endpoint products; it neither selects a simple
    blocking nor defines the public index.
    % Auxiliary to paper_v2.tex lines 690--704; issue 6958.
    % **Local fix (rank-product exponent):** Source line 703 prints $d^k$; the
    % consistent blocked identity is $d^{2k}$. See
    % docs/paper-gaps/mpu_blocking_rank_product_exponent.tex.
\end{lemma}

\begin{proof}\leanok
    \uses{thm:mpu_specified_source_index_identities, thm:mpu_source_cut_rank_blocking_saturation}
    Exact blocking growth gives the common positive factor $d^{k-k_0}$:
    \begin{align}
        r_k    & =d^{k-k_0}r_{k_0},    \\
        \ell_k & =d^{k-k_0}\ell_{k_0}.
    \end{align}
    Its logarithm occurs once in each rank term and therefore cancels from
    their difference.
\end{proof}

\begin{definition}[Index]
    \label{def:index}
    \lean{MPOTensor.IsMPU.index_eq_logb_of_isMPUCanonicalForm,
        MPOTensor.IsMPU.index_eq_logb_rightRank_div_of_isMPUCanonicalForm,
        MPOTensor.IsMPU.rightRank_mul_leftRank_blockTensor_of_isMPUCanonicalForm}
    \leanok
    \discussion{6014}
    \uses{ThmFund1, thm:mpu_bounded_simple_blocking, def:mpu_canonical_form,
        thm:mpu_canonical_form_similar_cfii}
    Let $\mathcal U$ be a tensor in canonical form generating matrix product
    unitaries, let $k\geq1$ be any blocking length for which $\mathcal U_k$ is
    simple, and let $r$ and $\ell$ be its right and left source-cut ranks. The
    source index is
    \begin{align}
        \operatorname{ind}(\mathcal U)
         & =\frac12(\log_2 r-\log_2\ell).
    \end{align}
    Since $r\ell=d^{2k}$, this is equivalently
    \begin{align}
        \operatorname{ind}(\mathcal U)
         & =\log_2\frac{r}{d^k}
        =-\log_2\frac{\ell}{d^k}.
    \end{align}
    In \cite[Definition~IV.1]{Cirac2017MPU}, this formula is stated
    for a tensor in canonical form with the ranks of its simple blocking, after
    the source has assumed without loss of generality that the tensor is in
    canonical form II (lines 294 and 356). The index of every matrix product
    unitary is taken through a canonical-form-II representative; for a tensor in
    canonical form this value is the displayed expression in the ranks of its
    own blocking, since these ranks are invariant under the similarity to
    canonical form II.
    % Source lines: paper_v2.tex 681--688.
\end{definition}

\begin{definition}[Index of a canonical-form-II tensor]
    \label{def:mpu_admissible_index}
    \lean{MPOTensor.IsMPUCanonicalFormII.index,
        MPOTensor.IsMPUCanonicalFormII.index_eq_logb_rightRank_div,
        MPOTensor.IsMPUCanonicalFormII.index_eq_neg_logb_leftRank_div}
    \leanok
    % **Scope restriction (canonical-form-II representative):** this index is
    % defined for a tensor already in canonical form II, not yet for every MPU.
    % Documented in docs/paper-gaps/mpu_canonical_form_full_support.tex.
    \uses{def:mpu_canonical_form_ii, def:mpu_specified_source_index_value}
    Let $\mathcal U$ be an MPU in canonical form II. Choose any positive
    blocking length $k$ for which $\mathcal U_k$ is simple, and let $r_k$ and
    $\ell_k$ be its right and left source-cut ranks. Such a blocking exists,
    and both ranks are positive. Define
    \begin{align}
        \operatorname{ind}(\mathcal U)
         & =\frac12(\log_2 r_k-\log_2\ell_k).
    \end{align}
    Since $r_k\ell_k=d^{2k}$, one also has
    \begin{align}
        \operatorname{ind}(\mathcal U)
         & =\log_2\frac{r_k}{d^k}
        =-\log_2\frac{\ell_k}{d^k}.
    \end{align}
    In \cite[Definition~IV.1]{Cirac2017MPU}, this formula is the canonical-form-II representative
    case.
    % Source lines: paper_v2.tex 681--688. The unrestricted source definition
    % remains at def:index above.
\end{definition}

\begin{lemma}[Comparison of canonical source indices]
    \label{thm:mpu_cfii_index_periodic_comparison}
    \lean{MPOTensor.IsMPUCanonicalFormII.sourceRanks_eq_of_mpo_eq,
        MPOTensor.IsMPUCanonicalFormII.index_eq_of_mpo_eq}
    \leanok
    \uses{def:mpu_canonical_form_ii, def:mpu_admissible_index,
        def:mpo_family}
    Let $U$ and $V$ be full-support canonical-form-II tensors on the same
    physical space, possibly with different bond dimensions. If
    $U^{(N)}=V^{(N)}$ for every $N>1$, then their right and left source-cut
    ranks agree, and $\operatorname{ind}(U)=\operatorname{ind}(V)$.
    In the setting of \cite[Definition~IV.1 and Proposition~IV.2]{Cirac2017MPU}, this result compares two canonical
    representatives.
\end{lemma}

\begin{proof}\leanok
    \uses{thm:mpu_cfii_equal_periodic_virtual_unitary,
        thm:mpu_virtual_sandwich_source_ranks,
        lem:mpu_common_simple_blocking_finite,
        prop:mpu_admissible_index_well_defined}
    Equality of the periodic operators gives a unitary virtual similarity,
    which preserves both source-cut ranks. For a common positive simple
    blocking length $k$ of $U$ and $V$, the periodic
    operators of the $k$-site blocks agree, so the ranks
    $r_k$ and $\ell_k$ of these blocks satisfy
    $r_k(U)=r_k(V)$ and $\ell_k(U)=\ell_k(V)$. Hence
    \begin{align}
        \operatorname{ind}(U)
         & =\frac12\bigl(\log_2 r_k(U)-\log_2\ell_k(U)\bigr)
        =\frac12\bigl(\log_2 r_k(V)-\log_2\ell_k(V)\bigr)
        =\operatorname{ind}(V).
    \end{align}
\end{proof}

\begin{definition}[Index through a reduced canonical representative]
    \label{def:mpu_reduced_representative_index}
    \lean{MPOTensor.IsMPU.index}
    \leanok
    \uses{thm:mpu_reduced_cfii_representative, def:mpu_admissible_index}
    Let $d,D>0$, and let $U$ be an MPU tensor. Choose a reduced
    canonical-form-II representative $U_{\mathrm{red}}$ with
    $U_{\mathrm{red}}^{(N)}=U^{(N)}$ for every $N>0$, and set
    \begin{align}
        \operatorname{ind}(U) & =\operatorname{ind}(U_{\mathrm{red}}).
    \end{align}
    The definition does not identify the source-cut ranks of $U$ itself
    with those of $U_{\mathrm{red}}$.
\end{definition}

\begin{lemma}[Independence of the reduced representative]
    \label{thm:mpu_reduced_representative_index_invariance}
    \lean{MPOTensor.IsMPU.index_eq_canonical_representative,
        MPOTensor.IsMPU.index_eq_of_mpo_eq}
    \leanok
    \uses{def:mpu_reduced_representative_index,
        def:mpu_canonical_form_ii, def:mpo_family}
    Every canonical-form-II representative with the same positive-length
    periodic operators as $U$ has index $\operatorname{ind}(U)$. In
    particular, positive-dimensional MPU tensors $U,V$ with
    $U^{(N)}=V^{(N)}$ for all $N>1$ have equal indices, even when their
    bond dimensions differ.
\end{lemma}

\begin{proof}\leanok
    \uses{thm:mpu_reduced_cfii_representative,
        thm:mpu_cfii_index_periodic_comparison}
    For reduced canonical-form-II representatives of $U$ and $V$,
    their periodic operators agree for $N>1$, so their canonical source
    indices agree by Lemma~\ref{thm:mpu_cfii_index_periodic_comparison}.
\end{proof}

\begin{lemma}[Canonical representative of an injective MPU]
    \label{thm:mpu_injective_cfii_representative}
    \lean{MPOTensor.IsMPU.exists_cfii_representative_of_isInjective}
    \leanok
    \uses{def:mpu, def:injective, def:mpu_canonical_form_ii}
    Let $d>0$ and let $\mathcal U$ be an injective MPU tensor of bond
    dimension $D$. Combining the canonical replacement of
    \cite[lines~319--356]{Cirac2017MPU} with the equality case of
    \cite[Proposition~20]{Molnar2018SemiInjective}, there are a
    canonical-form-II tensor $\mathcal V$ of the same bond dimension and an
    invertible matrix $X$ such that, for every $N>0$,
    \begin{align}
        \mathcal V^{ij} & =X\mathcal U^{ij}X^{-1}, \\
        V^{(N)}         & =U^{(N)}.
    \end{align}
\end{lemma}

\begin{proof}\leanok
    \uses{thm:mpu_reduced_cfii_representative,
        thm:mps_equality_reduction_exists, thm:mps_rectangular_reduction_identities}
    A reduced canonical representative $\mathcal V$ has bond dimension
    $E\leq D$, and equality of the positive-length periodic operators
    identifies the doubled-index matrix product vectors of $\mathcal U$ and
    $\mathcal V$. Since $\mathcal U$ is injective, there are
    $P\in\C^{D\times E}$ and $Q\in\C^{E\times D}$ with $PQ=I_D$ and
    $P\mathcal V^{ij}Q=\mathcal U^{ij}$. This left-invertibility implies
    $D\leq E$ and hence $D=E$, so for square matrices $PQ=I_D$ also gives
    $QP=I_D$. Multiplying the letter identity on the left by $Q$ and on the
    right by $P$ yields $\mathcal V^{ij}=Q\mathcal U^{ij}P$, as required
    with $X=Q$.
\end{proof}

\begin{lemma}[Source ranks and index of a simple injective MPU]
    \label{thm:mpu_simple_injective_source_index}
    \lean{MPOTensor.IsMPU.rightRank_mul_leftRank_of_isInjective_of_isMPUSimple,
        MPOTensor.IsMPU.index_eq_logb_of_isInjective_of_isMPUSimple}
    \leanok
    \uses{def:mpu, def:injective, def:mpu_simple,
        def:mpu_reduced_representative_index, def:mpu_right_rank, def:mpu_left_rank}
    Let $d>0$ and let $\mathcal U$ be a simple injective MPU tensor as in
    \cite[lines~1403--1407]{FrancoRubio2025MPUGauging}. Following
    \cite[lines~562--573 and~681--704]{Cirac2017MPU}, its source-cut ranks
    $r$ and $\ell$ satisfy
    \begin{align}
        r\ell & =d^2,                                    \\
        \operatorname{ind}(\mathcal U)
              & =\frac12\bigl(\log_2 r-\log_2\ell\bigr).
    \end{align}
\end{lemma}

\begin{proof}\leanok
    \uses{thm:mpu_injective_cfii_representative,
        thm:mpu_simple_rectangular_inclusion, thm:mpu_virtual_sandwich_source_ranks,
        thm:mpu_reduced_representative_index_invariance,
        thm:mpu_bounded_simple_blocking,
        thm:mpu_source_ranks_physical_reindexing,
        ThmFund1, prop:mpu_admissible_index_well_defined}
    Passing to the same-dimensional canonical representative
    $\mathcal V^{ij}=X\mathcal U^{ij}X^{-1}$ preserves both simplicity and
    the two source ranks. The simple canonical tensor $\mathcal V$ therefore
    has rank product $r\ell=d^2$. Since its one-site blocking only relabels
    the physical alphabet, its index is
    $\frac12(\log_2 r-\log_2\ell)$, and equality of the periodic operators
    identifies this value with $\operatorname{ind}(\mathcal U)$.
\end{proof}


\begin{lemma}[Blocking and adjunction of the representative index]
    \label{thm:mpu_representative_index_blocking_adjoint}
    \lean{MPOTensor.IsMPU.index_blockTensor,
        MPOTensor.IsMPU.index_physicalAdjointTensor}
    \leanok
    \uses{def:mpu_reduced_representative_index,
        def:mpdo_physical_blocking, def:mpdo_physical_adjoint}
    Let $U$ be a positive-dimensional MPU tensor. For every positive integer
    $p$, physical blocking preserves its representative index, while
    physical adjunction reverses its sign:
    \begin{align}
        \operatorname{ind}(U_p)      & =\operatorname{ind}(U),  \\
        \operatorname{ind}(U^\sharp) & =-\operatorname{ind}(U).
    \end{align}
\end{lemma}

\begin{proof}\leanok
    \uses{thm:mpu_reduced_cfii_representative,
        thm:mpu_reduced_representative_index_invariance,
        prop:mpu_admissible_index_well_defined,
        thm:mpu_cfii_index_physical_adjoint,
        thm:mpdo_closed_chain_physical_blocking, def:mpdo_physical_adjoint}
    Let $U_{\mathrm{red}}$ be a reduced canonical-form-II representative.
    Under the standard identification of blocked configurations with
    concatenated configurations, the periodic operators satisfy
    \begin{align}
        (U_p)^{(N)}      & = U^{(Np)},          &
        (U^\sharp)^{(N)} & = (U^{(N)})^\dagger.
    \end{align}
    Hence blocking and adjunction preserve equality of the positive-length
    periodic operators. The respective canonical-form-II index identities
    give
    $\operatorname{ind}((U_{\mathrm{red}})_p)
        =\operatorname{ind}(U_{\mathrm{red}})$ and
    $\operatorname{ind}(U_{\mathrm{red}}^\sharp)
        =-\operatorname{ind}(U_{\mathrm{red}})$.
    Representative independence then gives both displayed identities.
\end{proof}

\begin{lemma}[Well-definedness of the index]
    \label{index-well-defined}
    \lean{MPOTensor.IsMPU.blockTensor_add_of_isMPUCanonicalForm,
        MPOTensor.IsMPU.leftRank_blockTensor_add_of_isMPUCanonicalForm,
        MPOTensor.IsMPU.blockTensor_isMPUSimple_of_isMPUCanonicalForm}
    \leanok
    \discussion{6014}
    \uses{def:index, ThmFund1, thm:mpu_source_cut_rank_blocking, thm:mpu_source_cut_rank_blocking_saturation, thm:mpu_specified_source_index_blocking_invariance, thm:mpu_all_later_simple_blockings, thm:mpu_canonical_form_similar_cfii}
    Let $\mathcal U$ be a tensor in canonical form generating matrix product
    unitaries. The source index is independent of the chosen simple blocking.
    More precisely, if $k\geq k_0\geq1$ and $\mathcal U_{k_0}$ is simple, then
    $\mathcal U_k$ is simple and
    \begin{align}
        r_k    & =d^{k-k_0}r_{k_0},    &
        \ell_k & =d^{k-k_0}\ell_{k_0}.
    \end{align}
    In \cite[Proposition~IV.2]{Cirac2017MPU}, this statement is formulated for a tensor in canonical
    form, the hypothesis under which the source defines the index.
    % **Local fix (rank-product exponent):** the source's proof prints
    % $r_k\ell_k=d^k$; it is $d^{2k}$. Documented in
    % docs/paper-gaps/mpu_blocking_rank_product_exponent.tex.
    % Source lines: paper_v2.tex 690--704.
\end{lemma}

\begin{proof}\leanok
    \uses{thm:mpu_canonical_form_similar_cfii, thm:mpu_virtual_sandwich_blocking,
        prop:mpu_admissible_index_well_defined}
    Passing to the similar tensor in canonical form II preserves the
    operators, the simple blockings, and the ranks; for that representative the
    statement is the canonical-form-II version.
\end{proof}

\begin{lemma}[Well-definedness on a canonical-form-II representative]
    \label{prop:mpu_admissible_index_well_defined}
    \lean{MPOTensor.IsMPUCanonicalFormII.index_eq_logb_of_isMPUSimple_blockTensor,
        MPOTensor.IsMPUCanonicalFormII.sourceIndexValue_blockTensor_eq,
        MPOTensor.IsMPUCanonicalFormII.sourceIndexValue_eq_index,
        MPOTensor.IsMPUCanonicalFormII.rightRank_mul_leftRank_blockTensor,
        MPOTensor.IsMPUCanonicalFormII.sourceRanks_blockTensor_pos,
        MPOTensor.IsMPUCanonicalFormII.index_eq,
        MPOTensor.IsMPUCanonicalFormII.index_blockTensor}
    \leanok
    % **Scope restriction (canonical-form-II representative):** this proves
    % blocking and presentation independence for one such tensor, not equality
    % across arbitrary representatives of the same MPU. Documented in
    % docs/paper-gaps/mpu_canonical_form_full_support.tex.
    \uses{def:mpu_admissible_index, def:mpu_canonical_form_ii}
    Let $\mathcal U$ be an MPU in canonical form II. If $k_0$ and $k$ are
    positive blocking lengths for which $\mathcal U_{k_0}$ and $\mathcal U_k$
    are simple, then both blocks give the same index. When $k\geq k_0$, their
    ranks satisfy
    \begin{align}
        r_k    & =d^{k-k_0}r_{k_0},    &
        \ell_k & =d^{k-k_0}\ell_{k_0}.
    \end{align}
    The value is unchanged by positive physical blocking or by choosing a
    different canonical-form-II presentation of the same tensor. This is the
    canonical-form-II representative case in
    \cite[Proposition~IV.2]{Cirac2017MPU} of index well-definedness.
    % Source lines: paper_v2.tex 690--704. The unrestricted source proposition
    % remains at index-well-defined above.
\end{lemma}

\begin{proof}\leanok
    \uses{thm:mpu_canonical_form_ii_positive_blocking, ThmFund1,
        thm:mpu_source_cut_rank_blocking_saturation,
        thm:mpu_specified_source_index_blocking_invariance,
        lem:mpu_iterated_physical_blocking}
    Each simple block retains canonical form II. The source-gate criterion
    gives $r_j\ell_j=d^{2j}$ at $j=k_0,k$, and therefore both ranks are
    positive. The source-rank blocking bounds, together with these two
    products, force the displayed equalities when $k\geq k_0$; the common
    factor $d^{k-k_0}$ cancels from the logarithmic difference. The same
    calculation compares either block with the one used to define the index.
    A further physical block is the corresponding direct block after the
    canonical reindexing, so its index is unchanged. Two canonical-form-II
    presentations of the same tensor are compared at a common simple block.
\end{proof}

\begin{lemma}[Index under physical adjunction in canonical form II]
    \label{thm:mpu_cfii_index_physical_adjoint}
    \lean{MPOTensor.IsMPUCanonicalFormII.index_physicalAdjointTensor}
    \leanok
    \uses{def:mpu_admissible_index, def:mpdo_physical_adjoint}
    Let $\mathcal U$ be an MPU in canonical form II. Its physical adjoint is
    again in canonical form II, and
    \begin{align}
        \operatorname{ind}(\mathcal U^\sharp)
         & =-\operatorname{ind}(\mathcal U).
    \end{align}
    This is the canonical-form-II representative case of the index-sign
    statement in the time-reversal discussion of
    \cite{Cirac2017MPU}. It does not assert that $\mathcal U$ and
    $\mathcal U^\sharp$ generate the same MPU.
    % Source: paper_v2.tex 1196--1207.
\end{lemma}

\begin{proof}\leanok
    \uses{thm:mpu_physical_adjoint_source_ranks,
        thm:mpu_double_layer_blocking, ThmFund1,
        def:mpu_admissible_index}
    For any positive simple block $\mathcal U_k$, physical adjunction
    exchanges its two source-cut ranks, so the rank-product criterion makes
    $(\mathcal U_k)^\sharp$ simple as well. At this block the right-rank
    formula for $\mathcal U^\sharp$ is the negative of the left-rank formula
    for $\mathcal U$:
    \begin{align}
        \log_2\frac{r((\mathcal U_k)^\sharp)}{d^k}
         & =\log_2\frac{\ell(\mathcal U_k)}{d^k}
        =-\operatorname{ind}(\mathcal U).
    \end{align}
    Both values are independent of the positive simple block chosen.
\end{proof}

\begin{definition}[Right- and left-shift tensors]
    \label{def:mpu_shift_tensors}
    \lean{MPOTensor.rightShiftTensor}
    \lean{MPOTensor.leftShiftTensor}
    \leanok
    \uses{def:mpdo_physical_adjoint}
    The right shift on $(\mathbb C^d)^{\otimes N}$ is
    \begin{align}
        T^{(N)}\lvert n_1,\ldots,n_N\rangle
         & =\lvert n_N,n_1,\ldots,n_{N-1}\rangle.
    \end{align}
    It has bond dimension $d$ and local matrices
    $A^{ij}=\lvert i\rangle\langle j\rvert$.
    The physical adjoint of $A$ defines the left-shift tensor.
    % Source lines: paper_v2.tex 1980--1987.
\end{definition}

\begin{lemma}[Shift tensor formulas]
    \label{thm:mpu_shift_tensor_formulas}
    \lean{MPOTensor.mpo_rightShiftTensor_apply}
    \lean{MPOTensor.mpo_rightShiftTensor}
    \lean{MPOTensor.mpo_leftShiftTensor}
    \lean{MPOTensor.rightShiftTensor_isMPU}
    \lean{MPOTensor.leftShiftTensor_isMPU}
    \leanok
    \uses{def:mpu_shift_tensors, def:mpu, thm:mpu_physical_adjoint_identities}
    For every $N>0$, the tensors in Definition~\ref{def:mpu_shift_tensors}
    generate $T^{(N)}$ and $T^{(N)\dagger}$, respectively.  Both tensors are
    matrix product unitaries.
\end{lemma}

\begin{proof}\leanok
    \uses{def:mpo_eval_word, def:mpo_rotate_config, lem:eval_word_of_fn_prod, thm:trace_eval_word_eq_sum_cyclic, def:mpu, thm:mpu_physical_adjoint_identities}
    Closing the virtual indices of $A^{ij}=\lvert i\rangle\langle j\rvert$
    leaves one nonzero virtual configuration.  If
    $\tau_n=\sigma_{n+1}$ for every $n$ modulo $N$, then
    \begin{align}
        \langle\sigma\vert T^{(N)}\vert\tau\rangle & =1.
    \end{align}
    Otherwise the same matrix entry is zero.  Hence the output configuration is
    $(\tau_{N-1},\tau_0,\ldots,\tau_{N-2})$, which fixes the permutation-matrix
    convention and gives the right shift.  Its permutation matrix is unitary.
    Physical adjunction gives the conjugate transpose, hence the left shift and
    its unitarity.
\end{proof}

\begin{definition}[Three topological-insulator shift MPUs]
    \label{threeMPU}
    \lean{MPOTensor.identityMPUTensor}
    \lean{MPOTensor.shiftExampleU₁}
    \lean{MPOTensor.shiftExampleU₂}
    \lean{MPOTensor.shiftExampleU₃}
    \leanok
    \uses{def:mpu_shift_tensors, def:mpu_independent_tensor_product, def:mpo_stacked_layers}
    On two $d$-dimensional spins per site, define
    \begin{align}
        U_1^{(N)} & =(\Id\otimes\Id)^{\otimes N},   \\
        U_2^{(N)} & =T^{(N)\dagger}\otimes T^{(N)}, \\
        U_3^{(N)} & =T^{(N)}\otimes T^{(N)\dagger}.
    \end{align}
    % Source lines: paper_v2.tex 1988--1995.
\end{definition}

\begin{lemma}[Finite-chain formulas for the three shift MPUs]
    \label{thm:threeMPU_chain_formulas}
    \lean{MPOTensor.mpo_identityMPUTensor}
    \lean{MPOTensor.identityMPUTensor_isMPU}
    \lean{MPOTensor.mpo_shiftExampleU₁}
    \lean{MPOTensor.mpo_shiftExampleU₂}
    \lean{MPOTensor.mpo_shiftExampleU₃}
    \lean{MPOTensor.shiftExampleU₁_isMPU}
    \lean{MPOTensor.shiftExampleU₂_isMPU}
    \lean{MPOTensor.shiftExampleU₃_isMPU}
    \leanok
    \uses{threeMPU, thm:mpu_shift_tensor_formulas, thm:mpu_finite_tensor_product, thm:mpu_tensor_product, def:mpu}
    The three tensors generate the operators displayed in
    Definition~\ref{threeMPU} on every nonempty chain, and each tensor is a
    matrix product unitary.
\end{lemma}

\begin{proof}\leanok
    \uses{def:mpo_stacked_layers, thm:mpu_finite_tensor_product, thm:mpu_shift_tensor_formulas, thm:mpu_tensor_product, def:mpu}
    The bond-one identity tensor generates the identity at every length.  The
    finite-chain operator of an independent tensor product is the Kronecker
    product of the two finite-chain operators, after the sitewise identification
    $(\mathbb C^d\otimes\mathbb C^d)^{\otimes N}
        \cong(\mathbb C^d)^{\otimes N}\otimes(\mathbb C^d)^{\otimes N}$.  Applying
    this identity to the right and left shifts gives the three displayed
    formulas.  The identity tensor is unitary, and tensor products preserve
    unitarity, so all three tensors are matrix product unitaries.
\end{proof}

\begin{lemma}[Index of the bond-one identity tensor]
    \label{thm:mpu_identity_index}
    \lean{MPOTensor.IsMPUCanonicalFormII.index_identityMPUTensor,
        MPOTensor.identityMPUTensorCanonicalFormII,
        MPOTensor.identityMPUTensor_index,
        MPOTensor.rightRank_identityMPUTensor,
        MPOTensor.leftRank_identityMPUTensor}
    \leanok
    \uses{def:mpu_admissible_index, threeMPU}
    Let $d>0$. The bond-one tensor generating the identity operator has a
    canonical-form-II presentation with fixed matrix $\rho=1$. For every
    canonical-form-II presentation of this tensor,
    \begin{align}
        \operatorname{ind}(\Id) & =0.
    \end{align}
    This follows from the index definition of
    \cite[Definition~IV.1]{Cirac2017MPU}; it is distinct from the source's
    later examples of three two-spin tensors with vanishing index.
    % Source: paper_v2.tex 681--704. The three examples are at 2037--2041.
\end{lemma}

\begin{proof}\leanok
    \uses{thm:mpu_canonical_form_ii_stabilization_converse, ThmFund1,
        def:mpu_admissible_index, def:mpu_physical_reindexing}
    In bond dimension one the normalized transfer matrix is the identity;
    the positive fixed matrix $\rho=1$ therefore gives a canonical-form-II
    presentation. Both source ranks of the identity tensor equal $d$, so the
    source-gate criterion makes it simple. Blocking once only relabels the
    physical alphabet. The defining logarithmic difference for any
    canonical-form-II presentation is consequently
    $\frac12(\log_2d-\log_2d)=0$.
\end{proof}

\begin{lemma}[Tensor-product additivity]
    \label{lem:mpu_index_tensor}
    \notready
    \discussion{6015}
    \uses{def:index, index-well-defined, thm:mpu_tensor_product_source_ranks}
    For MPU tensors under the source's standing canonical-form convention,
    \begin{align}
        \operatorname{ind}(\mathcal U\otimes\mathcal V)
         & =\operatorname{ind}(\mathcal U)+\operatorname{ind}(\mathcal V).
    \end{align}
    In \cite[Theorem~IV.6(ii)]{Cirac2017MPU}, this additivity identity is the unrestricted tensor-product
    statement.
    % Source lines: paper_v2.tex 824--833.
\end{lemma}

\begin{lemma}[Tensor-product additivity in canonical form II]
    \label{thm:mpu_admissible_index_tensor}
    \lean{MPOTensor.IsMPUCanonicalFormII.index_tensorProduct,
        MPOTensor.IsMPUCanonicalFormII.index_tensorProduct_identityMPUTensor}
    \leanok
    % **Scope restriction (canonical representatives):** the two factors are
    % already in canonical form II. No comparison with an arbitrary unreduced
    % representative is claimed.
    % Documented in docs/paper-gaps/mpu_canonical_form_full_support.tex.
    \uses{def:mpu_admissible_index, def:mpu_canonical_form_ii,
        def:mpu_independent_tensor_product, threeMPU}
    Let $\mathcal U$ and $\mathcal V$ be MPUs in canonical form II. Their
    independent tensor product is again in canonical form II, and
    \begin{align}
        \operatorname{ind}(\mathcal U\boxtimes\mathcal V)
         & =\operatorname{ind}(\mathcal U)+\operatorname{ind}(\mathcal V).
    \end{align}
    In particular, if $n>0$ and $\Id_n$ is the bond-one tensor generating
    the identity on a physical space of dimension $n$, then
    \begin{align}
        \operatorname{ind}(\mathcal U\boxtimes\Id_n)
         & =\operatorname{ind}(\mathcal U).
    \end{align}
    The bridge to physical identity-ancilla attachment is
    Lemma~\ref{thm:mpu_physical_ancilla_index}.
    In \cite[Theorem~IV.6(ii)]{Cirac2017MPU}, this identity is the canonical-form-II representative case of the tensoring
    claim.
    % Source lines: paper_v2.tex 824--833.
\end{lemma}

\begin{proof}\leanok
    \uses{thm:mpu_canonical_form_ii_tensor_product,
        lem:mpu_common_simple_blocking_finite,
        thm:mpu_block_tensor_product,
        thm:mpu_tensor_product_source_ranks,
        thm:mpu_specified_source_index_identities,
        prop:mpu_admissible_index_well_defined,
        thm:mpu_identity_index}
    At a common positive simple blocking length $k$ for $\mathcal U$,
    $\mathcal V$, and their tensor product, blocking commutes with the independent tensor product
    up to a physical relabeling, so their source ranks satisfy
    \begin{align}
        r((\mathcal U\boxtimes\mathcal V)_k)
         & =r(\mathcal U_k)r(\mathcal V_k),       \\
        \ell((\mathcal U\boxtimes\mathcal V)_k)
         & =\ell(\mathcal U_k)\ell(\mathcal V_k).
    \end{align}
    Taking the half-difference of the base-two logarithms gives additivity;
    independence of $k$ identifies these values with the three indices.
    The bond-one identity has index zero by
    Lemma~\ref{thm:mpu_identity_index}, giving the last assertion.
\end{proof}

\begin{lemma}[Tensor-product additivity of the representative index]
    \label{thm:mpu_representative_index_tensor_product}
    \lean{MPOTensor.IsMPU.index_tensorProduct,
        MPOTensor.IsMPU.index_identityMPUTensor,
        MPOTensor.IsMPU.index_tensorProduct_identityMPUTensor}
    \leanok
    \uses{def:mpu_reduced_representative_index,
        def:mpu_independent_tensor_product, thm:mpu_tensor_product, threeMPU}
    Let $U$ and $V$ be positive-dimensional MPU tensors, possibly with
    different physical and bond dimensions. Their independent tensor product
    satisfies
    \begin{align}
        \operatorname{ind}(U\boxtimes V)
         & =\operatorname{ind}(U)+\operatorname{ind}(V).
    \end{align}
    The bond-one identity tensor $\Id_n$ has index zero for every $n>0$;
    hence $\operatorname{ind}(U\boxtimes\Id_n)=\operatorname{ind}(U)$.
    In \cite[Theorem~IV.6(ii)]{Cirac2017MPU}, this additivity is the reduced-representative form of the tensoring
    assertion.
\end{lemma}

\begin{proof}\leanok
    \uses{thm:mpu_reduced_cfii_representative,
        thm:mpu_reduced_representative_index_invariance,
        thm:mpu_admissible_index_tensor, thm:mpu_identity_index,
        thm:mpu_finite_tensor_product}
    For reduced canonical-form-II representatives $U_{\mathrm{red}}$ and
    $V_{\mathrm{red}}$, Lemma~\ref{thm:mpu_finite_tensor_product} gives, for every $N>0$,
    the periodic factorization
    \begin{align}
        (U_{\mathrm{red}}\boxtimes V_{\mathrm{red}})^{(N)}
         & \cong U_{\mathrm{red}}^{(N)}\otimes V_{\mathrm{red}}^{(N)}
        = U^{(N)}\otimes V^{(N)}
        \cong (U\boxtimes V)^{(N)}.
    \end{align}
    Representative independence and canonical tensor-product additivity give
    \begin{align}
        \operatorname{ind}(U\boxtimes V)
         & =\operatorname{ind}(U_{\mathrm{red}}\boxtimes V_{\mathrm{red}}) \\
         & =\operatorname{ind}(U_{\mathrm{red}})
        +\operatorname{ind}(V_{\mathrm{red}}).
    \end{align}
    The identity tensor itself is a canonical representative of index zero,
    and the ancilla assertion follows from additivity.
\end{proof}

\begin{lemma}[Physical identity ancillas and the representative index]
    \label{thm:mpu_physical_ancilla_index}
    \lean{MPOTensor.mpo_tensorProduct_identityMPUTensor,
        MPOTensor.mpo_tensorPhysicalId,
        MPOTensor.IsMPU.index_tensorPhysicalId}
    \leanok
    \uses{def:equivalent-tensors, def:mpu_independent_tensor_product, threeMPU,
        def:mpu_reduced_representative_index}
    Let $U$ be a tensor and let $U^{(n)}=U\otimes\Id_n$ denote attachment
    of an identity ancilla of dimension $n$ at each physical site, keeping
    the bond dimension fixed. For every chain length $N$, its periodic
    operator equals that of the independent tensor product with the
    bond-one identity tensor:
    \begin{align}
        (U^{(n)})^{(N)}=(U\boxtimes\Id_n)^{(N)}
        \cong U^{(N)}\otimes I_{n^N}.
    \end{align}
    Here the identification groups the original physical coordinates and
    the ancilla coordinates separately. If $U$ is a positive-dimensional
    MPU and $n>0$, then
    \begin{align}
        \operatorname{ind}(U^{(n)})=\operatorname{ind}(U).
    \end{align}
    % Source: paper_v2.tex, def:equivalent-tensors, lines 706--720;
    % IndexTh(ii), lines 824--845.
\end{lemma}

\begin{proof}\leanok
    \uses{thm:mpu_finite_tensor_product, thm:threeMPU_chain_formulas,
        thm:mpu_reduced_representative_index_invariance,
        thm:mpu_representative_index_tensor_product}
    Both constructions give the displayed Kronecker product in the same
    sitewise coordinates. Equality of the periodic operators identifies
    their representative indices, and the bond-one identity factor has
    index zero.
\end{proof}

\begin{lemma}[Source ranks under rectangular reduction]
    \label{lem:mpu_rectangular_reduction_source_ranks}
    \lean{MPOTensor.sourceCutM₁_rectangularSandwich,
        MPOTensor.sourceCutM₂_rectangularSandwich,
        MPOTensor.rightRank_rectangularSandwich_le,
        MPOTensor.leftRank_rectangularSandwich_le,
        MPOTensor.eq_rectangularSandwich_of_isReduction,
        MPOTensor.rightRank_le_of_isReduction,
        MPOTensor.leftRank_le_of_isReduction,
        MPOTensor.isReduction_blockTensor,
        MPOTensor.rightRank_blockTensor_le_of_isReduction,
        MPOTensor.leftRank_blockTensor_le_of_isReduction}
    \leanok
    \uses{def:mpu_source_matrix_one, def:mpu_source_matrix_two,
        def:mpu_right_rank, def:mpu_left_rank, def:mps_rectangular_reduction,
        def:mpdo_physical_blocking}
    Let $\mathcal U$ be a matrix product operator tensor of physical
    dimension $d$ and bond dimension $D$. For arbitrary matrices
    $A\in\mathbb C^{E\times D}$ and $B\in\mathbb C^{D\times E}$, put
    $\mathcal V^{ij}=A\mathcal U^{ij}B$. Then
    \begin{align}
        M_1(\mathcal V)
                         & =(I_d\otimes B^{\mathsf T})M_1(\mathcal U)
        (A^{\mathsf T}\otimes I_d),                                      \\
        M_2(\mathcal V)
                         & =(A\otimes I_d)M_2(\mathcal U)(I_d\otimes B), \\
        r[\mathcal V]    & \leq r[\mathcal U],                           \\
        \ell[\mathcal V] & \leq\ell[\mathcal U].
    \end{align}
    In particular, a rectangular reduction cannot increase either source
    rank. No equality of ranks is asserted, and no unitarity, simplicity,
    or canonical-form assumption is needed.
    If $\mathcal V$ is a rectangular reduction of $\mathcal U$, then for
    every common blocking length $L$,
    \begin{align}
        r[\mathcal V_L]    & \leq r[\mathcal U_L],   \\
        \ell[\mathcal V_L] & \leq\ell[\mathcal U_L].
    \end{align}
    These identities follow directly from the source cuts of
    \cite[Section~III, definition of $r,\ell$]{Cirac2017MPU} and permit
    upper bounds for composition to pass to a reduced representative.
\end{lemma}

\begin{proof}\leanok
    Entrywise expansion of the matrix products in the two source cuts gives both formulas.
    The first cut places the right virtual coordinate on its row side,
    which accounts for the two transposes. The second cut keeps the
    ordinary order of the virtual coordinates. Multiplication on either
    side cannot increase matrix rank, proving both inequalities.
    For a rectangular reduction, its single-letter word identity gives
    $\mathcal V^{ij}=A\mathcal U^{ij}B$.
    Its all-word identity gives the same relation for the blocked tensors,
    with the same matrices $A,B$, so the inequalities apply at every length.
\end{proof}

\begin{lemma}[Four-site composition source-rank bounds]
    \label{lem:mpu_four_site_composition_source_rank_bounds}
    \lean{MPOTensor.rightRank_blockTensor_mulTensor_four_le,
        MPOTensor.leftRank_blockTensor_mulTensor_four_le}
    \leanok
    \uses{def:mpdo_operator_product, def:mpdo_physical_blocking, def:mpu_right_rank, def:mpu_left_rank}
    Let $\mathcal U$ and $\mathcal V$ be arbitrary matrix product operator
    tensors of common physical dimension $d$. For their product tensor,
    four-site blocking satisfies
    \begin{align}
        r[(\mathcal U\mathbin{\cdot}\mathcal V)_4]
         & \leq d^2 r[\mathcal U]r[\mathcal V],       \\
        \ell[(\mathcal U\mathbin{\cdot}\mathcal V)_4]
         & \leq d^2 \ell[\mathcal U]\ell[\mathcal V].
    \end{align}
    No unitarity or simplicity assumption is required.
    % Source lines: paper_v2.tex 837--845; issue 6999.
\end{lemma}

\begin{proof}\leanok
    \uses{thm:mpu_source_v_open_two_site,
        thm:mpu_source_u_open_two_site,
        thm:mpdo_blocking_product,
        thm:mpdo_block_two_as_positive_blocking,
        thm:mpu_source_ranks_physical_reindexing,
        lem:mpu_iterated_physical_blocking}
    The canonical physical reindexing identifies iterated two-site
    blocking with direct four-site blocking. For the right cut, writing each
    two-site block in the $v$--$Y_1$--$Y_2$ form and grouping $v$ with $Y_2$ on one
    side of the cut while retaining $Y_1$ on the other shows that the cut carries
    one physical leg and one right source leg from each tensor. Thus, for
    suitable rectangular matrices $A_{1,4}$ and $B_{1,4}$,
    \begin{align}
        M_1((\mathcal U\mathbin{\cdot}\mathcal V)_4) & =A_{1,4}B_{1,4},
    \end{align}
    where the columns of $A_{1,4}$ and the rows of $B_{1,4}$ are indexed by
    \begin{align}
        \mathcal K_1
         & =(\mathbb C^d\otimes\mathbb C^{r[\mathcal U]})
        \otimes
        (\mathbb C^d\otimes\mathbb C^{r[\mathcal V]}).
    \end{align}
    Consequently,
    \begin{align}
        r[(\mathcal U\mathbin{\cdot}\mathcal V)_4]
         & =\operatorname{rank}(A_{1,4}B_{1,4})
        \leq\dim(\mathcal K_1)
        =d^2r[\mathcal U]r[\mathcal V].
    \end{align}
    For the left cut, write each two-site block in the
    $X_2$--$u$--$X_1$ form. Retain $X_2$ on one side and group $u$ with $X_1$
    on the other. The cut carries one physical leg and one left source leg from
    each tensor. For suitable rectangular matrices $A_{2,4}$ and $B_{2,4}$,
    \begin{align}
        M_2((\mathcal U\mathbin{\cdot}\mathcal V)_4) & =A_{2,4}B_{2,4},
    \end{align}
    where the intermediate space is
    \begin{align}
        \mathcal K_2
         & =(\mathbb C^d\otimes\mathbb C^{\ell[\mathcal U]})
        \otimes
        (\mathbb C^d\otimes\mathbb C^{\ell[\mathcal V]}).
    \end{align}
    Hence
    \begin{align}
        \ell[(\mathcal U\mathbin{\cdot}\mathcal V)_4]
         & =\operatorname{rank}(A_{2,4}B_{2,4})
        \leq\dim(\mathcal K_2)
        =d^2\ell[\mathcal U]\ell[\mathcal V].
    \end{align}
    Both source-cut ranks are unchanged by the canonical physical relabeling.
\end{proof}

\begin{lemma}[Composition additivity]
    \label{lem:mpu_index_composition}
    \notready
    \discussion{6016}
    \uses{SF, def:index, index-well-defined, lem:mpu_four_site_composition_source_rank_bounds}
    Let $\mathcal W$ generate the composition of the MPUs generated by
    $\mathcal U$ and $\mathcal U'$. Under the source's standing canonical-form
    convention,
    \begin{align}
        \operatorname{ind}(\mathcal W)
         & =\operatorname{ind}(\mathcal U)+\operatorname{ind}(\mathcal U').
    \end{align}
    In \cite[Theorem~IV.6(ii)]{Cirac2017MPU}, this additivity formula is the unrestricted composition
    statement.
    % Source lines: paper_v2.tex 833--845.
\end{lemma}

\begin{lemma}[Reduction to a canonical representative]
    \label{thm:mpu_canonical_representative_reduction}
    \lean{MPOTensor.exists_isReduction_of_mpo_eq_cfii}
    \leanok
    \uses{def:mpu_canonical_form_ii, def:mps_rectangular_reduction,
        def:mpo_family}
    Let $\mathcal U$ and $\mathcal V$ have the same physical dimension and
    bond dimensions $D$ and $E$, respectively. Suppose $\mathcal V$ is in
    canonical form II and their periodic operators agree at every positive
    length. Then there are matrices $A\in\mathbb C^{E\times D}$ and
    $B\in\mathbb C^{D\times E}$ such that, for every word $w$ in the paired
    physical alphabet,
    \begin{align}
        A\mathcal U^w B & =\mathcal V^w.
    \end{align}
    The empty word is included, so $AB=I_E$.
\end{lemma}

\begin{proof}\leanok
    \uses{thm:mps_equality_reduction_exists, def:mpu_canonical_form_ii,
        thm:mpu_canonical_form_ii_normal}
    The normalized flattening of $\mathcal V$ is normal. Multiplication by
    the nonzero scalar $\sqrt d$ preserves normality, so its unnormalized
    flattening is normal as well. Equality of the periodic operators gives
    equality of all positive-length word traces. The equality-case
    rectangular reduction theorem now gives $A$ and $B$.
\end{proof}

\begin{lemma}[Saturation of the composition rank bounds]
    \label{thm:mpu_composition_rank_saturation}
    \lean{MPOTensor.IsMPUCanonicalFormII.compositionRanks_of_isReduction}
    \leanok
    \uses{def:mpu_canonical_form_ii, def:mpu_simple, def:mps_rectangular_reduction,
        def:mpu_right_rank, def:mpu_left_rank}
    Let $\mathcal U,\mathcal V,\mathcal R$ be canonical-form-II tensors of
    physical dimension $d$. Suppose $\mathcal U$ and $\mathcal V$ are simple,
    $\mathcal R_4$ is simple, and $\mathcal R$ is a rectangular reduction of
    $\mathcal U\mathbin{\cdot}\mathcal V$. Then
    \begin{align}
        r[\mathcal R_4]    & =d^2r[\mathcal U]r[\mathcal V],       \\
        \ell[\mathcal R_4] & =d^2\ell[\mathcal U]\ell[\mathcal V].
    \end{align}
\end{lemma}

\begin{proof}\leanok
    \uses{lem:mpu_rectangular_reduction_source_ranks,
        lem:mpu_four_site_composition_source_rank_bounds,
        prop:mpu_admissible_index_well_defined, ThmFund1}
    Rectangular reduction and the four-site composition bounds give the
    corresponding upper bounds. Simplicity gives
    \begin{align}
        r[\mathcal U]\ell[\mathcal U]
                                          & =r[\mathcal V]\ell[\mathcal V]=d^2, \\
        r[\mathcal R_4]\ell[\mathcal R_4] & =d^8.
    \end{align}
    The product of the two upper bounds is therefore $d^8$. Since both ranks
    of $\mathcal R_4$ are positive, neither inequality can be strict.
\end{proof}

\begin{lemma}[Composition additivity of the representative index]
    \label{thm:mpu_admissible_index_composition}
    \lean{MPOTensor.IsMPU.index_mulTensor,
        MPOTensor.IsMPUCanonicalFormII.index_eq_add_of_mpo_eq_mulTensor}
    \leanok
    \uses{def:mpu_reduced_representative_index, def:mpdo_operator_product,
        def:mpo_family}
    Let $\mathcal U$ and $\mathcal V$ be MPUs with a common positive physical
    dimension and positive bond dimensions. For the index defined through a reduced canonical-form-II
    representative,
    \begin{align}
        \operatorname{ind}(\mathcal U\mathbin{\cdot}\mathcal V)
         & =\operatorname{ind}(\mathcal U)+\operatorname{ind}(\mathcal V).
    \end{align}
    In particular, if $\mathcal U,\mathcal V,\mathcal R$ are in canonical form
    II and the positive-length periodic operators of $\mathcal R$ equal the
    products of those of $\mathcal U$ and $\mathcal V$, then their canonical
    indices satisfy the same identity, following in \cite[Theorem~IV.6(ii)]{Cirac2017MPU} the rank-saturation
    proof.
    % The unreduced source-cut formula in def:index is not asserted here.
\end{lemma}

\begin{proof}\leanok
    \uses{thm:mpu_reduced_cfii_representative,
        thm:mpu_reduced_representative_index_invariance,
        thm:mpu_canonical_representative_reduction,
        lem:mpu_common_simple_blocking_finite,
        thm:mpu_composition_rank_saturation,
        thm:mpu_specified_source_index_identities,
        prop:mpu_admissible_index_well_defined}
    For reduced canonical representatives of the two factors and their
    composition, the positive-length periodic operators still satisfy the
    composition identity. Lemma~\ref{thm:mpu_canonical_representative_reduction} gives a rectangular
    reduction from the product tensor to its representative $\mathcal R$.
    At a common positive blocking length $k$ where the three
    representatives are simple, with $d_k=d^k$, the rectangular reduction
    persists under blocking. Rank saturation gives
    \begin{align}
        r[\mathcal R_{4k}]    & =d_k^2r[\mathcal U_k]r[\mathcal V_k],       \\
        \ell[\mathcal R_{4k}] & =d_k^2\ell[\mathcal U_k]\ell[\mathcal V_k].
    \end{align}
    Subtracting the logarithms cancels $d_k^2$ and adds the two indices.
    Independence of positive blocking and of the chosen canonical
    representatives gives the asserted identity.
\end{proof}

\begin{definition}[Strict equivalence]
    \label{def:strictly-equivalent-tensors}
    \lean{MPOTensor.StrictlyEquivalent, MPOTensor.padBond,
        MPOTensor.StrictlyEquivalent.of_fixedBond}
    \leanok
    \uses{def:mpu, def:mpu_physical_reindexing, def:mpu_canonical_form}
    Let $\mathcal U$ and $\mathcal V$ be canonical-form MPU tensors with
    physical dimensions $d_a$ and $d_b$ and virtual bond dimensions $D_a$ and
    $D_b$. For a tensor $\mathcal U$ of bond dimension $D$ and $D'\geq D$,
    write $\mathcal U_{\uparrow D'}$ for the tensor of
    bond dimension $D'$ whose letters are those of $\mathcal U$ in the upper
    left $D\times D$ corner and zero on the adjoined directions; adjoining
    unused bond directions changes no periodic operator of positive length.
    Given an equality $h\colon d_a=d_b$, reindex $\mathcal U$ from $d_a$ to
    $d_b$ along $h^{-1}$. The tensors are strictly equivalent with respect to
    $h$ if there are $D'\geq\max(D_a,D_b)$ and a continuous path
    $\mathcal W(p)$, $p\in[0,1]$, of MPU tensors of bond dimension $D'$ such
    that
    \begin{align}
        \mathcal W(0) & =\operatorname{reind}_{h^{-1}}(\mathcal U_{\uparrow D'}), \\
        \mathcal W(1) & =\mathcal V_{\uparrow D'}.
    \end{align}
    Only the canonical-form requirement is relaxed along the path
    \cite[lines~708--714]{Cirac2017MPU}. When $D_a=D_b=D$, a path of MPU
    tensors of bond dimension $D$ from $\mathcal U$ to $\mathcal V$ establishes
    strict equivalence with $D'=D$.
    % **Local fix:** the source does not say in which bond space the path runs
    % when $D_a\neq D_b$; its canonical form admits the adjoined directions and
    % ours does not. Documented in docs/paper-gaps/mpu_equivalence_fixed_bond.tex.
\end{definition}

\begin{lemma}[Normalized transfer data under identity ancillas]
    \label{thm:mpu_identity_ancilla_reduced_cfii}
    \lean{MPOTensor.normalizedFlattening_tensorPhysicalId,
        MPOTensor.transferMap_normalizedFlattening_tensorPhysicalId,
        MPOTensor.tensorPhysicalIdCFIIData,
        MPOTensor.hasFullSupport_tensorPhysicalIdCFIIData}
    \leanok
    \uses{def:mpu, def:mpu_normalized_flattening, def:cpsv_cfii, def:mpu_full_support}
    Let $x>0$. After adjoining an identity ancilla of dimension $x$, the
    normalized flattening is the physical reindexing of the original flattening
    with a normalized diagonal ancilla alphabet. Its transfer map is exactly the
    original transfer map. Consequently, chosen canonical-form-II data with full
    support determine canonical-form-II data with full support for
    $\mathcal U^{(x)}=\mathcal U\otimes\Id_x$.

    The product-index equivalence has the orientation
    \begin{align}
        ((i,a),(j,b)) & \longmapsto((i,j),(a,b)).
    \end{align}
    This result preserves the normalized transfer data.
    % Source lines: paper_v2.tex 706--724; normalized transfer lines 336--340.
\end{lemma}

\begin{proof}\leanok
    \uses{def:mpu_physical_reindexing}
    Only the $x$ diagonal ancilla letters are nonzero, each with coefficient
    $x^{-1/2}$. Hence their total transfer contribution is
    \begin{align}
        x\left(x^{-1/2}\right)^2 & =1.
    \end{align}
    Thus the transfer map and each diagonal positive fixed matrix are unchanged.
    The block dimensions, weights, and ambient coisometry are unchanged as well,
    so full support follows from the same block-dimension sum. The displayed
    physical reindexing gives the normalized flattening of the enlarged MPO.
\end{proof}

\begin{definition}[Equivalence after blocking and ancillas]
    \label{def:equivalent-tensors}
    \lean{MPOTensor.Equivalent}
    \leanok
    \uses{def:mpu, def:strictly-equivalent-tensors, lem:mpu_blocking}
    Let $\mathcal U$ and $\mathcal V$ be MPU
    tensors with physical dimensions $d_a$ and $d_b$. They are equivalent if
    there are a
    positive blocking length $k$ and coprime positive integers $p_a,p_b$ with
    \begin{align}
        p_a d_a & =p_b d_b
    \end{align}
    such that $(\mathcal U^{(p_a)})_k$ and $(\mathcal V^{(p_b)})_k$ are strictly
    equivalent. Here the positive-dimensional identity ancillas are attached
    before the common blocking, and
    \begin{align}
        \mathcal U^{(p)} & =\mathcal U\otimes\Id_p.
    \end{align}
    As in \cite[lines~716--724]{Cirac2017MPU}, canonical form of the blocked
    tensors with ancilla is a condition of the definition.
\end{definition}

\begin{definition}[Admissible equivalence and path datum]
    \label{def:mpu_admissible_equivalence_datum}
    \notready
    % **Scope restriction (no existence or preservation theorem):** this datum
    % records every reduced source hypothesis needed along a chosen path. It does
    % not assert preservation under blocking, ancillas, products, composition, or
    % continuous variation. Documented in
    % docs/paper-gaps/mpu_canonical_form_full_support.tex.
    \uses{def:strictly-equivalent-tensors, def:equivalent-tensors, def:mpu_reduced_full_support_source_datum}
    An admissible strict-equivalence datum for a path
    $[0,1]\ni p\mapsto\mathcal W(p)$ consists of a common positive blocking
    length $m$ such that the blocked path tensor $\mathcal W(p)_m$ itself is
    simple and carries a reduced full-support source datum for every $p$. The
    blocked tensors and their source data vary continuously with $p$, after
    fixed identifications of the finite-dimensional spaces involved. No
    equivalent or reduced representative may replace $\mathcal W(p)_m$. In
    particular, the blocked endpoint tensors themselves carry their own source data.

    When the path is used under a specified symmetry $\mathcal S$, the datum also
    includes continuously varying reduced full-support source data for the fixed
    family of comparison paths
    \begin{align}
        \overline{\mathcal W(p)_m},\qquad
        \mathcal W(p)_m^\sharp,\qquad
        \mathcal W(p)_m^{\mathsf T},\qquad
        \mathcal S[\mathcal W(p)_m].
    \end{align}
    Every tensor in this family is simple. This requirement is part of the datum
    independently of which comparisons a later proof uses; no preservation under
    conjugation, physical adjoint, transposition, or $\mathcal S$ is assumed.

    Two tensors are admissibly strictly equivalent if their strict equivalence
    is realized by such a path. They are admissibly equivalent if the blocked,
    ancilla-enlarged endpoints in Definition~\ref{def:equivalent-tensors} are
    joined by an admissible strict-equivalence path. The source arguments
    motivating these additional path hypotheses are Proposition~IV.5,
    Theorem~IV.6, Corollary~IV.7, and the lemma following the
    symmetry-equivalence definitions in \cite{Cirac2017MPU}; the datum itself is
    introduced here and has no source-labelled counterpart.
    % Source lines: paper_v2.tex 733--910 and 1366--1380.
\end{definition}

\begin{lemma}[Generic two-projection angle block]
    \label{thm:mpu_two_projection_angle_block}
    \lean{LinearMap.IsSymmetricProjection.rangeCompression_isSymmetric,
        LinearMap.IsSymmetricProjection.angleDefect_block,
        LinearMap.IsSymmetricProjection.exists_orthonormal_angle_block}
    \leanok

    Let $P$ and $Q$ be orthogonal projections on a complex inner-product space.
    The compression of $Q$ to $\operatorname{ran}P$ is self-adjoint. If a unit
    vector $x\in\operatorname{ran}P$ satisfies
    \begin{align}
        PQx & =\mu x, & 0 & <\mu<1,
    \end{align}
    then there is a unit vector $w\perp x$ such that
    \begin{align}
        Px & =x,                             & Pw & =0, \\
        Qx & =\mu x+\sqrt{\mu(1-\mu)}\,w,    &
        Qw & =\sqrt{\mu(1-\mu)}\,x+(1-\mu)w.
    \end{align}
    Thus $\operatorname{span}\{x,w\}$ is the two-dimensional angle block
    associated with the eigenvalue $\mu\in(0,1)$.
    % Formalization status: endpoint spaces and global Jordan assembly remain not ready.
    % Source: paper_v2.tex, Proposition IV.5 (prop:continuity-index), lines 756--764.
    % Local correction: line 762 repeats P where Q is intended; see
    % docs/paper-gaps/mpu_two_projection_projector_typo.tex.
\end{lemma}

\begin{proof}\leanok
    For $d=Qx-\mu x$, symmetry and idempotence give
    \begin{align}
        Pd               & =0,                    & \langle x,d\rangle & =0, &
        \lVert d\rVert^2 & =\mu(1-\mu),                                        \\
        Qd               & =\mu(1-\mu)x+(1-\mu)d.
    \end{align}
    Since $0<\mu<1$, the vector
    $w=d/\sqrt{\mu(1-\mu)}$ is well defined and has the stated properties.
\end{proof}

\begin{definition}[Compression spectral data]
    \label{def:mpu_two_projection_compression_spectral_data}
    \lean{LinearMap.IsSymmetricProjection.rangeCompressionEigenvalues,
        LinearMap.IsSymmetricProjection.rangeCompressionEigenbasis}
    \leanok
    \uses{thm:mpu_two_projection_angle_block}
    For orthogonal projections $P$ and $Q$ on a finite-dimensional complex
    inner-product space, let $(x_i)_i$ be an orthonormal eigenbasis of the
    compression of $Q$ to $\operatorname{ran}P$, with ordered real eigenvalues
    $(\mu_i)_i$.
    % Source lines: paper_v2.tex 759--762.
\end{definition}

\begin{lemma}[Compression-spectrum endpoint and angle-block trichotomy]
    \label{thm:mpu_two_projection_compression_spectrum}
    \lean{LinearMap.IsSymmetricProjection.rangeCompression_apply_eigenbasis,
        LinearMap.IsSymmetricProjection.rangeCompressionEigenbasis_apply_self,
        LinearMap.IsSymmetricProjection.rangeCompressionEigenbasis_apply_compression,
        LinearMap.IsSymmetricProjection.rangeCompressionEigenvalues_mem_unitInterval,
        LinearMap.IsSymmetricProjection.rangeCompressionEigenvalues_eq_zero_or_eq_one_or_exists_angle_block}
    \leanok
    \uses{def:mpu_two_projection_compression_spectral_data, thm:mpu_two_projection_angle_block}
    Let $P$ and $Q$ be orthogonal projections on a finite-dimensional complex
    inner-product space, with compression spectral data $(x_i,\mu_i)$. Then, for
    every $i$, $\mu_i\in[0,1]$ and exactly one of the following alternatives
    holds:
    \begin{enumerate}
        \item $\mu_i=0$ and $Qx_i=0$;
        \item $\mu_i=1$ and $Qx_i=x_i$;
        \item $0<\mu_i<1$, and $x_i$ belongs to the two-dimensional angle block
              described in Lemma~\ref{thm:mpu_two_projection_angle_block}.
    \end{enumerate}
    This theorem classifies the compression eigenvectors.
    % Source lines: paper_v2.tex 759--762.
    % Scope: no pairwise orthogonality of angle blocks, no ambient orthogonal
    % direct sum, and no commutation on the remaining summand.
\end{lemma}

\begin{proof}\leanok
    \uses{def:mpu_two_projection_compression_spectral_data, thm:mpu_two_projection_angle_block}
    The finite-dimensional spectral theorem applied to the self-adjoint compression
    of $Q$ to $\operatorname{ran}P$ diagonalizes $PQP|_{\operatorname{ran}P}$. For a unit compression eigenvector $x$ with
    eigenvalue $\mu$, the angle-defect identity gives
    \begin{align}
        \lVert Qx-\mu x\rVert^2 & =\mu(1-\mu),
    \end{align}
    hence $0\leq\mu\leq1$. At $\mu=0$ or $1$ the defect vanishes, giving the two
    common one-dimensional blocks. For $0<\mu<1$, apply
    Lemma~\ref{thm:mpu_two_projection_angle_block}.
\end{proof}

\begin{definition}[Compression defect block]
    \label{def:mpu_two_projection_angle_defect_block}
    \lean{LinearMap.IsSymmetricProjection.angleDefect,
        LinearMap.IsSymmetricProjection.angleDefectBlock}
    \leanok
    \uses{def:mpu_two_projection_compression_spectral_data}
    For compression spectral data $(x_i,\mu_i)$, define
    \begin{align}
        d_i & =Qx_i-\mu_i x_i,                 \\
        B_i & =\operatorname{span}\{x_i,d_i\}.
    \end{align}
\end{definition}

\begin{lemma}[Orthogonality of compression angle blocks]
    \label{thm:mpu_two_projection_angle_block_orthogonality}
    \lean{LinearMap.IsSymmetricProjection.rangeCompressionEigenbasis_inner_angleDefect_eq_zero,
        LinearMap.IsSymmetricProjection.angleDefect_inner_rangeCompressionEigenbasis_eq_zero,
        LinearMap.IsSymmetricProjection.angleDefect_inner_angleDefect,
        LinearMap.IsSymmetricProjection.angleDefect_inner_angleDefect_eq_zero,
        LinearMap.IsSymmetricProjection.finrank_angleDefectBlock_le_two,
        LinearMap.IsSymmetricProjection.map_angleDefectBlock_le_first_projection,
        LinearMap.IsSymmetricProjection.map_angleDefectBlock_le_second_projection,
        LinearMap.IsSymmetricProjection.angleDefectBlock_isOrtho}
    \leanok
    \uses{def:mpu_two_projection_angle_defect_block}
    For all indices $i,j$,
    \begin{align}
        \langle x_i,d_j\rangle & =0,                                    \\
        \langle d_i,x_j\rangle & =0,                                    \\
        \langle d_i,d_j\rangle
                               & =\mu_j(1-\mu_j)\langle x_i,x_j\rangle.
    \end{align}
    In particular, $\langle d_i,d_j\rangle=0$ when $i\ne j$, including when
    $\mu_i=\mu_j$. The block $B_i$ has dimension at most two, is invariant under
    both $P$ and $Q$, and satisfies $B_i\perp B_j$ for distinct indices.
    % This does not assert an ambient direct sum or any property of the remaining
    % orthogonal complement.
    % Source lines: paper_v2.tex 759--762.
\end{lemma}

\begin{proof}\leanok
    \uses{def:mpu_two_projection_angle_defect_block, thm:mpu_two_projection_compression_spectrum, thm:mpu_two_projection_angle_block}
    Since $Px_i=x_i$, symmetry of $P$ and the compression eigenvalue equation give
    \begin{align}
        \langle x_i,d_j\rangle
         & =\langle x_i,Qx_j\rangle-\mu_j\langle x_i,x_j\rangle     \\
         & =\langle Px_i,Qx_j\rangle-\mu_j\langle x_i,x_j\rangle    \\
         & =\langle x_i,PQx_j\rangle-\mu_j\langle x_i,x_j\rangle=0.
    \end{align}
    Hence $\langle d_i,x_j\rangle=0$ by conjugate symmetry. For the exact defect
    pairing, symmetry of $Q$ and the local action formula give
    \begin{align}
        Qd_j & =\mu_j(1-\mu_j)x_j+(1-\mu_j)d_j,                      \\
        \langle d_i,d_j\rangle
             & =\langle Qx_i,d_j\rangle-\mu_i\langle x_i,d_j\rangle, \\
        \langle Qx_i,d_j\rangle-\mu_i\langle x_i,d_j\rangle
             & =\langle x_i,Qd_j\rangle,                             \\
        \langle x_i,Qd_j\rangle
             & =\mu_j(1-\mu_j)\langle x_i,x_j\rangle
        +(1-\mu_j)\langle x_i,d_j\rangle,                            \\
        \langle d_i,d_j\rangle
             & =\mu_j(1-\mu_j)\langle x_i,x_j\rangle.
    \end{align}
    Orthogonality of the generated spans follows by linearity. The action formulas
    \begin{align}
        Px_i & =x_i,                           \\
        Pd_i & =0,                             \\
        Qx_i & =\mu_i x_i+d_i,                 \\
        Qd_i & =\mu_i(1-\mu_i)x_i+(1-\mu_i)d_i
    \end{align}
    show that both projections preserve
    $B_i=\operatorname{span}\{x_i,d_i\}$. Its dimension bound follows from
    its two generators.
\end{proof}
\begin{definition}[Defect-block sum and complementary summand]
    \label{def:mpu_two_projection_angle_defect_sum}
    \lean{LinearMap.IsSymmetricProjection.angleDefectSum,
        LinearMap.IsSymmetricProjection.angleDefectComplement}
    \leanok
    \uses{def:mpu_two_projection_angle_defect_block}
    For the compression defect blocks $(B_i)_i$, define
    \begin{align}
        B & =\bigvee_i B_i, & K & =B^\perp.
    \end{align}
    The index $i$ ranges over every vector in the chosen compression eigenbasis,
    including repeated eigenvalues and the endpoint eigenvalues $0$ and $1$.
    % Source lines: paper_v2.tex 759--762.
\end{definition}

\begin{lemma}[Orthogonal defect-block sum and commuting remainder]
    \label{thm:mpu_two_projection_angle_defect_sum}
    \lean{LinearMap.IsSymmetricProjection.angleDefectBlock_orthogonalFamily,
        LinearMap.IsSymmetricProjection.angleDefectBlock_isInternal,
        LinearMap.IsSymmetricProjection.angleDefectSum_isCompl_angleDefectComplement,
        LinearMap.IsSymmetricProjection.range_le_angleDefectSum,
        LinearMap.IsSymmetricProjection.map_angleDefectSum_le_first_projection,
        LinearMap.IsSymmetricProjection.map_angleDefectSum_le_second_projection,
        LinearMap.IsSymmetricProjection.map_angleDefectComplement_le_first_projection,
        LinearMap.IsSymmetricProjection.map_angleDefectComplement_le_second_projection,
        LinearMap.IsSymmetricProjection.projections_comp_apply_eq_zero_of_mem_angleDefectComplement,
        LinearMap.IsSymmetricProjection.firstProjectionOnAngleDefectComplement,
        LinearMap.IsSymmetricProjection.secondProjectionOnAngleDefectComplement,
        LinearMap.IsSymmetricProjection.firstProjectionOnAngleDefectComplement_isSymmetricProjection,
        LinearMap.IsSymmetricProjection.secondProjectionOnAngleDefectComplement_isSymmetricProjection,
        LinearMap.IsSymmetricProjection.firstProjectionOnAngleDefectComplement_commute_secondProjection}
    \leanok
    \uses{def:mpu_two_projection_angle_defect_sum}
    The blocks $(B_i)_i$ form an orthogonal internal direct sum in $B$, and
    \begin{align}
        E & =B\mathbin{\perp\!\oplus}K.
    \end{align}
    Both $B$ and $K$ are invariant under $P$ and $Q$, with
    $\operatorname{ran}P\subseteq B$, so for every $z\in K$,
    \begin{align}
        Pz & =0, & P(Qz) & =0, & Q(Pz) & =0.
    \end{align}
    Consequently the restrictions $P|_K$ and $Q|_K$ are commuting orthogonal
    projections on $K$.
    % Source lines: paper_v2.tex 759--762.
    % Local correction: line 762 repeats P where Q is intended; see
    % docs/paper-gaps/mpu_two_projection_projector_typo.tex.
\end{lemma}

\begin{proof}\leanok
    \uses{thm:mpu_two_projection_angle_block_orthogonality}
    Pairwise orthogonality of the blocks implies independence, while their
    supremum is $B$ by definition. Hence they form an internal direct sum in
    $B$. The standard orthogonal-complement decomposition gives
    $E=B\mathbin{\perp\!\oplus}B^\perp$.

    The local action formulas show that each $B_i$ is invariant under both
    projections, hence so is $B$. Symmetry of each projection then makes
    $K=B^\perp$ invariant. The compression eigenvectors form a basis of
    $\operatorname{ran}P$, and each lies in its block $B_i$; therefore
    $\operatorname{ran}P\subseteq B$. Thus $P$ vanishes on $K$. Since $Qz\in K$
    whenever $z\in K$, both $P(Qz)$ and $Q(Pz)$ vanish. The two restricted
    projections therefore commute.
\end{proof}
\begin{definition}[Reduced projection]
    \label{def:mpu_two_projection_reduced_projection}
    \lean{LinearMap.IsSymmetricProjection.reducedProjection}
    \leanok

    Let $P$ and $Q$ be orthogonal projections on a finite-dimensional complex
    inner-product space. Define $\widetilde P$ as the orthogonal projection onto
    $\operatorname{ran}(PQ)$.
    % The source defines Ptilde blockwise; issue #6294 owns the identification
    % of that projector with this canonical construction.
    % Source lines: paper_v2.tex 764--770.
\end{definition}

\begin{lemma}[Reduced-projection identities]
    \label{thm:mpu_two_projection_reduced_projection}
    \lean{LinearMap.IsSymmetricProjection.reducedProjection_isSymmetric,
        LinearMap.IsSymmetricProjection.range_reducedProjection,
        LinearMap.IsSymmetricProjection.reducedProjection_range_triple,
        LinearMap.IsSymmetricProjection.comp_reducedProjection,
        LinearMap.IsSymmetricProjection.reducedProjection_comp,
        LinearMap.IsSymmetricProjection.comp_reducedProjection_left,
        LinearMap.IsSymmetricProjection.range_reducedProjection_comp}
    \leanok
    \uses{def:mpu_two_projection_reduced_projection}
    The reduced projection satisfies
    \begin{align}
        \operatorname{ran}(\widetilde P)   & =\operatorname{ran}(PQ),           \\
        \operatorname{ran}(PQP)            & =\operatorname{ran}(PQ),           \\
        P\widetilde P                      & =\widetilde P,                     \\
        \widetilde P Q                     & =PQ,                               \\
        Q\widetilde P                      & =QP,                               \\
        \operatorname{ran}(\widetilde P Q) & =\operatorname{ran}(\widetilde P).
    \end{align}
    The first two identities describe the canonical construction, while the next
    three establish the properties in source lines 767--769. The final range
    equality is the derived input to the right-factor theorem, not item~(iv) of
    the source.
    % Identifying the canonical construction with the source's blockwise
    % projector remains in issue #6294.
    % Source lines: paper_v2.tex 764--769.
\end{lemma}

\begin{proof}\leanok
    With $A=QP$, the standard range identity for $A^\dagger A$ gives
    \begin{align}
        \operatorname{ran}(PQP)
         & =\operatorname{ran}(A^\dagger A)
        =\operatorname{ran}(A^\dagger)
        =\operatorname{ran}(PQ).
    \end{align}
    By definition, $\operatorname{ran}(\widetilde P)=\operatorname{ran}(PQ)$,
    which also gives $P\widetilde P=\widetilde P$. For all $x$ and $z$,
    \begin{align}
        \left\langle Qx-PQx,\,PQz\right\rangle & =0.
    \end{align}
    Hence $Qx-PQx\in\operatorname{ran}(PQ)^\perp$, so orthogonal projection
    onto this range yields $\widetilde P Q=PQ$. Taking adjoints gives
    $Q\widetilde P=QP$. Finally,
    \begin{align}
        \operatorname{ran}(\widetilde P Q)
         & =\operatorname{ran}(PQ)
        =\operatorname{ran}(\widetilde P).
    \end{align}
\end{proof}

\begin{lemma}[Transversality of the reduced range]
    \label{thm:mpu_two_projection_reduced_projection_transversality}
    \lean{LinearMap.IsSymmetricProjection.disjoint_ker_range_reducedProjection,
        Matrix.disjoint_ker_range_reducedProjection}
    \leanok
    \uses{def:mpu_two_projection_reduced_projection, thm:mpu_two_projection_reduced_projection}
    For orthogonal projections $P$ and $Q$, the reduced range contains no
    nonzero vector annihilated by $Q$:
    \begin{align}
        \ker Q\cap\operatorname{ran}(\widetilde P) & =\{0\}.
        \label{eq:mpu_reduced_projection_transversality}
    \end{align}
    % Source lines: paper_v2.tex 764--781.
\end{lemma}

\begin{proof}\leanok
    If $x=PQy$ and $Qx=0$, then $Px=x$ and symmetry of the two projections gives
    \begin{align}
        \langle x,x\rangle
         & =\langle PQy,PQy\rangle
        =\langle Qy,PQy\rangle
        =\langle y,QPQy\rangle
        =0.
    \end{align}
    Hence $x=0$.
\end{proof}

\begin{lemma}[Invertible right factor for the reduced projection]
    \label{thm:mpu_two_projection_reduced_projection_right_factor}
    \lean{LinearMap.IsSymmetricProjection.exists_reducedProjection_rightFactor}
    \leanok
    \uses{thm:mpu_two_projection_reduced_projection}
    There is an invertible linear map $Y$ such that
    \begin{align}
        \widetilde P QY=\widetilde P.
    \end{align}
    % Source line: paper_v2.tex 770.
\end{lemma}

\begin{proof}\leanok
    The equality
    \begin{align}
        \operatorname{ran}(\widetilde P Q) & =\operatorname{ran}(\widetilde P)
    \end{align}
    and the equal-ranges right-factor theorem give an invertible $Y$ satisfying
    $\widetilde P QY=\widetilde P$.
\end{proof}

\begin{definition}[Matrix reduced projection]
    \label{def:mpu_matrix_reduced_projection}
    \lean{Matrix.reducedProjection}
    \leanok
    \uses{def:mpu_two_projection_reduced_projection}
    In standard Euclidean coordinates, let $\widetilde P$ be the matrix of the
    reduced projection associated with two square matrices $P$ and $Q$.
\end{definition}

\begin{lemma}[Coordinate identities for the reduced projection]
    \label{thm:mpu_matrix_reduced_projection}
    \lean{Matrix.toEuclideanLin_reducedProjection,
        Matrix.reducedProjection_isOrthogonalProjection,
        Matrix.mul_reducedProjection,
        Matrix.reducedProjection_mul,
        Matrix.reducedProjection_mul_second,
        Matrix.second_mul_reducedProjection,
        Matrix.exists_reducedProjection_rightFactor}
    \leanok
    \uses{def:mpu_matrix_reduced_projection, thm:mpu_two_projection_reduced_projection, thm:mpu_two_projection_reduced_projection_right_factor}
    If $P$ and $Q$ are orthogonal projections, then $\widetilde P$ is an
    orthogonal projection and
    \begin{align}
        P\widetilde P  & =\widetilde P, &
        \widetilde P P & =\widetilde P, &
        \widetilde P Q & =PQ,           &
        Q\widetilde P  & =QP.
        \label{eq:mpu_reduced_projection_coordinate_identities}
    \end{align}
    There is also an invertible matrix $Y$ satisfying
    $\widetilde P QY=\widetilde P$.
    % The tagged Lean declarations prove that \widetilde P is an orthogonal
    % projection for arbitrary P and Q; the hypotheses above are those of the source.
    % Source lines: paper_v2.tex 764--770.
\end{lemma}

\begin{lemma}[Transfer map under virtual sandwiching]
    \label{lem:mpu_transfer_map_virtual_sandwich}
    \lean{MPOTensor.transferMap_virtualSandwich}
    \leanok
    \uses{def:mpo_transfer_map, def:mpu_virtual_sandwich}
    Let $U=(U^{ij})$ be an MPO tensor, and let
    $\widehat U^{ij}=AU^{ij}B$. Then
    \begin{align}
        E_{\widehat U}(X) & =A E_U(BXB^\dagger)A^\dagger.
    \end{align}
\end{lemma}

\begin{proof}\leanok
    Expanding the transfer map and collecting the factors $A$, $B$, $B^\dagger$, and
    $A^\dagger$ outside the finite sum gives the displayed identity.
\end{proof}

\begin{lemma}[Normalized transfer under virtual sandwiching]
    \label{thm:mpu_normalized_transfer_sandwich}
    \lean{MPOTensor.transferMap_normalizedFlattening_virtualSandwich,
        MPOTensor.transferMap_normalizedFlattening_reducedProjection}
    \leanok
    \uses{def:mpu_normalized_flattening, def:mpu_virtual_sandwich,
        def:mpu_matrix_reduced_projection}
    Let $W=(W^{ij})$ be an MPO tensor, let $\widehat E_W$ denote the transfer
    map of its normalized flattening, and set $\widehat W^{ij}=AW^{ij}B$.
    For every bond matrix $X$,
    \begin{align}
        \widehat E_{\widehat W}(X)
         & =A\widehat E_W(BXB^\dagger)A^\dagger.
    \end{align}
    In particular, suppose that $\widehat E_W(X)=\tr(LX)R$ for every $X$.
    For any bond matrices $P,Q$, let $T$ be their reduced projection and set
    $\widetilde W^{ij}=TW^{ij}T$. Then
    \begin{align}
        \widehat E_{\widetilde W}(X)
         & =\tr\!\bigl((TLT)X\bigr)\,TRT.
    \end{align}
    In the proof of \cite[Proposition~IV.5]{Cirac2017MPU}, this formula is the normalized compression
    identity.
    % Source: paper_v2.tex, eq:transfer-op, lines 336--340;
    % prop:continuity-index, lines 747--752 and 778--781.
\end{lemma}

\begin{proof}\leanok
    \uses{def:mpu_normalized_flattening, def:mpu_virtual_sandwich,
        thm:mpu_matrix_reduced_projection}
    Scalar normalization commutes with left and right multiplication.
    Expanding the normalized transfer map therefore gives the first identity
    by collecting the constant factors outside the finite sum. For the
    reduced projection, $T^\dagger=T$, and cyclicity of the trace gives
    $\tr(LTXT)=\tr(TLTX)$.
\end{proof}

\begin{lemma}[Positive definite support compression]
    \label{lem:mpu_support_compression_posdef}
    \lean{Matrix.PosSemidef.compression_posDef_of_support_action_ne_zero}
    \leanok

    Let $\rho\succeq0$. If
    \begin{align}
        \operatorname{supp}(\rho)(Vx) & \neq0
        \qquad\text{for every }x\neq0,
    \end{align}
    then
    \begin{align}
        V^\dagger\rho V & \succ0.
    \end{align}
\end{lemma}

\begin{proof}\leanok
    If the quadratic form of $V^\dagger\rho V$ vanished at a nonzero vector $x$,
    positive semidefiniteness would give $\rho Vx=0$. The support projection would
    then annihilate $Vx$, contrary to the hypothesis.
\end{proof}

\begin{lemma}[Reduced representative under a given positive fixed pair]
    \label{thm:mpu_reduced_representative}
    \lean{MPOTensor.mpo_virtualSandwich_reducedProjection_eq,
        MPOTensor.transferMap_virtualSandwich_reducedProjection,
        MPOTensor.trace_compressed_fixed_pair_reducedProjection,
        MPOTensor.compressed_fixed_pair_posDef_reducedProjection}
    \leanok
    \uses{def:mpo_family, def:mpo_transfer_map, def:mpu_virtual_sandwich, def:mpu_matrix_reduced_projection, lem:mpu_transfer_map_virtual_sandwich, lem:mpu_support_compression_posdef, thm:mpu_matrix_reduced_projection, thm:mpu_two_projection_reduced_projection_transversality}
    Let $W=(W^{ij})$ be an MPO tensor. Let $L,R\succeq0$ satisfy
    \begin{align}
        E_W(X)                & =\operatorname{tr}(LX)R, &
        \operatorname{tr}(LR) & =1,
    \end{align}
    and let $P$ and $Q$ be their support projections. Assume, letter by letter,
    \begin{align}
        W^{ij}P & =W^{ij}, & QW^{ij} & =W^{ij}.
    \end{align}
    Set $T=\widetilde P$ and $\widetilde W^{ij}=TW^{ij}T$. Then, for every
    $N\geq0$,
    \begin{align}
        \operatorname{MPO}_N(\widetilde W)        & =\operatorname{MPO}_N(W),                    \\
        E_{\widetilde W}(X)
                                                  & =\operatorname{tr}\!\bigl((TLT)X\bigr)\,TRT, \\
        \operatorname{tr}\!\bigl((TLT)(TRT)\bigr) & =1.
    \end{align}
    If $V^\dagger V=I$ and $VV^\dagger=T$, then
    \begin{align}
        V^\dagger LV & \succ0, & V^\dagger RV & \succ0.
    \end{align}
    % Scope restriction: the positivity, support, absorption, transfer, and
    % normalization assumptions are supplied explicitly. The source derives them
    % by blocking D^4 sites and using the isometry in the proof of Proposition IV.5.
    % See docs/paper-gaps/mpu_reduced_representative_supplied_fixed_pair.tex.
    % Source lines: paper_v2.tex 747--783.
\end{lemma}

\begin{proof}\leanok
    For every nonempty word $w$, idempotence of $T$ and
    (\ref{eq:mpu_reduced_projection_coordinate_identities}) give
    \begin{align}
        \widetilde W_w & =TW_wT.
    \end{align}
    Cyclicity of the trace, together with $QW_w=W_w=W_wP$, therefore gives
    $\operatorname{tr}(TW_wT)=\operatorname{tr}(W_w)$. Length zero is immediate.
    Lemma~\ref{lem:mpu_transfer_map_virtual_sandwich}, with $A=B=T$ and
    $T^\dagger=T$, gives
    \begin{align}
        E_{\widetilde W}(X)
         & =\operatorname{tr}(LTXT)\,TRT
        =\operatorname{tr}((TLT)X)\,TRT.
    \end{align}
    Since $P$ and $Q$ are the support projections of $L$ and $R$, respectively,
    \begin{align}
        PL & =LP=L, & QR & =RQ=R.
    \end{align}
    Combining these absorptions with
    (\ref{eq:mpu_reduced_projection_coordinate_identities}), $T^2=T$, and
    cyclicity of the trace gives
    \begin{align}
        \operatorname{tr}\!\bigl((TLT)(TRT)\bigr)
         & =\operatorname{tr}(LTRT)
        =\operatorname{tr}(LRT)
        =\operatorname{tr}(RTL)
        =\operatorname{tr}(RL)
        =\operatorname{tr}(LR)
        =1.
    \end{align}
    Let $x\neq0$ and set $u=Vx$. Since $V^\dagger V=I$, one has $u\neq0$;
    since $VV^\dagger=T$, one has $Tu=u$. Equation
    (\ref{eq:mpu_reduced_projection_coordinate_identities}) gives $Pu=u$, and
    (\ref{eq:mpu_reduced_projection_transversality}) gives $Qu\neq0$. Thus
    \begin{align}
        \operatorname{supp}(L)u & =Pu=u\neq0, &
        \operatorname{supp}(R)u & =Qu\neq0.
    \end{align}
    Applying Lemma~\ref{lem:mpu_support_compression_posdef} to $L$ and $R$ yields
    \begin{align}
        x^\dagger(V^\dagger LV)x & =u^\dagger Lu>0, &
        x^\dagger(V^\dagger RV)x & =u^\dagger Ru>0.
    \end{align}
\end{proof}

\begin{lemma}[Simplicity under a rectangular bond inclusion]
    \label{thm:mpu_simple_rectangular_inclusion}
    \lean{MPOTensor.IsMPUSimple.rectangularSandwich,
        MPOTensor.isMPUSimple_of_rectangularSandwich,
        MPOTensor.isMPUSimple_supportCoordinates}
    \leanok
    \uses{def:mpu_simple}
    Let $W$ have bond dimension $D$, and let
    $A\in\mathbb C^{E\times D}$ and $B\in\mathbb C^{D\times E}$ satisfy
    $BA=I_D$. The tensor with letters $AW^{ij}B$ is simple if and only if
    $W$ is simple. In particular, if $V^\dagger V=I_k$ and $VV^\dagger=T$,
    simplicity of $(TW^{ij}T)$ implies simplicity of $(V^\dagger W^{ij}V)$.
\end{lemma}

\begin{proof}\leanok
    \uses{def:mpu_simple}
    With $\mathcal A=\overline A\otimes A$ and
    $\mathcal B=\overline B\otimes B$ in the doubled bond coordinates,
    the double-layer letters $F^{ij}$ of $W$ transform to
    $\mathcal A F^{ij}\mathcal B$, and $\mathcal B\mathcal A=I$.
    If $(a|,|b)$ are simplicity boundary vectors for $W$, then
    $(a|\mathcal B,\mathcal A|b)$ are boundary vectors for the included tensor.
    Conversely, boundary vectors $(a|,|b)$ for the included tensor give
    $(a|\mathcal A,\mathcal B|b)$ for $W$. Substitution in
    (\ref{eq:mpu_simple1}) and (\ref{eq:mpu_simple2}) proves both claims.
    For the assertion about support coordinates, use
    $V(V^\dagger W^{ij}V)V^\dagger=TW^{ij}T$.
\end{proof}

\begin{lemma}[Simplicity under reduced support compression]
    \label{thm:mpu_simple_reduced_support}
    \lean{MPOTensor.IsMPUSimple.reducedProjection}
    \leanok
    \uses{def:mpu_simple, def:mpu_matrix_reduced_projection}
    Let $W$ be simple with $d>0$. Let $P$ be an orthogonal projection and
    $Q$ a square matrix on its bond space, with $W^{ij}P=W^{ij}=QW^{ij}$.
    Suppose the normalized diagonal $E$ of its double layer has a rank-one
    factorization $E=|\rho)(\Phi|$. If $T$ is the reduced projection
    associated with $P,Q$, then $(TW^{ij}T)$ is simple.
\end{lemma}

\begin{proof}\leanok
    \uses{thm:mpu_matrix_reduced_projection, thm:mpu_reduced_representative,
        thm:mpu_supplied_projector_simple2, thm:mpu_all_later_simple_blockings,
        thm:mpu_simple_is_mpu}
    With $F^{ij}$ denoting the double-layer letters and
    $\mathcal P=\overline P\otimes P$,
    $\mathcal Q=\overline Q\otimes Q$, and $H=\overline T\otimes T$,
    the letterwise support identities also hold for their normalized diagonal
    $E$. Since $TQ=PQ$, they give
    \begin{align}
        XHY & =XH\mathcal QY=X\mathcal P\mathcal QY=XY
        \label{eq:mpu_simple_support_insertion}
    \end{align}
    whenever $X,Y$ are double-layer letters or $E$.
    The projector identity is $F^{ij}F^{k\ell}=F^{ij}EF^{k\ell}$.
    Using (\ref{eq:mpu_simple_support_insertion}) and $H^2=H$ therefore yields
    \begin{align}
        (HF^{ij}H)(HF^{k\ell}H)
         & =(HF^{ij}H)(HEH)(HF^{k\ell}H).
    \end{align}
    These are the double-layer letters of the compressed tensor, and
    $HEH=|H\rho)(\Phi H|$ gives its two simplicity boundary vectors.
    Support compression preserves the periodic operators, so the compressed
    tensor still generates an MPU. Its two-letter insertion identity then
    implies the first simplicity contraction with the same boundary vectors.
\end{proof}

\begin{lemma}[Invertible factor for a reduced projection]
    \label{thm:reduced_projection_invertible_factor}
    \lean{Matrix.exists_units_supportInvExtension_reducedProjection_rightFactor}
    \leanok
    \uses{def:mpu_matrix_reduced_projection, def:mpu_ambient_support_inverse_extension}
    Let $L,R\succeq0$, let $P$ and $Q$ be their support projections, and let
    $\widetilde P$ be the reduced projection of $P$ and $Q$. There are
    invertible support-inverse extensions of $L$ and $R$, together with an
    invertible matrix $Y$, such that $PQY=\widetilde P$.
\end{lemma}

\begin{proof}\leanok
    Extend the inverses on the supports by the identity on their orthogonal
    complements, and choose an invertible right factor for the reduced
    projection.
\end{proof}

\begin{lemma}[Source-rank comparison of the hatted and reduced tensors]
    \label{thm:mpu_reduced_to_hat_source_ranks}
    \lean{MPOTensor.virtualSandwich_hat_eq_reducedProjection,
        MPOTensor.rightRank_hat_eq_reducedProjection,
        MPOTensor.leftRank_hat_eq_reducedProjection}
    \leanok
    \uses{def:mpo_family, def:mpu_virtual_sandwich, def:mpu_matrix_reduced_projection, def:mpu_right_rank, def:mpu_left_rank, def:mpu_ambient_support_inverse_extension, thm:reduced_projection_invertible_factor}
    Let $W=(W^{ij})$ be an MPO tensor, let $L$, $R$, $X$, $Z$, $Y$ be square
    matrices of the same size, let $P$ be an orthogonal projection, let $Q$
    be a square matrix of the same size, and set $\widetilde P$ equal to
    their reduced projection.
    Assume, letter by letter, $W^{ij}P=W^{ij}=QW^{ij}$. Set
    $\widehat W^{ij}=LW^{ij}R$ and
    $\widetilde W^{ij}=\widetilde P W^{ij}\widetilde P$. If $X,Z,Y$ satisfy
    $XL=P$, $RZ=Q$, and $PQY=\widetilde P$, then
    \begin{align}
        X\widehat W ZY             & =PWQY
        \label{eq:mpu_reduced_hat_chain_pwqy}                    \\
        PWQY                       & =PQWPQY
        \label{eq:mpu_reduced_hat_chain_pqwpqy}                  \\
        PQWPQY                     & =\widetilde P W\widetilde P
        \label{eq:mpu_reduced_hat_chain_reduced}                 \\
        \widetilde P W\widetilde P & =\widetilde W.
        \label{eq:mpu_reduced_hat_chain_tilde}
    \end{align}
    If in addition $X$, $Z$, and $Y$ are invertible, then
    \begin{align}
        r[\widehat W]    & =r[\widetilde W],\label{eq:mpu_reduced_hat_right_rank} \\
        \ell[\widehat W] & =\ell[\widetilde W].
        \label{eq:mpu_reduced_hat_left_rank}
    \end{align} If $L,R\succeq0$ and
    $P=\operatorname{supp}(L)$, $Q=\operatorname{supp}(R)$, such invertible
    matrices may be chosen as $X=L^++(\Id-P)$, $Z=R^++(\Id-Q)$, together
    with an invertible right factor $Y$ of the reduced projection.
    % Scope restriction: the fixed pair, the equations XL=P and RZ=Q, and the
    % letterwise absorptions are supplied hypotheses, not consequences of a bare
    % MPU assumption. See
    % docs/paper-gaps/mpu_reduced_representative_supplied_fixed_pair.tex.
    % The identity chain is formalized without the invertibility hypotheses
    % (the paper takes X, Y, Z invertible at line 792); invertibility of X, Z, Y
    % is used only for the rank conclusion.
    % Source lines: paper_v2.tex 786--804.
\end{lemma}

\begin{proof}\leanok
    \uses{def:mpu_virtual_sandwich, thm:mpu_matrix_reduced_projection, thm:mpu_virtual_sandwich_source_ranks, thm:mpu_ambient_support_inverse_extension}
    Expanding each virtual sandwich and using $XL=P$ and $RZ=Q$ gives
    (\ref{eq:mpu_reduced_hat_chain_pwqy}). The absorptions
    $QW^{ij}=W^{ij}=W^{ij}P$ give, letter by letter,
    \begin{align}
        PW^{ij}QY
         & =P(QW^{ij})QY=PQ(W^{ij}P)QY=PQW^{ij}PQY,
        \notag
    \end{align}
    which is (\ref{eq:mpu_reduced_hat_chain_pqwpqy}). Since
    $\widetilde P Q=PQ$ and $PQY=\widetilde P$,
    \begin{align}
        PQW^{ij}PQY
         & =(\widetilde P Q)W^{ij}\widetilde P
        =\widetilde P W^{ij}\widetilde P.
    \end{align}
    This proves (\ref{eq:mpu_reduced_hat_chain_reduced}) and
    (\ref{eq:mpu_reduced_hat_chain_tilde}). Thus $\widetilde W$ is obtained
    from $\widehat W$ by left multiplication by $X$ and right multiplication
    by $ZY$. Invertibility of these factors and the source ranks of a
    virtual sandwich (Lemma~\ref{thm:mpu_virtual_sandwich_source_ranks})
    give
    \begin{align}
        r[\widehat W]    & =r[X\widehat W(ZY)]=r[\widetilde W],
        \notag                                                        \\
        \ell[\widehat W] & =\ell[X\widehat W(ZY)]=\ell[\widetilde W],
        \notag
    \end{align}
    by invertibility of $X$ and $ZY$ together with the relation $X\widehat W(ZY)=\widetilde W$
    from~(\ref{eq:mpu_reduced_hat_chain_pwqy})--(\ref{eq:mpu_reduced_hat_chain_tilde}), yielding~(\ref{eq:mpu_reduced_hat_right_rank}) and
    (\ref{eq:mpu_reduced_hat_left_rank}). For positive semidefinite $L$ and
    $R$, their ambient support-inverse extensions give $X$ and $Z$; the
    invertible right-factor identity for $\widetilde P Q$ gives $Y$.
\end{proof}

\begin{lemma}[Lower semicontinuity of finite matrix rank]
    \label{thm:mpu_rank_lower_semicontinuous}
    \lean{Matrix.lowerSemicontinuous_rank,
        Continuous.lowerSemicontinuous_matrix_rank}
    \leanok
    For finite index sets $M,N$, matrix rank defines a lower-semicontinuous map
    \begin{align}
        \operatorname{rank}\colon
        \operatorname{Mat}_{M\times N}(\mathbb C) & \longrightarrow\mathbb N.
    \end{align}
    Consequently, if $A\colon X\to\operatorname{Mat}_{M\times N}(\mathbb C)$
    is continuous, then $x\mapsto\operatorname{rank}A(x)$ is lower
    semicontinuous.
    % Source lines: paper_v2.tex 813--814.
\end{lemma}

\begin{proof}\leanok
    For every $k\in\mathbb N$, the sublevel set is the closed determinantal set
    \begin{align}
        \{A:\operatorname{rank}A\leq k\}.
    \end{align}
    These closed sublevel sets characterize lower semicontinuity. Composition
    with a continuous map preserves lower semicontinuity.
\end{proof}

\begin{lemma}[Fixed positive product gives local constancy]
    \label{thm:mpu_fixed_product_locally_constant}
    \lean{isLocallyConstant_pair_of_lowerSemicontinuous_nat_mul_eq}
    \leanok
    Let $r,\ell\colon X\to\mathbb N$ be lower semicontinuous. Suppose that
    $c>0$ and
    \begin{align}
        r(x)\ell(x) & =c
    \end{align}
    for every $x\in X$. Then both $r$ and $\ell$ are locally constant.
\end{lemma}

\begin{proof}\leanok
    For any $x\in X$, since $c>0$, both $r(x)$ and $\ell(x)$ are positive. Lower
    semicontinuity gives a neighborhood $U$ of $x$ on which
    \begin{align}
        r(y) & \geq r(x), & \ell(y) & \geq\ell(x).
    \end{align}
    For $y\in U$, monotonicity of multiplication and the fixed-product identity
    give
    \begin{align}
        r(x)\ell(x)
         & \leq r(y)\ell(x)
        \leq r(y)\ell(y)
        =r(x)\ell(x).
    \end{align}
    Positivity permits cancellation, first of $\ell(x)$ and then of $r(x)$, so
    $r(y)=r(x)$ and $\ell(y)=\ell(x)$.
\end{proof}

\begin{lemma}[Fixed positive product on a preconnected set]
    \label{thm:mpu_fixed_product_preconnected}
    \lean{eq_of_lowerSemicontinuous_nat_mul_eq_of_isPreconnected,
        eq_of_lowerSemicontinuous_nat_mul_eq,
        unitInterval_eq_of_lowerSemicontinuous_nat_mul_eq}
    \leanok
    Under the hypotheses of
    Lemma~\ref{thm:mpu_fixed_product_locally_constant}, if $S\subseteq X$ is
    preconnected, then for all $x,y\in S$,
    \begin{align}
        r(x) & =r(y), & \ell(x) & =\ell(y).
    \end{align}
    In particular, both functions are constant when $X$ is preconnected and
    when $X=[0,1]$.
\end{lemma}

\begin{proof}\leanok
    \uses{thm:mpu_fixed_product_locally_constant}
    A locally constant function has open fibers. A preconnected set cannot meet
    two disjoint nonempty open fibers, so each of $r$ and $\ell$ has one value
    on $S$. The unit interval is connected, hence preconnected.
\end{proof}

\begin{lemma}[Blocking, simplicity and ranks under a bond similarity]
    \label{thm:mpu_virtual_sandwich_blocking}
    \lean{MPOTensor.blockTensor_virtualSandwich, MPOTensor.isMPUSimple_virtualSandwich_iff,
        MPOTensor.mpo_virtualSandwich, MPOTensor.IsMPU.virtualSandwich,
        MPOTensor.rightRank_blockTensor_virtualSandwich,
        MPOTensor.leftRank_blockTensor_virtualSandwich}
    \leanok
    \uses{thm:mpu_virtual_sandwich_source_ranks, def:mpu_simple, def:mpdo_physical_blocking}
    Let $A,B$ be square matrices on the bond space with $BA=I$. The tensor with
    letters $AU^{ij}B$ has the blockings $A\,U_L\,B$, the same periodic
    operators as $\mathcal U$, and is simple if and only if $\mathcal U$ is;
    the ranks of its blockings equal those of $\mathcal U$.
    % Source: paper_v2.tex, Proposition IV.5, lines 773--804.
\end{lemma}

\begin{proof}\leanok
    The letters of a blocking are products, and $BA=I$ telescopes; the trace is
    cyclic; the double layer is conjugated by $\bar A\otimes A$, which carries
    the boundary vectors along; the ranks by the source-ranks theorem.
\end{proof}

\begin{lemma}[Source factors under a unitary virtual conjugation]
    \label{thm:mpu_source_factors_virtual_unitary_gauge}
    \lean{MPOTensor.source_factors_virtual_unitary_sandwich,
        MPOTensor.source_rank_virtual_unitary_sandwich,
        MPOTensor.transported_source_factor_premises_at_selected_ranks,
        MPOTensor.IsMPUCanonicalFormII.exists_selected_source_factor_unitary_gauges}
    \leanok
    \uses{def:mpu_simple, def:mpu_canonical_form_ii,
        def:mpu_virtual_sandwich, def:mpu_weighted_source_factors,
        thm:mpu_virtual_sandwich_source_ranks}
    Let $\mathcal U$ and $\mathcal V$ be simple tensors in canonical form II
    with $V^{ij}=zU^{ij}z^\dagger$ for a unitary $z$ on the virtual space.
    Their recorded positive matrices need not be related. Write
    $X_a(\mathcal U),Y_a(\mathcal U)$ and
    $X_a(\mathcal V),Y_a(\mathcal V)$ for their selected source factors,
    $a=1,2$. After identifying the equal source ranks, there are unitaries
    $W_1$ and $W_2$ on the respective source-rank spaces such that
    \begin{align}
        X_1(\mathcal V)
         & =(I_d\otimes(z^\dagger)^T)X_1(\mathcal U)W_1,
         & Y_1(\mathcal V)
         & =W_1^\dagger Y_1(\mathcal U)(z^T\otimes I_d),       \\
        X_2(\mathcal V)
         & =(z\otimes I_d)X_2(\mathcal U)W_2,
         & Y_2(\mathcal V)
         & =W_2^\dagger Y_2(\mathcal U)(I_d\otimes z^\dagger).
    \end{align}
    This is the forward factor
    comparison associated with the diagrams in
    \cite[lines~624--648]{Cirac2017MPU}; it does not assert
    a converse from the local relations to equality at every chain length.
    In \cite[lines~1052--1066]{FrancoRubio2025MPUGauging}, the scalar-free step is the isometric
    specialization.
    % Source: paper_v2.tex 624--648 and II_RelationStandard1/2.png;
    % raw cut/rank transport at lines 786--812.
\end{lemma}

\begin{proof}\leanok
    \uses{thm:mpu_virtual_sandwich_source_cuts,
        thm:mpu_virtual_sandwich_source_ranks,
        thm:mpu_source_factorizations,
        thm:mpu_source_factor_normalizations,
        thm:mpu_source_factor_right_inverses, ThmFund1,
        cor:mpug_normalized_source_decomposition_gauges}
    The two cut identities of
    Lemma~\ref{thm:mpu_virtual_sandwich_source_cuts}, applied to the
    selected factorizations of Lemma~\ref{thm:mpu_source_factorizations},
    give transformed factorizations with the four exterior matrices
    displayed in the statement. The right exterior matrices are unitary, so
    the transformed right factors retain the right inverses of
    Lemma~\ref{thm:mpu_source_factor_right_inverses}. The left exterior
    matrix of the second cut is unitary, so its transformed left factor
    remains an isometry. The contraction defining $u$ is unchanged: its
    interior virtual product is $z^\dagger z=I_D$. Rank preservation,
    Lemma~\ref{thm:mpu_virtual_sandwich_source_ranks}, identifies the
    transformed factors with the selected source-rank spaces. Simplicity and
    the source-gate equivalence of Theorem~\ref{ThmFund1} give unitarity of
    both $u$ contractions. The selected factors have left inverses from their
    weighted first-cut and ordinary second-cut normalizations in
    Lemma~\ref{thm:mpu_source_factor_normalizations}. The normalized
    source-decomposition comparison of
    Corollary~\ref{cor:mpug_normalized_source_decomposition_gauges} now
    gives the two unitary changes $W_1,W_2$ and the four displayed
    equalities.
\end{proof}

\begin{remark}
    In the setting of
    Lemma~\ref{thm:mpu_source_factors_virtual_unitary_gauge}, write
    $u_{\mathcal U},u_{\mathcal V}$ for the source gates of
    Definition~\ref{def:mpu_source_uv} built from the selected factors.
    Substituting the two $Y_a(\mathcal V)$ identities into that contraction,
    the interior virtual product is $\sum_\beta(z^\dagger)_{\beta'\beta}
        z_{\beta\beta''}=(z^\dagger z)_{\beta'\beta''}=\delta_{\beta'\beta''}$,
    so
    \begin{align}
        u_{\mathcal V} & =(W_2^\dagger\otimes W_1^\dagger)u_{\mathcal U}.
    \end{align}
    Thus the unitaries $x=W_2^\dagger$ and $y=W_1^\dagger$ realize the local
    relation $u_{\mathcal V}=(x\otimes y)u_{\mathcal U}$ of
    Theorem~\ref{FundamentalMPU}. This identification of the gauges is a
    consequence of the displayed factor identities and is not asserted by
    Lemma~\ref{thm:mpu_source_factors_virtual_unitary_gauge}. It is
    stated, together with the corresponding relation for $v$, in
    Lemma~\ref{thm:mpu_source_gates_virtual_unitary_gauge}.
\end{remark}

\begin{lemma}[Source gates under a unitary virtual conjugation]
    \label{thm:mpu_source_gates_virtual_unitary_gauge}
    \lean{MPOTensor.IsMPUCanonicalFormII.exists_selected_source_gate_unitary_gauges}
    \leanok
    \uses{def:mpu_simple, def:mpu_canonical_form_ii,
        def:mpu_virtual_sandwich, def:mpu_source_uv,
        thm:mpu_virtual_sandwich_source_ranks}
    Let $\mathcal U$ and $\mathcal V$ be simple tensors in canonical form II
    with $V^{ij}=zU^{ij}z^\dagger$ for a unitary $z$. Let
    $e_r:r[\mathcal U]\simeq r[\mathcal V]$ and
    $e_\ell:\ell[\mathcal U]\simeq\ell[\mathcal V]$ be the source-rank
    bijections. Reindex the source gates of $\mathcal U$ by
    \begin{align}
        \widetilde u_{(\ell,r),(p,q)}
         & =u_{\mathcal U,(e_\ell^{-1}\ell,e_r^{-1}r),(p,q)},
         & \widetilde v_{(p,q),(r,\ell)}
         & =v_{\mathcal U,(p,q),(e_r^{-1}r,e_\ell^{-1}\ell)}.
    \end{align}
    There are unitaries $x$ on $\ell[\mathcal V]$ and $y$ on
    $r[\mathcal V]$ such that
    \begin{align}
        u_{\mathcal V} & =(x\otimes y)\widetilde u,
        \label{eq:mpu_uv_rel_u}                                     \\
        v_{\mathcal V} & =\widetilde v(y^\dagger\otimes x^\dagger).
        \label{eq:mpu_uv_rel_v}
    \end{align}
    Relation~(\ref{eq:mpu_uv_rel_u}) is the forward $u$ relation of
    \cite[Equation~(SFuu)]{Cirac2017MPU}. The source's
    \cite[Equation~(SFvv)]{Cirac2017MPU} displays two relations between the
    open half-factors involving $z$; contraction of those half-factors yields~(\ref{eq:mpu_uv_rel_v}). No converse for the original periodic
    operators at odd length is asserted.
    % Source: paper_v2.tex 624--648; II_RelationStandard1/2.png.
\end{lemma}

\begin{proof}\leanok
    \uses{thm:mpu_source_factors_virtual_unitary_gauge}
    By the four selected-factor identities of
    Lemma~\ref{thm:mpu_source_factors_virtual_unitary_gauge}, the
    source-rank changes are $x=W_2^\dagger$ and $y=W_1^\dagger$. The
    two source-gate contractions have entries
    \begin{align}
        u_{(\ell,r),(p,q)}
         & =\sum_\beta (Y_2)_{\ell,(p,\beta)}(Y_1)_{r,(\beta,q)},    \\
        v_{(p,q),(r,\ell)}
         & =\sum_\alpha (X_1)_{(p,\alpha),r}(X_2)_{(\alpha,q),\ell}.
    \end{align}
    In the first contraction, the two virtual matrices cancel as
    $z^\dagger z=I_D$, leaving $(x\otimes y)\widetilde u$. In the second,
    the first-cut transpose and second-cut left action likewise contract
    to $z^\dagger z=I_D$, leaving
    $\widetilde v(y^\dagger\otimes x^\dagger)$.
\end{proof}

\begin{definition}[Ambient support-inverse extension]
    \label{def:mpu_ambient_support_inverse_extension}
    \lean{Matrix.PosSemidef.supportInvExtension,
        Matrix.PosSemidef.supportInvExtensionInverse}
    \leanok

    Let $A\succeq0$ have support projection $P$, and let $A^+$ denote its
    inverse on the support. Define
    \begin{align}
        X_A & =A^++(\Id-P).
    \end{align}
    Define also
    \begin{align}
        Z_A & =A+(\Id-P).
    \end{align}
    For Proposition~IV.5 of \cite{Cirac2017MPU}, this construction produces the required invertible
    matrices.
    % Source requirement: paper_v2.tex, Proposition IV.5, lines 786--797.
    % The displayed construction is a formal witness, not a formula from the paper.
\end{definition}

\begin{lemma}[Invertibility and support identities]
    \label{thm:mpu_ambient_support_inverse_extension}
    \lean{Matrix.PosSemidef.supportInvExtension_mul_self,
        Matrix.PosSemidef.self_mul_supportInvExtension,
        Matrix.PosSemidef.supportInvExtension_mul_inverse,
        Matrix.PosSemidef.supportInvExtension_inverse_mul,
        Matrix.PosSemidef.supportInvExtensionUnit,
        Matrix.PosSemidef.isUnit_supportInvExtension,
        Matrix.PosSemidef.isUnit_det_supportInvExtension}
    \leanok
    \uses{def:mpu_ambient_support_inverse_extension}
    The matrices $X_A$ and $Z_A$ are mutually inverse, and
    \begin{align}
        X_AA=AX_A & =P.
    \end{align}
    Consequently, $X_A$ is a unit in the matrix algebra, with inverse $Z_A$,
    and $\det(X_A)$ is a unit in $\mathbb C$.

    In the notation of Proposition~IV.5 of \cite{Cirac2017MPU}, let $L$ and
    $R$ be the positive semidefinite outer factors with support projections
    $P$ and $Q$, respectively. The required matrices are
    \begin{align}
        X & =X_L=L^++(\Id-P), \\
        Z & =X_R=R^++(\Id-Q).
    \end{align}
    Thus
    \begin{align}
        X_L L & =P, \\
        R X_R & =Q,
    \end{align}
    or equivalently $XL=P$ and $RZ=Q$. Here the paper's matrix $Z$ is $X_R$,
    not the inverse-extension matrix $Z_R$ from
    Definition~\ref{def:mpu_ambient_support_inverse_extension}.
    % This proves the existence step used in paper_v2.tex, lines 786--797.
\end{lemma}

\begin{proof}\leanok
    The support identities give
    \begin{align}
        A^+A=AA^+ & =P.
    \end{align}
    Their complementary forms are
    \begin{align}
        A^+(\Id-P)=(\Id-P)A^+ & =0,
    \end{align}
    \begin{align}
        A(\Id-P)=(\Id-P)A & =0,
    \end{align}
    and idempotence of $P$ gives
    \begin{align}
        (\Id-P)^2 & =\Id-P.
    \end{align}
    Therefore
    \begin{align}
        X_AZ_A & =P+(\Id-P)=\Id.
    \end{align}
    Likewise,
    \begin{align}
        Z_AX_A & =P+(\Id-P)=\Id,
    \end{align}
    while multiplication by $A$ gives $X_AA=AX_A=P$.
\end{proof}

\begin{lemma}[Rank-one transfer maps have positive trace factors]
    \label{lem:mpu_rank_one_transfer_trace_pair}
    \lean{Kraus.exists_posSemidef_trace_pair_of_transferMatrix_eq_vecMulVec,
        Kraus.mul_supportProj_of_map_eq_trace_smul,
        Kraus.supportProj_mul_of_map_eq_trace_smul}
    \leanok
    Let $A_1,\dots,A_n\in\MN{D}$ and let $\mathcal E(X)=\sum_aA_aXA_a^\dagger$.
    Suppose that the matrix of $\mathcal E$ is $|\rho)(\Phi|$ with
    $(\Phi|\rho)=1$. Then there are $L,R\succeq0$ with
    \begin{align}
        \mathcal E(X) & =\operatorname{tr}(LX)R
    \end{align}
    for every $X$. Whenever $\mathcal E(X)=\operatorname{tr}(LX)R$ with
    $L,R\succeq0$, the support projections $P$ of $L$ and $Q$ of $R$ satisfy
    \begin{align}
        A_aP & =A_a, & QA_a & =A_a
    \end{align}
    for every $a$.
    % Source lines: paper_v2.tex 747--756.
\end{lemma}

\begin{proof}\leanok
    Set $R=\mathcal E(I)=\sum_aA_aA_a^\dagger\succeq0$. Every output of
    $\mathcal E$ is a multiple of the matrix whose vectorization is $\rho$, and
    $R\neq0$: otherwise every $A_a$ vanishes, contrary to $(\Phi|\rho)=1$.
    Hence $\operatorname{tr}R>0$ and
    $\mathcal E(X)=\operatorname{tr}\mathcal E(X)\,(\operatorname{tr}R)^{-1}R$.
    Since $\operatorname{tr}\mathcal E(X)=\operatorname{tr}\bigl(
        (\sum_aA_a^\dagger A_a)X\bigr)$, the matrix
    $L=(\operatorname{tr}R)^{-1}\sum_aA_a^\dagger A_a\succeq0$ works.

    For the absorptions, $L(I-P)=0$ gives
    $\sum_a\bigl(A_a(I-P)\bigr)\bigl(A_a(I-P)\bigr)^\dagger
        =\mathcal E(I-P)=0$, so $A_a(I-P)=0$. Similarly $(I-Q)R=0$ gives
    $\sum_a\bigl((I-Q)A_a\bigr)\bigl((I-Q)A_a\bigr)^\dagger
        =(I-Q)\mathcal E(I)(I-Q)=0$, so $(I-Q)A_a=0$.
\end{proof}

\begin{lemma}[Simultaneous congruence of two positive definite matrices]
    \label{lem:mpu_posdef_congruence_diagonal}
    \lean{Matrix.PosDef.exists_congruence_isDiag}
    \leanok
    Let $L,R\in\MN{m}$ be positive definite. There is an invertible $S$ such
    that $S^\dagger S=L$ and $SRS^\dagger$ is positive definite and diagonal.
    % Source lines: paper_v2.tex 798--804.
\end{lemma}

\begin{proof}\leanok
    Let $G=L^{1/2}$ and let $U$ be a unitary diagonalizing the positive
    definite matrix $GRG$. Then $S=U^\dagger G$ satisfies
    $S^\dagger S=G^2=L$ and $SRS^\dagger=U^\dagger(GRG)U$.
\end{proof}

\begin{lemma}[Canonical form from a rank-one transfer map]
    \label{lem:mpu_rank_one_transfer_canonical_form}
    \lean{MPOTensor.mpo_rectangularSandwich_of_absorb,
        MPOTensor.exists_isMPUCanonicalFormII_of_transferMap_eq_trace_smul}
    \leanok
    \uses{def:mpu, def:mpu_canonical_form_ii, def:mpu_virtual_sandwich, def:mpu_right_rank, def:mpu_left_rank}
    Let $\mathcal W$ generate MPU with $d>0$, and suppose that its normalized
    transfer map is $X\mapsto\operatorname{tr}(LX)R$ with $L,R\succeq0$. Then
    there is a tensor $\mathcal B$ in canonical form II such that
    $\operatorname{MPO}_N(\mathcal B)=\operatorname{MPO}_N(\mathcal W)$ for
    every $N\geq1$ and
    \begin{align}
        r[\mathcal B]    & =r[L\mathcal WR],    &
        \ell[\mathcal B] & =\ell[L\mathcal WR].
    \end{align}
    % Source lines: paper_v2.tex 756--804.
\end{lemma}

\begin{proof}\leanok
    \uses{lem:mpu_rank_one_transfer_trace_pair, lem:mpu_posdef_congruence_diagonal, thm:mpu_reduced_representative, thm:mpu_reduced_to_hat_source_ranks, thm:reduced_projection_invertible_factor, lem:mpu_rectangular_reduction_source_ranks, thm:mpu_canonical_form_ii_stabilization_converse}
    Let $P,Q$ be the support projections of $L,R$ and let $\widetilde P$ be
    their reduced projection. By
    Lemma~\ref{lem:mpu_rank_one_transfer_trace_pair}, the letters satisfy
    $W^{ij}P=W^{ij}=QW^{ij}$. Choose an isometry $J$ with
    $JJ^\dagger=\widetilde P$. By
    Lemma~\ref{thm:mpu_reduced_representative}, $J^\dagger LJ$ and
    $J^\dagger RJ$ are positive definite; let $S$ be the congruence of
    Lemma~\ref{lem:mpu_posdef_congruence_diagonal} and put
    $\Lambda=S(J^\dagger RJ)S^\dagger$. Define
    \begin{align}
        B^{ij} & =SJ^\dagger W^{ij}JS^{-1}.
    \end{align}
    Since $(JS^{-1})(SJ^\dagger)=\widetilde P$ and the letters of
    $\widetilde W=\widetilde PW\widetilde P$ absorb $\widetilde P$ on the
    right, nonempty words of $\mathcal B$ are compressions of words of
    $\widetilde{\mathcal W}$ with the same trace, so $\mathcal B$ and
    $\mathcal W$ generate the same periodic operators. The tensors
    $\mathcal B$ and $\widetilde{\mathcal W}$ are rectangular compressions of
    each other, so they have the same source ranks, and
    Lemma~\ref{thm:mpu_reduced_to_hat_source_ranks} identifies these with the
    ranks of $L\mathcal WR$. Finally
    $(S^{-1})^\dagger J^\dagger LJS^{-1}=I$, so the normalized transfer map of
    $\mathcal B$ is $X\mapsto\operatorname{tr}(X)\Lambda$ with $\Lambda$
    positive definite and diagonal. Lemma~\ref{thm:mpu_canonical_form_ii_stabilization_converse}
    puts $\mathcal B$ in canonical form II.
\end{proof}

\begin{lemma}[Blocked transfer matrices of an MPU are rank one]
    \label{lem:mpu_blocked_transfer_rank_one}
    \lean{MPOTensor.transferMatrix_normalizedFlattening_blockTensor,
        MPOTensor.IsMPU.exists_transferMatrix_blockTensor_eq_vecMulVec}
    \leanok
    \uses{def:mpu, def:mpu_normalized_flattening}
    The normalized transfer matrix of the block of $k$ sites is the $k$-th
    power of the normalized transfer matrix. If $\mathcal W$ generates MPU
    with $d>0$ and $k\geq D^2$, it equals $|\rho)(\Phi|$ for vectors with
    $(\Phi|\rho)=1$.
    % Source lines: paper_v2.tex 397--409 and 747--752.
\end{lemma}

\begin{proof}\leanok
    \uses{thm:mpu_blocked_double_layer_diagonal, thm:mpu_normalized_transfer_stabilization}
    Write $E_{\mathcal W}$ for the normalized transfer matrix of
    $\mathcal W$. The normalized transfer matrix of the block is the
    normalized physical diagonal of its double layer, so
    Lemma~\ref{thm:mpu_blocked_double_layer_diagonal} gives
    \begin{align}
        E_{\mathcal W_k} & =E_{\mathcal W}^k,
        \label{eq:mpu_blocked_transfer_power}
    \end{align}
    which is the first claim. For $D>1$ every power of exponent at least $D^2-1$ equals the
    stabilized rank-one projector. For $D=1$ the normalized transfer matrix is
    the identity of a one-dimensional space.
\end{proof}

\begin{lemma}[Canonical source ranks from the factor-free sandwich]
    \label{lem:mpu_canonical_ranks_factor_free}
    \lean{MPOTensor.IsMPUCanonicalFormII.bondDim_le_of_mpo_eq,
        MPOTensor.IsMPUCanonicalFormII.rightRank_mul_leftRank_blockTensor_of_mpo_eq,
        MPOTensor.IsMPU.sourceRanks_blockTensor_eq_factorFreeSandwich}
    \leanok
    \uses{def:mpu, def:mpu_canonical_form_ii, def:mpu_factor_free_sandwich, def:mpu_right_rank, def:mpu_left_rank,
        thm:mpu_bounded_simple_blocking}
    Let $\mathcal W$ generate MPU with $d>0$, let $k\geq D^2$, and let
    $\mathcal U$, of bond dimension $D'$, be in canonical form II and generate
    the same MPU as $\mathcal W$. Then $D'\leq D$, the block $\mathcal U_k$ is
    simple with $r[\mathcal U_k]\ell[\mathcal U_k]=d^{2k}$, and
    \begin{align}
        r[\mathcal U_k]    & =r[\widehat{\mathcal W}_{k,\mathrm{ff}}],    &
        \ell[\mathcal U_k] & =\ell[\widehat{\mathcal W}_{k,\mathrm{ff}}],
    \end{align}
    where $\widehat{\mathcal W}_{k,\mathrm{ff}}$ is the factor-free sandwich
    of the block $\mathcal W_k$.
    % Source lines: paper_v2.tex 747--812.
\end{lemma}

\begin{proof}\leanok
    \uses{thm:mpu_reduced_cfii_representative, thm:mpu_cfii_equal_periodic_virtual_unitary, thm:mpu_bounded_simple_blocking, prop:mpu_admissible_index_well_defined, lem:mpu_blocked_transfer_rank_one, lem:mpu_rank_one_transfer_trace_pair, lem:mpu_rank_one_transfer_canonical_form, thm:mpu_cfii_index_periodic_comparison, thm:mpu_canonical_form_ii_positive_blocking, thm:mpu_factor_free_sandwich_continuous}
    Canonical-form-II tensors generating the same MPU are related by a unitary
    virtual gauge, so $D'$ equals the bond dimension of the reduced
    representative, which is at most $D$. Hence $D'^2\leq k$ and
    $\mathcal U_k$ is simple. By Lemmas~\ref{lem:mpu_blocked_transfer_rank_one}
    and~\ref{lem:mpu_rank_one_transfer_trace_pair}, the normalized transfer map
    of $\mathcal W_k$ is $X\mapsto\operatorname{tr}(LX)R$ with $L,R\succeq0$,
    and $\widehat{\mathcal W}_{k,\mathrm{ff}}=L\mathcal W_kR$.
    Lemma~\ref{lem:mpu_rank_one_transfer_canonical_form} gives a
    canonical-form-II tensor $\mathcal B$ with the ranks of $L\mathcal W_kR$
    and the periodic operators of $\mathcal W_k$, which are those of
    $\mathcal U_k$. Canonical-form-II tensors with equal periodic operators
    have equal source ranks.
\end{proof}

\begin{lemma}[Source ranks along a path with fixed rank product]
    \label{lem:mpu_path_fixed_product_source_ranks}
    \lean{MPOTensor.lowerSemicontinuous_rightRank,
        MPOTensor.lowerSemicontinuous_leftRank,
        MPOTensor.sourceRanks_const_of_fixed_product}
    \leanok
    \uses{def:mpu_right_rank, def:mpu_left_rank}
    Both source ranks are lower semicontinuous along a continuous family of
    tensors. If $[0,1]\ni x\mapsto\mathcal V(x)$ is continuous and
    $r[\mathcal V(x)]\ell[\mathcal V(x)]=c>0$ for every $x$, then both ranks
    are constant.
    % Source lines: paper_v2.tex 813--818.
\end{lemma}

\begin{proof}\leanok
    \uses{thm:mpu_rank_lower_semicontinuous, thm:mpu_fixed_product_preconnected}
    The source cuts are coordinate rearrangements of the tensor, so they vary
    continuously and their ranks are lower semicontinuous. The fixed positive
    product makes both ranks constant on the connected interval.
\end{proof}

\begin{lemma}[Common simple blocking along a continuous MPU path]
    \label{thm:mpu_common_pow_four_simple_blocking}
    \lean{MPOTensor.IsMPU.blockTensor_pow_four_isMPUSimple,
        MPOTensor.continuous_evalWord,
        MPOTensor.continuous_blockTensor,
        Continuous.blockTensor}
    \leanok
    \uses{def:mpu, def:mpu_simple, def:mpdo_physical_blocking}
    Let $d,D>0$.  Every bond-$D$ MPU tensor becomes simple after blocking
    exactly $D^4$ sites.  For each fixed $L$, the map
    \begin{align}
        \mathcal U & \longmapsto \mathcal U_L
    \end{align}
    is continuous.  Consequently, if $x\mapsto\mathcal W(x)$ is a continuous
    family of bond-$D$ MPU tensors, then
    $x\mapsto\mathcal W(x)_{D^4}$ is a continuous family of simple tensors.
    This provides the common-blocking step for
    Lemma~\ref{prop:continuity-index}.
    % Source lines: paper_v2.tex 747--752.
\end{lemma}

\begin{proof}\leanok
    \uses{thm:mpu_bounded_simple_blocking, thm:mpu_all_later_simple_blockings}
    Every tensor is simple after blocking any $k\geq D^2$ sites, in
    particular after blocking $D^4\geq D^2$ sites.
    For fixed blocked ket and bra words, every matrix entry of the blocked
    tensor is a finite sum of products of entries of the original tensor.
    These entries are continuous, and the finite product and sum operations
    preserve continuity.  Applying this fixed blocking map to the original
    path proves the path statement.
\end{proof}

\begin{lemma}[Constancy of the source ranks]
    \label{prop:continuity-index}
    \lean{MPOTensor.IsMPU.sourceRanks_canonicalFormII_blockTensor_pow_four_const,
        MPOTensor.IsMPU.sourceRanks_canonicalFormII_blockTensor_const,
        MPOTensor.IsMPU.sourceRanks_canonicalForm_blockTensor_const}
    \leanok
    \discussion{6022}
    \uses{def:mpu, def:mpu_canonical_form, def:mpu_canonical_form_ii, def:mpu_right_rank, def:mpu_left_rank}
    Let $[0,1]\ni x\mapsto\mathcal W(x)$ be a continuous path of tensors of
    physical dimension $d>0$ and bond dimension $D$ generating MPUs, without
    requiring the path to be in canonical form. For each $x$, let
    $\mathcal U(x)$ be a tensor in canonical form generating the same MPU as
    $\mathcal W(x)$, and let $r(x)$ and $\ell(x)$ be the two source-cut ranks
    of its block of $D^4$ sites. Then $r(x)$ and $\ell(x)$ are
    constant on $[0,1]$. The same holds at every blocking length $k\geq D^2$.
    % Source lines: paper_v2.tex 733--818.
\end{lemma}

\begin{proof}\leanok
    \uses{lem:mpu_canonical_ranks_factor_free, lem:mpu_path_fixed_product_source_ranks, thm:mpu_factor_free_sandwich_continuous, thm:mpu_canonical_form_similar_cfii, thm:mpu_reduced_cfii_representative}
    A tensor in canonical form is similar to one in canonical form II, and
    blocking commutes with this similarity without changing the source ranks,
    so it suffices to take $\mathcal U(x)$ in canonical form II. Fix
    $k\geq D^2$. By Lemma~\ref{lem:mpu_canonical_ranks_factor_free}, the source
    ranks of $\mathcal U(x)_k$ are those of the factor-free sandwich of
    $\mathcal W(x)_k$, which is a continuous function of $x$. Applying the same
    lemma to a reduced representative shows that the product of these ranks is
    $d^{2k}$ at every point. Lemma~\ref{lem:mpu_path_fixed_product_source_ranks}
    makes both ranks constant. The blocking length $D^4$ satisfies
    $D^4\geq D^2$.
\end{proof}

\begin{lemma}[Index along an arbitrary continuous path]
    \label{thm:mpu_path_index_constant}
    \lean{MPOTensor.IsMPU.index_const_of_continuous}
    \leanok
    \uses{def:mpu_reduced_representative_index, prop:continuity-index}
    Under the hypotheses of Lemma~\ref{prop:continuity-index}, the index of
    $\mathcal W(x)$ is independent of $x\in[0,1]$.
    % Source lines: paper_v2.tex 681--688 and 733--818.
\end{lemma}

\begin{proof}\leanok
    \uses{prop:continuity-index, thm:mpu_reduced_representative_index_invariance, prop:mpu_admissible_index_well_defined, thm:mpu_reduced_cfii_representative}
    The index of $\mathcal W(x)$ is the index of a reduced canonical-form-II
    representative $\mathcal U(x)$, computed from the source ranks of its
    simple block of $D^2$ sites as
    $\tfrac12(\log_2r-\log_2\ell)$. These ranks are constant by
    Lemma~\ref{prop:continuity-index}.
\end{proof}

\begin{lemma}[Admissible constancy of the source ranks]
    \label{prop:mpu_admissible_continuity_index}
    \notready
    \discussion{6022}
    % **Scope restriction (pathwise reduced source data):** every associated
    % simple tensor must carry a supplied fixed pair, to which the
    % same-tensor converse applies before Theorem~\ref{ThmFund1}.
    % Documented in docs/paper-gaps/mpu_canonical_form_full_support.tex.
    \uses{ThmFund1, thm:mpu_canonical_form_ii_stabilization_converse, thm:mpu_bounded_simple_blocking, def:mpu_reduced_full_support_source_datum, def:mpu_admissible_equivalence_datum}
    Let $[0,1]\ni x\mapsto\mathcal W(x)$ be a continuous path of tensors
    generating MPUs, without requiring the path to be in canonical form.
    Suppose that the path carries an admissible equivalence datum with blocking
    length $m$. Then the ranks $r(x)$ and $\ell(x)$ of the blocked path
    tensors $\mathcal W(x)_m$ are constant on $[0,1]$.
    % Source lines: paper_v2.tex 733--818.
\end{lemma}

\begin{lemma}[Source ranks along a canonical-form-II path]
    \label{thm:mpu_cfii_path_source_ranks}
    \lean{MPOTensor.IsMPUCanonicalFormII.sourceRanks_blockTensor_pow_four_const,
        MPOTensor.sourceRanks_blockTensor_const_of_fixed_product}
    \leanok
    % **Scope restriction (canonical-form-II path):** Proposition IV.5 permits
    % arbitrary MPU paths and uses ranks of reduced representatives. Here each
    % path tensor itself has canonical-form-II data. See
    % docs/paper-gaps/mpu_canonical_form_full_support.tex.
    \uses{def:mpu_canonical_form_ii, def:mpu_right_rank, def:mpu_left_rank}
    Let $[0,1]\ni x\mapsto\mathcal U(x)$ be a continuous path of tensors of
    fixed physical dimension $d$ and bond dimension $D$. Suppose that each
    $\mathcal U(x)$ is in canonical form II. Then for all $x,y\in[0,1]$,
    \begin{align}
        r[\mathcal U(x)_{D^4}]    & =r[\mathcal U(y)_{D^4}],    &
        \ell[\mathcal U(x)_{D^4}] & =\ell[\mathcal U(y)_{D^4}].
    \end{align}
    This is the canonical-form-II path case of
    \cite[Proposition~IV.5]{Cirac2017MPU}. Its hypotheses are not comparable
    with those of Lemma~\ref{prop:mpu_admissible_continuity_index}: here
    each tensor of the path is in canonical form II, but no continuously
    varying source data are assumed; there the path may leave canonical form,
    but continuously varying source data at a common blocking length are
    assumed.
    % Source lines: paper_v2.tex 739--818.
\end{lemma}

\begin{proof}\leanok
    \uses{thm:mpu_common_pow_four_simple_blocking,
        prop:mpu_admissible_index_well_defined,
        thm:mpu_rank_lower_semicontinuous,
        thm:mpu_fixed_product_preconnected}
    Blocking is continuous, and the two source cuts are coordinate
    rearrangements of the blocked tensor. Thus their ranks are lower
    semicontinuous functions of $x$. The common block of length $D^4$ is
    simple at every point. Its ranks obey
    \begin{align}
        r[\mathcal U(x)_{D^4}]\ell[\mathcal U(x)_{D^4}]
         & =d^{2D^4}>0.
    \end{align}
    Lower semicontinuity and this fixed positive product make both ranks
    constant on the connected interval.
\end{proof}

\begin{lemma}[Index along a canonical-form-II path]
    \label{thm:mpu_cfii_path_index}
    \lean{MPOTensor.IsMPUCanonicalFormII.index_const_of_continuous}
    \leanok
    \uses{def:mpu_admissible_index, def:mpu_canonical_form_ii,
        thm:mpu_cfii_path_source_ranks}
    Under the hypotheses of Lemma~\ref{thm:mpu_cfii_path_source_ranks},
    the index of $\mathcal U(x)$ is independent of $x\in[0,1]$.
\end{lemma}

\begin{proof}\leanok
    \uses{thm:mpu_cfii_path_source_ranks,
        prop:mpu_admissible_index_well_defined}
    Every $D^4$-site block is simple, so for every $x\in[0,1]$,
    \begin{align}
        \operatorname{ind}(\mathcal U(x))
         & =\frac12\bigl(\log_2 r[\mathcal U(x)_{D^4}]
        -\log_2\ell[\mathcal U(x)_{D^4}]\bigr).
    \end{align}
    Both ranks on the right-hand side are independent of $x$ by
    Lemma~\ref{thm:mpu_cfii_path_source_ranks}.
\end{proof}

\begin{theorem}[Index Theorem]
    \label{IndexTh}
    \notready
    \discussion{6023}
    \uses{def:index, index-well-defined, lem:mpu_index_tensor, lem:mpu_index_composition, prop:continuity-index, def:equivalent-tensors}
    For MPU tensors under the source's standing canonical-form convention, take
    the index defined from any simple blocking. This index has the following
    properties.
    \begin{enumerate}
        \item It is invariant under blocking.
        \item It is additive under tensor products and composition.
        \item It is constant under continuous deformations through MPU tensors.
        \item Two MPU tensors have the same index if and only if they are
              equivalent.
    \end{enumerate}
    In \cite[Theorem~IV.6]{Cirac2017MPU}, these three clauses form the unrestricted
    index theorem.
    % Source lines: paper_v2.tex 824--865.
\end{theorem}

\begin{corollary}[Admissible Index Theorem]
    \label{thm:mpu_admissible_index}
    \notready
    \discussion{6023}
    % **Scope restriction (admissible tensors and operations):** every simple
    % block and standard form used below carries reduced full-support source
    % data. Documented in docs/paper-gaps/mpu_canonical_form_full_support.tex.
    \uses{prop:mpu_admissible_index_well_defined, thm:mpu_admissible_index_tensor, thm:mpu_representative_index_tensor_product, thm:mpu_admissible_index_composition, def:equivalent-tensors, def:mpu_admissible_equivalence_datum, prop:mpu_admissible_continuity_index, SF, FundamentalMPU, thm:mpu_canonical_form_ii_stabilization_converse, def:mpu_reduced_full_support_source_datum}
    Assume the following operation-specific admissibility hypotheses. Every pair
    of simple blockings compared for blocking independence carries reduced
    full-support source data. Tensor products need no hypothesis, by
    Lemma~\ref{thm:mpu_representative_index_tensor_product}.
    For a composition tensor $\mathcal W$ of $\mathcal U$ and $\mathcal U'$, there
    is one positive length $k$ such that
    \begin{align}
        \mathcal U_k,\quad \mathcal U'_k,\quad \mathcal W_k,\quad
        \mathcal U_{2k},\quad \mathcal U'_{2k},\quad \mathcal W_{4k}
    \end{align}
    are simple and carry their own data. Continuous-deformation and equivalence
    claims use the path data of
    Definition~\ref{def:mpu_admissible_equivalence_datum}. Under these
    hypotheses, the MPU index has the following properties.
    \begin{enumerate}
        \item It is invariant under blocking.
        \item It is additive under tensor products and composition.
        \item It is constant under continuous deformations through MPU tensors.
        \item Admissibly equivalent MPU tensors have the same index. Conversely,
              tensors with the same index are admissibly equivalent whenever the
              standard-form interpolation is equipped with an admissible path datum.
    \end{enumerate}
    In \cite[Theorem~IV.6]{Cirac2017MPU}, these three properties are the index-classification conclusions, here restricted
    to the explicit admissibility hypotheses above.
    % Source lines: paper_v2.tex 824--865.
\end{corollary}

\begin{corollary}[Continuous choice of standard form]
    \label{coruvSF}
    \notready
    \discussion{6024}
    \uses{prop:continuity-index, IndexTh, SF}
    For every continuous path $p\mapsto\mathcal U(p)$ of MPU tensors, not
    necessarily in canonical form, there are an integer $k_0\leq D^4$ and
    continuous paths of unitaries $u(p)$ and $v(p)$ such that
    $U^{(2k_0N)}(p)$ has standard form given by $u(p)$ and $v(p)$.
    % Source lines: paper_v2.tex 867--910.
\end{corollary}

\begin{corollary}[Admissible continuous choice of standard form]
    \label{cor:mpu_admissible_continuous_standard_form}
    \notready
    \discussion{6024}
    % **Scope restriction (continuous source-derived standard forms):** the chosen
    % simple blocks carry reduced full-support source data varying continuously
    % along the path. Documented in
    % docs/paper-gaps/mpu_canonical_form_full_support.tex.
    \uses{prop:mpu_admissible_continuity_index, thm:mpu_admissible_index, SF, ThmFund1, thm:mpu_canonical_form_ii_stabilization_converse, def:mpu_reduced_full_support_source_datum, def:mpu_admissible_equivalence_datum}
    Let $k_0$ be the common blocking length in an admissible equivalence
    datum for a continuous path $\mathcal U(p)$ of MPU tensors, and suppose that
    this given length satisfies $k_0\leq D^4$. Then there are continuous
    paths of unitaries $u(p)$ and $v(p)$ such that the MPU
    $U^{(2k_0N)}(p)$ has standard form given by $u(p)$ and $v(p)$ for every
    $p\in[0,1]$. For \cite[Corollary~IV.7]{Cirac2017MPU}, this statement is the reduced-source
    version.
    % Source lines: paper_v2.tex 867--910.
\end{corollary}

\subsection{Simple reduced representatives at the original scale}

\begin{lemma}[Doubled rectangular bond maps]
    \label{lem:mpu_doubled_rectangular_bond_map}
    \lean{MPOTensor.doubledBondMap, MPOTensor.doubleLayer_rectangularSandwich}
    \leanok
    \uses{def:mpu_double_layer}
    For a rectangular matrix $A$, write $\widehat A=\overline A\otimes A$,
    with the canonical enumeration of the product indices.
    If $V^{ij}=AU^{ij}B$, then its double-layer letters satisfy
    \begin{align}
        E_V^{ij} & =\widehat A E_U^{ij}\widehat B.
        \notag
    \end{align}
    This identity holds for arbitrary rectangular matrices; it does not
    require an inverse or an isometry.
\end{lemma}

\begin{lemma}[Normalized diagonal insertions]
    \label{lem:mpu_normalized_diagonal_insertions}
    \lean{MPOTensor.normalizedDiagonal_eq_smul_sum_diagonal,
        MPOTensor.IsMPUSimple.letter_mul_normalizedDiagonal_mul_letter,
        MPOTensor.simple2_of_normalizedDiagonal_insertion}
    \leanok
    \uses{def:mpu_simple, def:mpu_double_layer}
    For double-layer letters $E^{ij}$, let
    $E=d^{-1}\sum_{q=0}^{d-1}E^{qq}$, where $d>0$.
    A simple tensor satisfies
    \begin{align}
        E^{ij} E E^{kl} & =E^{ij}E^{kl}.
        \notag
    \end{align}
    Conversely, suppose this insertion identity holds and
    $E^J=\rho\Phi$ for some integer $J>0$, with a column $\rho$ and row $\Phi$.
    Then
    \begin{align}
        E^{ij}E^{kl} & =E^{ij}\rho\Phi E^{kl}.
        \notag
    \end{align}
\end{lemma}

\begin{proof}\leanok
    \uses{def:mpu_simple}
    Let $a,b$ be the simplicity vectors and put $R=|b)(a|$. Since
    $(a|E^{qq}|b)=1$, we have $RE^{qq}R=R$, and two insertions of $R$ give
    \begin{align}
        E^{ij}E^{qq}E^{kl} & =E^{ij}RE^{qq}RE^{kl}=E^{ij}RE^{kl}=E^{ij}E^{kl}.
        \notag
    \end{align}
    Averaging over $q$ gives the first identity. For the converse, expanding
    the last copy of $E$ gives $E^{ij}E^2=E^{ij}E$, so induction permits
    inserting every positive power, in particular $E^J=\rho\Phi$.
\end{proof}

\begin{lemma}[Doubled open words under a rectangular reduction]
    \label{lem:mpu_doubled_open_word_reduction}
    \lean{MPOTensor.doubleLayer_evalWord_of_isReduction}
    \leanok
    \uses{lem:mpu_doubled_rectangular_bond_map}
    Suppose $AB=\Id$ and $AU^wB=V^w$ for every open physical word $w$,
    including the empty word. For every pair of words $\sigma,\tau$ of the
    same length $N\geq0$,
    \begin{align}
        \widehat A E_U^{\sigma_1\tau_1}\cdots
        E_U^{\sigma_N\tau_N}\widehat B
         & =E_V^{\sigma_1\tau_1}\cdots E_V^{\sigma_N\tau_N}.
        \notag
    \end{align}
    % The full open-word identities are required, rather than their
    % single-letter consequences alone.
\end{lemma}

\begin{proof}\leanok
    \uses{lem:mpu_doubled_rectangular_bond_map,
        lem:mpu_rectangular_reduction_source_ranks,
        thm:mpu_double_layer_blocking}
    The open-word identities of length $N$ say that the $N$-site blocks
    satisfy $V_N^{IJ}=AU_N^{IJ}B$. Apply the preceding lemma to these
    blocks; the double layer of an $N$-site block at the pair of words
    $(\sigma,\tau)$ is the product $E^{\sigma_1\tau_1}\cdots E^{\sigma_N\tau_N}$.
\end{proof}

\begin{lemma}[Simplicity descends to a canonical reduction]
    \label{lem:mpu_simple_canonical_reduction}
    \lean{MPOTensor.IsMPUSimple.normalizedDiagonal_insertion_of_isReduction,
        MPOTensor.IsMPUCanonicalFormII.isMPUSimple_of_isReduction}
    \leanok
    \uses{def:mpu_simple, def:mpu_double_layer, def:mpu_canonical_form_ii}
    Let $d>0$, let $U$ be simple, and suppose $AB=\Id$ and $AU^wB=V^w$
    for every open word $w$. Then the double layer of $V$ satisfies
    \begin{align}
        E_V^{ij}E_VE_V^{kl} & =E_V^{ij}E_V^{kl}.
        \notag
    \end{align}
    If $V$ is in canonical form II, then $V$ is simple at the same physical
    scale.
\end{lemma}

\begin{proof}\leanok
    \uses{lem:mpu_doubled_open_word_reduction,
        lem:mpu_normalized_diagonal_insertions,
        thm:mpu_canonical_form_ii_transfer_stabilization,
        thm:mpu_all_later_simple_blockings}
    Write $E_U$ and $E_V$ for the normalized diagonals. The three- and
    two-letter cases of Lemma~\ref{lem:mpu_doubled_open_word_reduction},
    averaged over the middle letter, and the insertion identity for $U$
    give
    \begin{align}
        E_V^{ij}E_VE_V^{kl}
         & =\widehat A\,E_U^{ij}E_UE_U^{kl}\,\widehat B
        =\widehat A\,E_U^{ij}E_U^{kl}\,\widehat B
        =E_V^{ij}E_V^{kl}.
        \notag
    \end{align}
    This is the first assertion. If $V$ is in canonical form II, its
    normalized diagonal has a rank-one power $E_V^J=\rho'\Phi'$ with
    $J>0$. The converse part of
    Lemma~\ref{lem:mpu_normalized_diagonal_insertions} then gives
    \begin{align}
        E_V^{ij}E_V^{kl} & =E_V^{ij}\rho'\Phi'E_V^{kl},
        \notag
    \end{align}
    and this two-letter rank-one insertion identity makes the MPU tensor
    $V$ simple.
    % This argument does not replace V by a further physical block.
\end{proof}

\begin{lemma}[Simple reduced canonical representative]
    \label{lem:mpu_simple_reduced_canonical_representative}
    \lean{MPOTensor.IsMPUSimple.exists_reduced_cfii_representative}
    \leanok
    \uses{def:mpu_simple, def:mpu_canonical_form_ii}
    Let $d,D>0$ and let $U$ be a simple MPU tensor of bond dimension $D$.
    There is a simple tensor $V$ in canonical form II, of bond dimension
    $0<E\leq D$, such that $V^{(N)}=U^{(N)}$ for every $N>0$.
\end{lemma}

\begin{proof}\leanok
    \uses{thm:mpu_reduced_cfii_representative,
        thm:mpu_canonical_representative_reduction,
        lem:mpu_simple_canonical_reduction}
    Take the reduced canonical-form-II representative $V$ of $U$. The
    canonical representative is normal, so equality of the periodic
    families supplies a rectangular reduction of all open words onto it.
    The preceding lemma then supplies simplicity at the original scale.
\end{proof}

\begin{lemma}[Literal standard form of a reduced simple MPU]
    \label{lem:mpu_simple_reduced_literal_standard_form}
    \lean{MPOTensor.IsMPUSimple.exists_reduced_twoSiteStandardFormData}
    \leanok
    \uses{def:mpu_simple, def:mpu_canonical_form_ii,
        def:mpu_supplied_two_site_standard_form, def:mpu_source_uv}
    Every simple MPU tensor $U$ with $d,D>0$ has a simple canonical
    representative $V$, with $0<E\leq D$ and $V^{(N)}=U^{(N)}$ for all $N>0$,
    whose two-site block has standard form with the literal corrected source
    gates
    \begin{align}
        u & =Y_2Y_1, & v & =X_1X_2.
        \notag
    \end{align}
    The input tensor need not be in canonical form II or have full bond
    support. The sole physical blocking in this conclusion is the two-site
    blocking in the standard form of \cite[Definition~III.9]{Cirac2017MPU}.
    % No virtual similarity between the potentially unequal bond spaces is
    % asserted.
\end{lemma}

\begin{proof}\leanok
    \uses{lem:mpu_simple_reduced_canonical_representative,
        def:mpu_source_open_two_site_standard_form}
    Take the simple canonical representative $V$ of the preceding lemma.
    Its two-site block has the open standard form built from the source
    cuts of $V$, whose gates are the source operators of $V$.
\end{proof}
